\documentclass[12pt,a4paper]{report}
\usepackage[utf8]{inputenc}
\usepackage[spanish]{babel}
\usepackage{amsmath}
\usepackage{amsfonts}
\usepackage{amssymb}
\usepackage{lmodern}
\usepackage{amsmath}
\usepackage{enumerate}
\usepackage[left=2cm,right=2cm,top=2cm,bottom=2cm]{geometry}
\usepackage{graphicx}
\usepackage{mathtools}
\usepackage{stackrel}
\renewcommand{\theequation}{\arabic{equation}}
\newcounter{neq}
\providecommand{\abs}[1]{\lvert#1\rvert}
\newcommand{\QED}{\hfill \textit{\textbf{Q.E.D.}}}
\author{Agustín Curto, agucurto95@gmail.com \\
			 Francisco Nievas, frannievas@gmail.com}
\title{Resumen de teoremas para el final \\ de Lenguajes Formales y Computabilidad}
\date{2017}

\begin{document}
\maketitle
\tableofcontents

\section{Conjuntos $\Sigma$-mixtos}

  \textbf{\underline{Lemma 1:}} Sea $S\subseteq \omega \times \Sigma ^{\ast }$, entonces $S$ es rectangular si y
    solo si se cumple la siguiente propiedad:

    \[
      \textup{Si } (x, \alpha), (y, \beta) \in S \Rightarrow (x, \beta) \in S
    \]

  \textbf{\underline{Proof:}} Ejercicio.

  \QED

\section{Ordenes naturales sobre $\Sigma^{\ast}$}

  \textbf{\underline{Lemma 2:}} La relación $<$ es un orden total estricto sobre $\Sigma^{\ast}$.

  \textbf{\underline{Proof:}} Ejercicio.

  \QED

  \vspace{5mm}
  \textbf{\underline{Lemma 3:}} La función $s^{<}: \Sigma^{\ast} \rightarrow \Sigma^{\ast}$, definida recursivamente
    de la siguiente manera:

    \begin{eqnarray}
			\nonumber s^{<}(\varepsilon) &=& a_{1} \\
			\nonumber s^{<}(\alpha a_{i}) &=& \alpha a_{i + 1} \text{, } i < n \\
			\nonumber s^{<}(\alpha a_{n}) &=& s^{<}(\alpha) a_{1}
		\end{eqnarray}

    tiene la siguiente propiedad:

    \[
      \displaystyle s^{<}(\alpha) = \min \{\beta \in \Sigma^{\ast}: \alpha < \beta \}\text{.}
    \]

  \textbf{\underline{Proof:}} Supongamos que $\alpha < \beta $. Probaremos entonces que $s^{<}(\alpha )\leq \beta $.

    \vspace{3mm}
    \begin{tabular}{|c|}
      \hline Caso $\left\vert \alpha \right\vert < \left\vert \beta \right\vert$\\\hline
    \end{tabular}

    \par Se puede ver fácilmente que $\left\vert \alpha \right\vert = \left\vert s^{<}(\alpha) \right\vert$ salvo en el
    caso en que $\alpha \in \{a_{n}\}^{\ast}$, por lo cual solo resta ver el caso $\alpha \in \{a_{n}\}^{\ast}$.
    Supongamos $\alpha = a_{n}^{\left\vert \alpha \right\vert}$, entonces $s^{<}(\alpha) = a_{1}^{\left\vert \alpha
    \right\vert + 1}$. Si $\left\vert \beta \right\vert = \left\vert \alpha \right\vert + 1$ entonces es fácil ver
    usando el ítem 2 de la definición del orden de $\Sigma^{\ast}$ que $s^{<}(\alpha) = a_{1}^{\left\vert \alpha
    \right\vert + 1} \leq \beta$. Si $\left\vert \beta \right\vert > \left\vert \alpha \right\vert + 1$, entonces por
    el ítem 1, de tal definición tenemos que $s^{<}(\alpha) = a_{1}^{\left\vert \alpha \right\vert +1}< \beta$.


    \vspace{3mm}
    \begin{tabular}{|c|}
      \hline Caso $\left\vert \alpha \right\vert =\left\vert \beta \right\vert$\\\hline
    \end{tabular}

    \par Tenemos entonces que:

    \begin{eqnarray}
      \nonumber \alpha &=& \alpha_{1} a_{i} \gamma_{1} \\
      \nonumber \beta &=& \alpha_{1} a_{j} \gamma_{2}
    \end{eqnarray}

    con $i < j$ y $\left\vert \gamma_{1} \right\vert = \left\vert \gamma_{2}\right\vert$. Si $\gamma_{1} = \gamma_{2} =
    \varepsilon$ entonces es claro que $s^{<}(\alpha) \leq \beta$. El caso en el que $\gamma_{1}$ termina con $a_{l}$
    para algún $l < n$ es fácil. Veamos el caso en que $\gamma_{1} = a_{n}^{k}$ con $k \geq 1.$ Tenemos que:

    \begin{eqnarray}
      \nonumber s^{<}(\alpha ) & = & s^{<}(\alpha _{1}a_{i}a_{n}^{k}) \\
      \nonumber &=& s^{<}(\alpha_{1} a_{i} a_{n}^{k - 1}) a_{1} \\
      \nonumber &\vdots& \;\;\;\;\;\;\vdots \\
      \nonumber &=& s^{<}(\alpha_{1} a_{i}) a_{1}^{k} \\
      \nonumber &=& \alpha_{1} a_{i + 1} a_{1}^{k} \\
      \nonumber &\leq& \alpha_{1} a_{j} \gamma_{2} = \beta
    \end{eqnarray}

    \par Supongamos finalmente que $\gamma_{1} = \rho_{1} a_{l} a_{n}^{k}$ con $k \geq 1$ y $l < n$. Tenemos que:

    \begin{eqnarray}
      \nonumber s^{<}(\alpha) &=& s^{<}(\alpha_{1} a_{i} \rho_{1} a_{l} a_{n}^{k}) \\
      \nonumber &=& s^{<}(\alpha_{1} a_{i} \rho_{1} a_{l} a_{n}^{k - 1}) a_{1} \\
      \nonumber &\vdots& \;\;\;\;\;\;\vdots \\
      \nonumber &=& s^{<}(\alpha_{1} a_{i} \rho_{1} a_{l}) a_{1}^{k} \\
      \nonumber &=& \alpha_{1} a_{i} \rho_{1} a_{l + 1} a_{1}^{k} \\
      \nonumber &\leq& \beta
    \end{eqnarray}

    \par Para completar nuestra demostración debemos probar que $\alpha < s^{<}(\alpha)$, para cada $\alpha \in
    \Sigma ^{\ast}$. Dejamos al lector como ejercicio esta Proof la cual puede ser hecha por inducción en $\left\vert
    \alpha \right\vert$ usando argumentos parecidos a los usados anteriormente. $\Box$

  \QED


  \vspace{5mm}
  \textbf{\underline{Corollary 4:}} $s^{<}$ es inyectiva.

  \textbf{\underline{Proof:}} Supongamos $\alpha \neq \beta$. Ya que el orden de $\Sigma^{\ast}$ es total podemos
    suponer sin pérdida de generalidad que $\alpha < \beta$. Por el lema anterior tenemos que $s^{<}(\alpha) \leq \beta
    < s^{<}(\beta)$ y ya que $<$ es transitiva obtenemos que $s^{<}(\alpha)< s^{<}(\beta)$, lo cual nos dice $s^{<}
    (\alpha) \neq s^{<}(\beta)$.

  \QED


  \textbf{\underline{Lemma 5:}} Se tiene que:
    \begin{enumerate}
      \item $\varepsilon \neq s^{<}(\alpha)$, para cada $\alpha \in \Sigma^{\ast}$.
      \item Si $\alpha \neq \varepsilon$, entonces $\alpha = s^{<}(\beta)$ para algún $\beta$.
      \item Si $S\subseteq \Sigma^{\ast}$ es no vacío, entonces $\exists \alpha \in S$ tal que $\alpha < \beta$, para
      cada $\beta \in S-\{\alpha\}$.
    \end{enumerate}

  \textbf{\underline{Proof:}}

    \begin{enumerate}
      \item Ejercicio
      \item Ejercicio
      \item Sea $k = \min \{\left\vert \alpha \right\vert: \alpha \in S\}$. Notese que hay una cantidad finita de
      palabras de $S$ con longitud igual a $k$ y que la menor de ellas es justamente la menor palabra de $S$.
    \end{enumerate}

  \QED


  \textbf{\underline{Lemma 6:}} Tenemos que:

    \[
      \Sigma ^{\ast }=\{\ast^{<}(0),\ast^{<}(1),...\}
    \]

    \par Mas aún la función $\ast^{<}$ es biyectiva.

  \textbf{\underline{Proof:}} Supongamos $\ast^{<}(x) = \ast^{<}(y)$ con $x > y$. Note que $y \neq 0$ ya que
    $\varepsilon$ no es el sucesor de ninguna palabra. Osea que $s^{<}(\ast^{<}(x - 1)) = s^{<}(\ast^{<}(y - 1))$ lo
    cual ya que $\ast^{<}$ es inyectiva nos dice que $\ast^{<}(x - 1) = \ast^{<}(y - 1)$. Iterando este razonamiento
    llegamos a que $\ast^{<}(z) = \ast^{<}(0) = \varepsilon$ para algún $z > 0$, lo cual es absurdo.

    \par Veamos que $\ast^{<}$ es sobreyectiva. Supongamos no lo es, es decir supongamos que
    $\Sigma^{\ast}-I_{\ast^{<}} \neq \varnothing$. Por (3) del lema anterior $\Sigma^{\ast}-I_{\ast^{<}}$ tiene un
    menor elemento $\alpha$. Ya que $\alpha \neq \varepsilon$, tenemos que $\alpha = s^{<}(\beta)$, para algún $\beta$.
    Ya que $\beta < \alpha$ tenemos que $\beta \notin \Sigma^{\ast}-I_{\ast^{<}}$, es decir que $\beta = \ast^{<}(x)$,
    para algún $x \in \omega$. Esto nos dice que $\alpha = s^{<}(\ast^{<}(x))$, lo cual por la definición de
    $\ast^{<}$ nos dice que $\alpha = \ast^{<}(x + 1)$. Pero esto es absurdo ya que $\alpha \notin I_{\ast^{<}}$.

  \QED


  \textbf{\underline{Lemma 7:}} Sea $n \geq 1$ fijo, entonces cada $x \geq 1$ se escribe en forma única de la siguiente
    manera:

    \[
      x = i_{k} n^{k} + i_{k - 1} n^{k - 1} + ... + i_{0} n^{0}
    \]

    \par con $k\geq 0$ y $1\leq i_{k},i_{k-1},...,i_{0}\leq n$.

  \textbf{\underline{Proof:}} Primero la unicidad. Supongamos que:

    \[
      i_{k} n^{k} + i_{k - 1} n^{k - 1} + ... + i_{0} n^{0} = j_{m} n^{m} + j_{m - 1} n^{m - 1} + ... + j_{0} n^{0}
    \]

    \par con $k, m \geq 0$ y $1 \leq i_{k}, i_{k - 1}, ..., i_{0}, j_{m}, ..., j_{0} \leq n$. Supongamos $k < m$.
    Llegaremos a un absurdo. Notese que:

    \begin{eqnarray}
			\nonumber i_{k}n^{k}+i_{k-1}n^{k-1}+...+i_{0}n^{0} & \leq & n.n^{k}+n.n^{k-1}+...+n.n^{0} \\
			\nonumber & \leq & n^{k+1}+n^{k}+...+n^{1} \\ & < & n^{k+1}+n^{k}+...+n^{1}+n^{0} \\
			\nonumber & \leq & n^{m}+n^{m-1}+...+n^{0} \\
      \nonumber & \leq & j_{m}n^{m}+j_{m-1}n^{m-1}+...+j_{0}n^{0}
		\end{eqnarray}

    \par lo cual contradice la primera igualdad.

    \par Probaremos por inducción en $x$ que: (1) existen $k \geq 0$ y $i_{k}, i_{k - 1}, ..., i_{0} \in \{1, ..., n\}$
    tales que

    \[
      x = i_{k} n^{k} + i_{k - 1} n^{k - 1} + ... + i_{0} n^{0}
    \]

    El caso $x=1$ es trivial. Supongamos (1) vale para $x$, probaremos que vale para $x + 1$. Hay varios casos:

    Caso $i_{0} < n$. Entonces:

    \begin{eqnarray}
      \nonumber x+1 & =\left( i_{k}n^{k}+i_{k-1}n^{k-1}+...+i_{0}n^{0}\right) +1 \\
      \nonumber & =i_{k}n^{k}+i_{k-1}n^{k-1}+...+(i_{0}+1)n^{0}
    \end{eqnarray}

    Caso $i_{k}=i_{k-1}=...=i_{0}=n$. Tenemos que:

    \begin{eqnarray}
      \nonumber x+1 & =\left( i_{k}n^{k}+i_{k-1}n^{k-1}+...+i_{0}n^{0}\right) +1 \\
      \nonumber & =\left( nn^{k}+nn^{k-1}+...+nn^{0}\right) +1 \\
      \nonumber & =1n^{k+1}+1n^{k}+...+1n^{1}+1n^{0}
    \end{eqnarray}

    Caso $i_{0}=i_{1}=...=i_{h}=n$, $\;i_{h+1}\not=n$, para algun $ 0\leq h< k$. Tenemos:

    \begin{eqnarray}
      \nonumber x+1 & =\left( i_{k}n^{k}+...+i_{h+2}n^{h+2}+i_{h+1}n^{h+1}+nn^{h}+...+nn^{0}\right) +1 \\
      \nonumber & =\left( i_{k}n^{k}+...+i_{h+2}n^{h+2}+i_{h+1}n^{h+1}+n^{h+1}+n^{h}+...+n^{1}\right) +1 \\
      \nonumber & =i_{k}n^{k}+...+i_{h+2}n^{h+2}+(i_{h+1}+1)n^{h+1}+1n^{h}+...+1n^{1}+1n^{0}
    \end{eqnarray}

  \QED


  \textbf{\underline{Lemma 8:}} La función $\#^{<}$ es biyectiva.

  \textbf{\underline{Lemma 9:}} Las funciones $\#^{<}$ y $\ast^{<}$ son una inversa de la otra.

  \textbf{\underline{Proof:}} Probaremos por inducción en $x$ que para cada $x \in \omega$, se tiene que
    $\#^{<}(\ast^{<}(x)) = x$. El caso $x = 0$ es trivial. Supongamos que $\#^{<}(\ast^{<}(x)) = x$, veremos entonces
    que $\#^{<}(\ast^{<}(x + 1)) = x + 1$. Sean $k \geq 0$ y $i_{k}, ..., i_{0}$ tales que $\ast^{<}(x) = a_{i_{0}} ...
    a_{i_{0}}$. Ya que $\#^{<}(\ast^{<}(x)) = x$ tenemos que $x = i_{k} n^{k} + ... + i_{0} n^{0}$. Hay varios casos:

    Caso $i_{0}< n$. Entonces $\ast^{<}(x + 1)=s^{<}(\ast^{<}(x)) = a_{i_{k}} ... a_{i_{0} + 1}$ por lo cual:

    \begin{eqnarray}
      \nonumber \#^{<}(\ast^{<}(x+1)) &=& i_{k}n^{k}+i_{k-1}n^{k-1}+...+(i_{0}+1)n^{0} \\
      \nonumber &=& \left( i_{k}n^{k}+i_{k-1}n^{k-1}+...+i_{0}n^{0}\right) +1 \\
      \nonumber &=& x + 1
    \end{eqnarray}

    Caso $i_{k}=i_{k-1}=...=i_{0}=n$. Entonces $\ast^{<}(x + 1) = s^{<}(\ast^{<}(x)) = a_{1}^{k + 2}$ por lo cual:

    \begin{eqnarray}
      \nonumber \#^{<}(\ast^{<}(x+1)) &=& 1n^{k+1}+1n^{k}+...+1n^{1}+1n^{0} \\
      \nonumber &=& \left( nn^{k}+nn^{k-1}+...+nn^{0}\right) +1 \\
      \nonumber &=& x + 1
    \end{eqnarray}

    Caso $i_{0}=i_{1}=...=i_{h}=n$, $\;i_{h+1}\not=n$, para algun $ 0\leq h< k$. Entonces $\ast^{<}(x+1)=s^{<}(\ast^{<}(x))=a_{i_{k}}...a_{i_{h+2}}a_{i_{h+1}+1}a_{1}...a_{1}$ por lo cual

    \begin{eqnarray}
      \nonumber \#^{<}(\ast^{<}(x+1)) &=& i_{k}n^{k}+...+i_{h+2}n^{h+2}+(i_{h+1}+1)n^{h+1}+1n^{h}+...+1n^{1}+1n^{0} \\
      \nonumber &=& \left( i_{k}n^{k}+...+i_{h+2}n^{h+2}+i_{h+1}n^{h+1}+n^{h+1}+n^{h}+...+n^{1}\right) +1 \\
      \nonumber &=& \left( i_{k}n^{k}+...+i_{h+2}n^{h+2}+i_{h+1}n^{h+1}+nn^{h}+...+nn^{0}\right) +1 \\
      \nonumber &=& x + 1
    \end{eqnarray}

  \QED

\section{Codificacion de sucesiones infinitas de numeros}

  \textbf{\underline{Lemma 10:}} Si \(p,p_{1},...,p_{n}\) son numeros primos y \(p\) divide a \(p_{1}...p_{n}\), entonces \(p=p_{i}\), para algun \(i\).
Ahora la version del Teorema Fundamental de la Aritmetica.

  \textbf{\underline{Theorem 11:}} Para cada \(x\in \mathbf{N}\), hay una unica sucesion \((s_{1},s_{2},...)\in \omega ^{\left[ \mathbf{N}\right] }\) tal que
\(\displaystyle x=\underset{i=1}{\overset{\infty }{\Pi }}pr(i)^{s_{i}} \)
(Notese que \(\underset{i=1}{\overset{\infty }{\Pi }}pr(i)^{s_{i}}\) tiene sentido ya que es un producto que solo tiene una cantidad finita de factores no iguales a \(1\). )

\textbf{\underline{Proof:}} Primero probaremos la existencia por induccion en \(x\). Claramente \(1= \underset{i=1}{\overset{\infty }{\Pi }}pr(i)^{0}\), con lo cual el caso \(x=1\) esta probado. Supongamos la existencia vale para cada \(y\) menor que \(x\), veremos que entonces vale para \(x\). Si \(x\) es primo, entonces \(x=pr(i_{0})\) para algun \(i_{0}\) por lo cual tenemos que \(x=\underset{i=1}{\overset{\infty }{\Pi }}pr(i)^{s_{i}}\), tomando \(s_{i}=0\) si \(i\neq i_{0}\) y \(s_{i_{0}}=1\). Si \(x\) no es primo, entonces \(x=y_{1}.y_{2}\) con \(y_{1},y_{2}< x\). Por hipotesis inductiva tenemos que hay \((s_{1},s_{2},...),(t_{1},t_{2},...)\in \omega ^{\left[ \mathbf{N}\right] }\) tales que \(y_{1}=\underset{i=1}{\overset {\infty }{\Pi }}pr(i)^{s_{i}}\) y \(y_{2}=\underset{i=1}{\overset{\infty }{\Pi }}pr(i)^{t_{i}}\). Tenemos entonces que \(x=\underset{i=1}{\overset{\infty }{ \Pi }}pr(i)^{s_{i}+t_{i}}\) lo cual concluye la prueba de la existencia.

Veamos ahora la unicidad. Suponganos que

\(\displaystyle \underset{i=1}{\overset{\infty }{\Pi }}pr(i)^{s_{i}}=\underset{i=1}{\overset{ \infty }{\Pi }}pr(i)^{t_{i}} \)

Si \(s_{i} >t_{i}\) entonces dividiendo ambos miembros por \(pr(i)^{t_{i}}\) obtenemos que \(pr(i)\) divide a un producto de primos todos distintos de el, lo cual es absurdo por el lema anterior. Analogamente llegamos a un absurdo si suponemos que \(t_{i} >s_{i}\), lo cual nos dice que \(s_{i}=t_{i}\), para cada \(i\in \mathbf{N}\) \(\Box\)

  \textbf{\underline{Lemma 12:}} Las funciones
\(\displaystyle \begin{array}{lll} \mathbf{N} & \rightarrow & \omega ^{\left[ \mathbf{N}\right] } \\ x & \rightarrow & ((x)_{1},(x)_{2},...) \end{array} \ \ \ \ \ \ \ \ \ \ \ \ \ \ \ \ \ \ \begin{array}{rll} \omega ^{\left[ \mathbf{N}\right] } & \rightarrow & \mathbf{N} \\ (s_{1},s_{2},...) & \rightarrow & \left\langle s_{1},s_{2},...\right\rangle \end{array} \)
son biyecciones una inversa de la otra.
\textbf{\underline{Proof:}} Notese que para cada \(x\in \mathbf{N}\), tenemos que \(\left\langle (x)_{1},(x)_{2},...\right\rangle =x\). Ademas para cada \((s_{1},s_{2},...)\in \omega ^{\left[ \mathbf{N}\right] }\), tenemos que \(((\left\langle s_{1},s_{2},...\right\rangle )_{1},(\left\langle s_{1},s_{2},...\right\rangle )_{2},...)=(s_{1},s_{2},...)\). Es claro que lo anterior garantiza que los mapeos en cuestion son uno inversa del otro \(\Box\)

  \textbf{\underline{Lemma 13}} Para cada \(x\in \mathbf{N}\):
\(Lt(x)=0\) sii \(x=1\)
\(x=\prod\nolimits_{i=1}^{Lt(x)}pr(i)^{(x)_{i}}\)
Cabe destacar entonces que la funcion \(\lambda ix[(x)_{i}]\) tiene dominio igual a \(\mathbf{N}^{2}\) y la funcion \(\lambda ix[Lt(x)]\) tiene dominio igual a \(\mathbf{N}\).

\section{Conjuntos $\Sigma$-efectivamente enumerables}

  \textbf{\underline{Lemma 14:}} Sean \(S_{1},S_{2}\subseteq \omega ^{n}\times \Sigma ^{\ast m}\) conjuntos \( \Sigma \)-efectivamente enumerables. Entonces \(S_{1}\cup S_{2}\) y \(S_{1}\cap S_{2}\) son \(\Sigma \)-efectivamente enumerables.

  \textbf{\underline{Proof:}} El caso en el que alguno de los conjuntos es vacio es trivial. Supongamos que ambos conjuntos son no vacios y sean \(\mathbb{P}_{1}\) y \(\mathbb{P}_{2}\) procedimientos que enumeran a \(S_{1}\) y \(S_{2}\). El siguiente procedimiento enumera al conjunto \(S_{1}\cup S_{2}\):

- Si \(x\) es par realizar \(\mathbb{P}_{1}\) partiendo de \(x/2\) y dar el elemento de \(S_{1}\) obtenido como salida. Si \(x\) es impar realizar \(\mathbb{P }_{2}\) partiendo de \((x-1)/2\) y dar el elemento de \(S_{2}\) obtenido como salida.
Veamos ahora que \(S_{1}\cap S_{2}\) es \(\Sigma \)-efectivamente enumerable. Si \(S_{1}\cap S_{2}=\varnothing \) entonces no hay nada que probar. Supongamos entonces que \(S_{1}\cap S_{2}\) en no vacio. Sea \(z_{0}\) un elemento fijo de \( S_{1}\cap S_{2}.\) Sea \(\mathbb{P}\) un procedimiento efectivo el cual enumere a \(\omega \times \omega \) (ver el ejemplo de mas arriba). Un procedimiento que enumera a \(S_{1}\cap S_{2}\) es el siguiente

Etapa 1

Realizar \(\mathbb{P}\) con dato de entrada \(x\), para obtener un par \((x_{1},x_{2})\in \omega \times \omega \).

Etapa 2

Realizar \(\mathbb{P}_{1}\) con dato de entrada \(x_{1}\) para obtener un elemento \(z_{1}\in S_{1}\)

Etapa 3

Realizar \(\mathbb{P}_{2}\) con dato de entrada \(x_{2}\) para obtener un elemento \(z_{2}\in S_{2}\)

Etapa 4

Si \(z_{1}=z_{2}\), entonces dar como dato de salida \(z_{1}.\) En caso contrario dar como dato de salida \(z_{0}\). \(\Box\)

\section{Conjuntos $\Sigma$-efectivamente computables}

  \textbf{\underline{Lemma 15:}} Si \(S\subseteq \omega ^{n}\times \Sigma ^{\ast m}\) es \(\Sigma \) -efectivamente computable entonces \(S\) es \(\Sigma \)-efectivamente enumerable.

  \textbf{\underline{Proof:}} Supongamos \(S\neq \varnothing \). Sea \((\vec{z},\gamma )\in S\), fijo. Sea \( \mathbb{P}\) un procedimiento efectivo que compute a \(\chi _{S}\). Ya vimos en el ejemplo anterior que \(\omega ^{2}\times \Sigma ^{\ast 3}\) es \(\Sigma \) -efectivamente enumerable. En forma similar se puede ver que \(\omega ^{n}\times \Sigma ^{\ast m}\) lo es. Sea \(\mathbb{P}_{1}\) un procedimiento efectivo que enumere a \(\omega ^{n}\times \Sigma ^{\ast m}\). Entonces el siguiente procedimiento enumera a \(S\):

Etapa 1

Realizar \(\mathbb{P}_{1}\) con \(x\) de entrada para obtener \((\vec{x} ,\vec{\alpha})\in \omega ^{n}\times \Sigma ^{\ast m}.\)

Etapa 2

Realizar \(\mathbb{P}\) con \((\vec{x},\vec{\alpha})\) de entrada para obtener el valor Booleano \(e\) de salida\(.\)

Etapa 3

Si \(e=1\) dar como dato de salida \((\vec{x},\vec{\alpha}).\) Si \(e=0\) dar como dato de salida \((\vec{z},\gamma )\). \(\Box\)


  \textbf{\underline{Theorem 16:}} Sea \(S\subseteq \omega ^{n}\times \Sigma ^{\ast m}\). Son equivalentes
  (a) \(S\) es \(\Sigma \)-efectivamente computable
  (b) \(S\) y \((\omega ^{n}\times \Sigma ^{\ast m})-S\) son \(\Sigma \) -efectivamente enumerables

  \textbf{\underline{Proof:}} (a)\(\Rightarrow \)(b). Por el lema anterior tenemos que \(S\) es \(\Sigma \) -efectivamente enumerable. Notese ademas que, dado que \(S\) es \(\Sigma \) -efectivamente computable, \((\omega ^{n}\times \Sigma ^{\ast m})-S\) tambien lo es (por que?). Es decir que aplicando nuevamente el lema anterior tenemos que \((\omega ^{n}\times \Sigma ^{\ast m})-S\) es \(\Sigma \)-efectivamente enumerable.

(b)\(\Rightarrow \)(a). Sea \(\mathbb{P}_{1}\) un procedimiento efectivo que enumere a \(S\) y sea \(\mathbb{P}_{2}\) un procedimiento efectivo que enumere a \((\omega ^{n}\times \Sigma ^{\ast m})-S\). Es facil ver que el siguiente procedimiento computa el predicado \(\chi _{S}\):

Etapa 1

Darle a la variable \(T\) el valor \(0\).

Etapa 2

Realizar \(\mathbb{P}_{1}\) con el valor de \(T\) como entrada para obtener de salida la upla \((\vec{y},\vec{\beta})\).

Etapa 3

Realizar \(\mathbb{P}_{2}\) con el valor de \(T\) como entrada para obtener de salida la upla \((\vec{z},\vec{\gamma})\).

Etapa 4

Si \((\vec{y},\vec{\beta})=(\vec{x},\vec{\alpha})\), entonces detenerse y dar como dato de salida el valor \(1\). Si \((\vec{z},\vec{\gamma} )=(\vec{x},\vec{\alpha})\), entonces detenerse y dar como dato de salida el valor \(0.\) Si no suceden ninguna de las dos posibilidades antes mensionadas, aumentar en \(1\) el valor de la variable \(T\) y dirijirse a la Etapa 2. \(\Box\)

  \textbf{\underline{Theorem 17:}} Dado \(S\subseteq \omega ^{n}\times \Sigma ^{\ast m}\), son equivalentes
(1) \(S\) es \(\Sigma \)-efectivamente enumerable
(2) \(S=\varnothing \) o \(S=I_{F}\), para alguna \(F:\omega \rightarrow \omega ^{n}\times \Sigma ^{\ast m}\) tal que cada \(F_{i}\) es \(\Sigma \) -efectivamente computable.
(3) \(S=I_{F}\), para alguna \(F:D_{F}\subseteq \omega ^{k}\times \Sigma ^{\ast l}\rightarrow \omega ^{n}\times \Sigma ^{\ast m}\) tal que cada \(F_{i}\) es \(\Sigma \)-efectivamente computable.
(4) \(S=D_{f}\), para alguna funcion \(f\) la cual es \(\Sigma \) -efectivamente computable.

  \textbf{\underline{Proof:}} (1)\(\Rightarrow \)(2) y (2)\(\Rightarrow \)(1) son muy naturales y son dejadas al lector. (2)\(\Rightarrow \)(3) es trivial.

(3)\(\Rightarrow \)(4). Para \(i=1,...,n+m\), sea \(\mathbb{P}_{i}\) un procedimiento el cual computa a \(F_{i}\) y sea \(\mathbb{P}\) un procedimiento el cual enumere a \(\omega \times \omega ^{k}\times \Sigma ^{\ast l}.\) El siguiente procedimiento computa la funcion \(f:I_{F}\rightarrow \{1\}\):

Etapa 1

Darle a la variable \(T\) el valor 0.

Etapa 2

Hacer correr \(\mathbb{P}\) con dato de entrada \(T\) y obtener \( (t,z_{1},...,z_{k},\gamma _{1},...,\gamma _{l})\) como dato de salida.

Etapa 3

Para cada \(i=1,...,n+m\), hacer correr \(\mathbb{P}_{i}\) durante \(t\) pasos, con dato de entrada \((z_{1},...,z_{k},\gamma _{1},...,\gamma _{l}).\) Si cada procedimiento \(\mathbb{P}_{i}\) al cabo de los \(t\) pasos termino y dio como resultado el valor \(o_{i}\), entonces comparar \((\vec{x},\vec{\alpha} )\) con \((o_{1},...,o_{n+m})\) y en caso de que sean iguales detenerse y dar como dato de salida el valor \(1\). En el caso en que no son iguales, aumentar en \(1\) el valor de la variable \(T\) y dirijirse a la Etapa 2. Si algun procedimiento \(\mathbb{P}_{i}\) al cabo de los \(t\) pasos no termino, entonces aumentar en \(1\) el valor de la variable \(T\) y dirijirse a la Etapa 2.

(4)\(\Rightarrow \)(1). Supongamos \(S\neq \varnothing .\) Sea \((\vec{z},\vec{ \gamma})\) un elemento fijo de \(S.\) Sea \(\mathbb{P}\) un procedimiento el cual compute a \(f\). Sea \(\mathbb{P}_{1}\) un procedimiento el cual enumere a \( \omega \times \omega ^{n}\times \Sigma ^{\ast m}.\) Dejamos al lector el dise \~{n}o de un procedimiento efectivo el cual enumere \(D_{f}\). \(\Box\)

\section{Funciones $\Sigma$-recursivas primitivas}

  \textbf{\underline{Lemma 18:}} Si \(f,f_{1},...,f_{n+m}\) son \(\Sigma \)-efectivamente computables, entonces \( f\circ (f_{1},...,f_{n+m})\) lo es.

  \textbf{\underline{Proof:}} Sean \(\mathbb{P},\mathbb{P}_{1},...,\mathbb{P}_{n+m}\) procedimientos efectivos los cuales computen las funciones \(f,f_{1},...,f_{n+m}\), respectivamente. Usando estos procedimientos es facil definir un procedimiento efectivo el cual compute a \(f\circ (f_{1},...,f_{n+m})\). \(\Box\)

  \textbf{\underline{Lemma 19:}} Si \(f\) y \(g\) son \(\Sigma \)-efectivamente computables, entonces \(R(f,g)\) lo es.

  \textbf{\underline{Proof:}} La Proof es dejada al lector. \(\Box\)

  \textbf{\underline{Lemma 20:}} Si \(f\) y cada \(\mathcal{G}_{a}\) son \(\Sigma \)-efectivamente computables, entonces \(R(f,\mathcal{G})\) lo es.

  \textbf{\underline{Proof:}} Es dejada al lector con la recomendacion de que haga la Proof para el caso \( \Sigma =\{@,\& \}\) \(\Box\)

  \textbf{\underline{Theorem 21:}} Si \(f\in \mathrm{PR}^{\Sigma }\), entonces \(f\) es \(\Sigma \)-efectivamente computable.

  \textbf{\underline{Proof:}} Dejamos al lector la Proof por induccion en \(k\) de que si \(f\in \mathrm{PR} _{k}^{\Sigma }\), entonces \(f\) es \(\Sigma \)-efectivamente computable, la cual sale en forma directa usando los lemas anteriores que garantizan que los constructores de composicion y recursion primitiva preservan la computabilidad efectiva \(\Box\)

  \textbf{\underline{Lemma 22:}}
  (1) \(\varnothing \in \mathrm{PR}^{\varnothing }\).
  (2) \(\lambda xy\left[ x+y\right] \in \mathrm{PR}^{\varnothing }\).
  (3) \(\lambda xy\left[ x.y\right] \in \mathrm{PR}^{\varnothing }\).
  (4) \(\lambda x\left[ x!\right] \in \mathrm{PR}^{\varnothing }\).

  \textbf{\underline{Proof:}} (1) Notese que \(\varnothing =Pred\circ C_{0}^{0,0}\in \mathrm{PR} _{1}^{\varnothing }\)

  (2) Notar que

  \(\displaystyle \begin{array}{rcl} \lambda xy\left[ x+y\right] (0,x_{1}) & =& x_{1}=p_{1}^{1,0}(x_{1}) \\ \lambda xy\left[ x+y\right] (t+1,x_{1}) & =& \lambda xy\left[ x+y\right] (t,x_{1})+1 \\ & =& \left( Suc\circ p_{1}^{3,0}\right) \left( \lambda xy\left[ x+y\right] (t,x_{1}),t,x_{1}\right) \end{array} \)

  lo cual implica que \(\lambda xy\left[ x+y\right] =R\left( p_{1}^{1,0},Suc\circ p_{1}^{3,0}\right) \in \mathrm{PR}_{2}^{\varnothing }.\)
  (3) Primero note que

  \(\displaystyle \begin{array}{rcl} C_{0}^{1,0}(0) & =& C_{0}^{0,0}(\Diamond ) \\ C_{0}^{1,0}(t+1) & =& C_{0}^{1,0}(t) \end{array} \)

  lo cual implica que \(C_{0}^{1,0}=R\left( C_{0}^{0,0},p_{1}^{2,0}\right) \in \mathrm{PR}_{1}^{\varnothing }.\) Tambien note que
  \(\displaystyle \lambda tx\left[ t.x\right] =R\left( C_{0}^{1,0},\lambda xy\left[ x+y\right] \circ \left( p_{1}^{3,0},p_{3}^{3,0}\right) \right) , \)

  lo cual por (1) implica que \(\lambda tx\left[ t.x\right] \in \mathrm{PR} _{3}^{\varnothing }\).
  (4) Note que

  \(\displaystyle \begin{array}{rcl} \lambda x\left[ x!\right] (0) & =& 1=C_{1}^{0,0}(\Diamond ) \\ \lambda x\left[ x!\right] (t+1) & =& \lambda x\left[ x!\right] (t).(t+1), \end{array} \)

  lo cual implica que
  \(\displaystyle \lambda x\left[ x!\right] =R\left( C_{1}^{0,0},\lambda xy\left[ x.y\right] \circ \left( p_{1}^{2,0},Suc\circ p_{2}^{2,0}\right) \right) . \)

  Ya que \(C_{1}^{0,0}=\) \(Suc\circ C_{0}^{0,0}\), tenemos que \(C_{1}^{0,0}\in \mathrm{PR}_{1}^{\varnothing }\). Por (2), tenemos que
  \(\displaystyle \lambda xy\left[ x.y\right] \circ \left( p_{1}^{2,0},Suc\circ p_{2}^{2,0}\right) \in \mathrm{PR}_{4}^{\varnothing }, \)

  obteniendo que \(\lambda x\left[ x!\right] \in \mathrm{PR}_{5}^{\varnothing }\). \(\Box\)


  \textbf{\underline{Lemma 23:}} Supongamos \(\Sigma \) es no vacio.
  (a) \(\lambda \alpha \beta \left[ \alpha \beta \right] \in \mathrm{PR} ^{\Sigma }\)
  (b) \(\lambda \alpha \left[ \left\vert \alpha \right\vert \right] \in \mathrm{PR}^{\Sigma }\)

  \textbf{\underline{Proof:}} (a) Ya que

  \(\displaystyle \begin{array}{rcl} \lambda \alpha \beta \left[ \alpha \beta \right] (\alpha _{1},\varepsilon ) & =& \alpha _{1}=p_{1}^{0,1}(\alpha _{1}) \\ \lambda \alpha \beta \left[ \alpha \beta \right] (\alpha _{1},\alpha a) & =& d_{a}(\lambda \alpha \beta \left[ \alpha \beta \right] (\alpha _{1},\alpha )),a\in \Sigma \end{array} \)

  tenemos que \(\lambda \alpha \beta \left[ \alpha \beta \right] =R\left( p_{1}^{0,1},\mathcal{G}\right) \), donde \(\mathcal{G}_{a}=d_{a}\circ p_{3}^{0,3}\), para cada \(a\in \Sigma \).
  (b) Ya que

  \(\displaystyle \begin{array}{rcl} \lambda \alpha \left[ \left\vert \alpha \right\vert \right] (\varepsilon ) & =& 0=C_{0}^{0,0}(\Diamond ) \\ \lambda \alpha \left[ \left\vert \alpha \right\vert \right] (\alpha a) & =& \lambda \alpha \left[ \left\vert \alpha \right\vert \right] (\alpha )+1 \end{array} \)

  tenemos que \(\lambda \alpha \left[ \left\vert \alpha \right\vert \right] =R\left( C_{0}^{0,0},\mathcal{G}\right) \), donde \(\mathcal{G}_{a}=\) \( Suc\circ p_{1}^{1,1}\), para cada \(a\in \Sigma .\). \(\Box\)


  \textbf{\underline{Lemma 24:}}
  (a) \(C_{k}^{n,m},C_{\alpha }^{n,m}\in \mathrm{PR}^{\Sigma }\), para \( n,m,k\geq 0\), \(\alpha \in \Sigma ^{\ast }\).
  (b) \(C_{k}^{n,0}\in \mathrm{PR}^{\varnothing }\), para \(n,k\geq 0\).


  \textbf{\underline{Proof:}} (a) Note que \(C_{k+1}^{0,0}=\) \(Suc\circ C_{k}^{0,0}\), lo cual implica \( C_{k}^{0,0}\in \mathrm{PR}_{k}^{\Sigma }\), para \(k\geq 0\). Tambien note que \( C_{\alpha a}^{0,0}=d_{a}\circ C_{\alpha }^{0,0}\), lo cual dice que \( C_{\alpha }^{0,0}\in \mathrm{PR}^{\Sigma }\), para \(\alpha \in \Sigma ^{\ast } \). Para ver que \(C_{k}^{0,1}\in \mathrm{PR}^{\Sigma }\) notar que

  \(\displaystyle \begin{array}{rcl} C_{k}^{0,1}(\varepsilon ) & =& k=C_{k}^{0,0}(\Diamond ) \\ C_{k}^{0,1}(\alpha a) & =& C_{k}^{0,1}(\alpha )=p_{1}^{1,1}\left( C_{k}^{0,1}(\alpha ),\alpha \right) \end{array} \)

  lo cual implica que \(C_{k}^{0,1}=R\left( C_{k}^{0,0},\mathcal{G}\right) \), con \(\mathcal{G}_{a}=p_{1}^{1,1}\), \(a\in \Sigma \). En forma similar podemos ver que \(C_{k}^{1,0},C_{\alpha }^{1,0},C_{\alpha }^{0,1}\in \mathrm{PR} ^{\Sigma }\). Supongamos ahora que \(m >0\). Entonces
  \(\displaystyle \begin{array}{rcl} C_{k}^{n,m} & =& C_{k}^{0,1}\circ p_{n+1}^{n,m} \\ C_{\alpha }^{n,m} & =& C_{\alpha }^{0,1}\circ p_{n+1}^{n,m} \end{array} \)

  de lo cual obtenemos que \(C_{k}^{n,m},C_{\alpha }^{n,m}\in \mathrm{PR} ^{\Sigma }\). El caso \(n >0\) es similar
  (b) Use argumentos similares a los usados en la Proof de (a). \(\Box\)


  \textbf{\underline{Lemma 25:}}
  (a) \(\lambda xy\left[ x^{y}\right] \in \mathrm{PR}^{\varnothing }\).
  (b) \(\lambda t\alpha \left[ \alpha ^{t}\right] \in \mathrm{PR} ^{\Sigma }\).


  \textbf{\underline{Proof:}} (a) Note que

  \(\displaystyle \lambda tx\left[ x^{t}\right] =R\left( C_{1}^{1,0},\lambda xy\left[ x.y \right] \circ \left( p_{1}^{3,0},p_{3}^{3,0}\right) \right) \in \mathrm{PR} ^{\varnothing }. \)

  O sea que \(\lambda xy\left[ x^{y}\right] =\lambda tx\left[ x^{t}\right] \circ \left( p_{2}^{2,0},p_{1}^{2,0}\right) \in \mathrm{PR}^{\varnothing }\).
  (b) Note que

  \(\displaystyle \lambda t\alpha \left[ \alpha ^{t}\right] =R\left( C_{\varepsilon }^{0,1},\lambda \alpha \beta \left[ \alpha \beta \right] \circ \left( p_{3}^{1,2},p_{2}^{1,2}\right) \right) \in \mathrm{PR}^{\Sigma }. \)

  \(\Box\)


  \textbf{\underline{Lemma 26:}} Si \(< \) es un orden total estricto sobre un alfabeto no vacio \( \Sigma \), entonces \(s^{< }\), \(\#^{< }\) y \(\ast ^{< }\) pertenecen a \(\mathrm{PR} ^{\Sigma }\)

  \textbf{\underline{Proof:}} Supongamos \(\Sigma =\{a_{1},...,a_{k}\}\) y \(< \) dado por \(a_{1}< ...< a_{k}\). Ya que

  \(\displaystyle \begin{array}{rcl} s^{< }(\varepsilon ) & =& a_{1} \\ s^{< }(\alpha a_{i}) & =& \alpha a_{i+1}\text{, para }i< k \\ s^{< }(\alpha a_{k}) & =& s^{< }(\alpha )a_{1} \end{array} \)

  tenemos que \(s^{< }=R\left( C_{a_{1}}^{0,0},\mathcal{G}\right) \), donde \( \mathcal{G}_{a_{i}}=d_{a_{i+1}}\circ p_{1}^{0,2}\), para \(i=1,...,k-1\) y \( \mathcal{G}_{a_{k}}=d_{a_{1}}\circ p_{2}^{0,2}.\) O sea que \(s^{< }\in \mathrm{ PR}^{\Sigma }.\) Ya que
  \(\displaystyle \begin{array}{rcl} \ast ^{< }(0) & =& \varepsilon \\ \ast ^{< }(t+1) & =& s^{< }(\ast ^{< }(t)) \end{array} \)

  podemos ver que \(\ast ^{< }\in \mathrm{PR}^{\Sigma }.\) Ya que
  \(\displaystyle \begin{array}{rcl} \#^{< }(\varepsilon ) & =& 0 \\ \#^{< }(\alpha a_{i}) & =& \#^{< }(\alpha ).k+i\text{, para }i=1,...,k, \end{array} \)

  tenemos que \(\#^{< }=R\left( C_{0}^{0,0},\mathcal{G}\right) \), donde
  \(\displaystyle \mathcal{G}_{a_{i}}=\lambda xy\left[ x+y\right] \circ \left( \lambda xy\left[ x.y\right] \circ \left( p_{1}^{1,1},C_{k}^{1,1}\right) ,C_{i}^{1,1}\right) \text{, para }i=1,...,k\text{.} \)

  O sea que \(\#^{< }\in \mathrm{PR}^{\Sigma }\). \(\Box\)

  \textbf{\underline{Lemma 27:}}
  (a) \(\lambda xy\left[ x\dot{-}y\right] \in \mathrm{PR}^{\varnothing }.\)
  (b) \(\lambda xy\left[ \max (x,y)\right] \in \mathrm{PR}^{\varnothing }.\)
  (c) \(\lambda xy\left[ x=y\right] \in \mathrm{PR}^{\varnothing }.\)
  (d) \(\lambda xy\left[ x\leq y\right] \in \mathrm{PR}^{\varnothing }.\)
  (e) Si \(\Sigma \) es no vacio, entonces \(\lambda \alpha \beta \left[ \alpha =\beta \right] \in \mathrm{PR}^{\Sigma }\)


  \textbf{\underline{Proof:}} (a) Primero notar que \(\lambda x\left[ x\dot{-}1\right] =R\left( C_{0}^{0,0},p_{2}^{2,0}\right) \in \mathrm{PR}^{\varnothing }.\) Tambien note que

  \(\displaystyle \lambda tx\left[ x\dot{-}t\right] =R\left( p_{1}^{1,0},\lambda x\left[ x\dot{ -}1\right] \circ p_{1}^{3,0}\right) \in \mathrm{PR}^{\varnothing }. \)

  O sea que \(\lambda xy\left[ x\dot{-}y\right] =\lambda tx\left[ x\dot{-}t \right] \circ \left( p_{2}^{2,0},p_{1}^{2,0}\right) \in \mathrm{PR} ^{\varnothing }.\)
  (b) Note que \(\lambda xy\left[ \max (x,y)\right] =\lambda xy\left[ (x+(y\dot{ -}x)\right] .\)

  (c) Note que \(\lambda xy\left[ x=y\right] =\lambda xy\left[ 1\dot{-}((x\dot{- }y)+(y\dot{-}x))\right] .\)

  (d) Note que \(\lambda xy\left[ x\leq y\right] =\lambda xy\left[ 1\dot{-}(x \dot{-}y)\right] .\)

  (e) Sea \(< \) un orden total estricto sobre \(\Sigma .\) Ya que

  \(\displaystyle \alpha =\beta \text{ sii }\#^{< }(\alpha )=\#^{< }(\beta ) \)

  tenemos que
  \(\displaystyle \lambda \alpha \beta \left[ \alpha =\beta \right] =\lambda xy\left[ x=y \right] \circ \left( \#^{< }\circ p_{1}^{0,2},\#^{< }\circ p_{2}^{0,2}\right) . \)

  O sea que podemos aplicar (c) y Lema 28 implica que \(\chi _{S}\) es \( \Sigma \)-p.r.. \(\Box\)


  \textbf{\underline{Lemma 28:}}
    \textbf{Hacer}
  \textbf{\underline{Proof:}}

  \textbf{\underline{Lemma 29:}}
    \textbf{Hacer}

  \textbf{\underline{Proof:}}

  \textbf{\underline{Corollary 30:}}
    \textbf{Hacer}

  \textbf{\underline{Proof:}}

  %% =========== FALTA UNA PARTE ===========%%

  \textbf{\underline{Lemma 31:}} Supongamos \(S_{1},...,S_{n}\subseteq \omega \), \( L_{1},...,L_{m}\subseteq \Sigma ^{\ast }\) son conjuntos no vacios. Entonces \( S_{1}\times ...\times S_{n}\times L_{1}\times ...\times L_{m}\) es \(\Sigma \) -p.r. sii \(S_{1},...,S_{n},L_{1},...,L_{m}\) son \(\Sigma \)-p.r.

  \textbf{\underline{Proof:}} (\(\Rightarrow \)) Veremos por ejemplo que \(L_{1}\) es \(\Sigma \)-p.r.. Sea \( (z_{1},...,z_{n},\zeta _{1},...,\zeta _{m})\) un elemento fijo de \( S_{1}\times ...\times S_{n}\times L_{1}\times ...\times L_{m}.\) Note que

  \(\displaystyle \alpha \in L_{1}\text{ sii }(z_{1},...,z_{n},\alpha ,\zeta _{2},...,\zeta _{m})\in S_{1}\times ...\times S_{n}\times L_{1}\times ...\times L_{m}, \)

  lo cual implica que
  \(\displaystyle \chi _{L_{1}}=\chi _{S_{1}\times ...\times S_{n}\times L_{1}\times ...\times L_{m}}\circ \left( C_{z_{1}}^{0,1},...,C_{z_{n}}^{0,1},p_{1}^{0,1},C_{\zeta _{2}}^{0,1},...,C_{\zeta _{m}}^{0,1}\right) . \)

  (\(\Leftarrow \)) Note que \(\chi _{S_{1}\times ...\times S_{n}\times L_{1}\times ...\times L_{m}}\) es el predicado
  \(\displaystyle \left( \chi _{S_{1}}\circ p_{1}^{n,m}\wedge ...\wedge \chi _{S_{n}}\circ p_{n}^{n,m}\wedge \chi _{L_{1}}\circ p_{n+1}^{n,m}\wedge ...\wedge \chi _{L_{m}}\circ p_{n+m}^{n,m}\right) . \)

  \(\Box\)


  \textbf{\underline{Lemma 32:}} Supongamos \(f:D_{f}\subseteq \omega ^{n}\times \Sigma ^{\ast m}\rightarrow O\) es \(\Sigma \)-p.r., donde \(O\in \{\omega ,\Sigma ^{\ast }\}.\) Si \(S\subseteq D_{f}\) es \(\Sigma \)-p.r., entonces \(f\mid _{S}\) es \(\Sigma \)-p.r..

  \textbf{\underline{Proof:}} Supongamos \(O=\Sigma ^{\ast }\). Entonces

  \(\displaystyle f\mid _{S}=\lambda x\alpha \left[ \alpha ^{x}\right] \circ \left( Suc\circ Pred\circ \chi _{S},f\right) \)

  es \(\Sigma \)-p.r.. El caso \(O=\omega \) es similar usando \(\lambda xy\left[ x^{y}\right] \) en lugar de \(\lambda x\alpha \left[ \alpha ^{x}\right] \). \(\Box\)


  \textbf{\underline{Lemma 33:}} Si \(f:D_{f}\subseteq \omega ^{n}\times \Sigma ^{\ast m}\rightarrow O\) es \(\Sigma \)-p.r., entonces existe una funcion \(\Sigma \) -p.r. \(\bar{f}:\omega ^{n}\times \Sigma ^{\ast m}\rightarrow O\), tal que \(f= \bar{f}\mid _{D_{f}}\).

  \textbf{\underline{Proof:}} Es facil ver por induccion en \(k\) que el enunciado se cumple para cada \(f\in \mathrm{PR}_{k}^{\Sigma }\) \(\Box\)


  \textbf{\underline{Proposition 34:}} Un conjunto \(S\) es \(\Sigma \)-p.r. sii \(S\) es el dominio de una funcion \(\Sigma \)-p.r.\(.\)

  \textbf{\underline{Proof:}} (\(\Rightarrow \)) Note que \(S=D_{Pred\circ \chi _{S}}.\)

  (\(\Leftarrow \)) Probaremos por induccion en \(k\) que \(D_{F}\) es \(\Sigma \) -p.r., para cada \(F\in \mathrm{PR}_{k}^{\Sigma }.\) El caso \(k=0\) es facil\(.\) Supongamos el resultado vale para un \(k\) fijo y supongamos \(F\in \mathrm{PR} _{k+1}^{\Sigma }.\) Veremos entonces que \(D_{F}\) es \(\Sigma \)-p.r.. Hay varios casos. Consideremos primero el caso en que \(F=R(f,g)\), donde

  \(\displaystyle \begin{array}{rcl} f & :& S_{1}\times ...\times S_{n}\times L_{1}\times ...\times L_{m}\rightarrow \Sigma ^{\ast } \\ g & :& \omega \times S_{1}\times ...\times S_{n}\times L_{1}\times ...\times L_{m}\times \Sigma ^{\ast }\rightarrow \Sigma ^{\ast }, \end{array} \)

  con \(S_{1},...,S_{n}\subseteq \omega \) y \(L_{1},...,L_{m}\subseteq \Sigma ^{\ast }\) conjuntos no vacios y \(f,g\in \mathrm{PR}_{k}^{\Sigma }\). Notese que por definicion de \(R(f,g)\), tenemos que
  \(\displaystyle D_{F}=\omega \times S_{1}\times ...\times S_{n}\times L_{1}\times ...\times L_{m}. \)

  Por hipotesis inductiva tenemos que \(D_{f}=S_{1}\times ...\times S_{n}\times L_{1}\times ...\times L_{m}\) es \(\Sigma \)-p.r., lo cual por el Lema 31 nos dice que los conjuntos \(S_{1},...,S_{n}\), \( L_{1},...,L_{m}\) son \(\Sigma \)-p.r.. Ya que \(\omega \) es \(\Sigma \)-p.r., el Lema 31 nos dice que \(D_{F}\) es \(\Sigma \)-p.r..
  Los otros casos de recursion primitiva son dejados al lector.

  Supongamos ahora que \(F=g\circ (g_{1},...,g_{n+m})\), donde

  \(\displaystyle \begin{array}{rcl} g & :& D_{g}\subseteq \omega ^{n}\times \Sigma ^{\ast m}\rightarrow O \\ g_{i} & :& D_{g_{i}}\subseteq \omega ^{k}\times \Sigma ^{\ast l}\rightarrow \omega \text{, }i=1,...,n \\ g_{i} & :& D_{g_{i}}\subseteq \omega ^{k}\times \Sigma ^{\ast l}\rightarrow \Sigma ^{\ast },i=n+1,...,n+m \end{array} \)

  estan en \(\mathrm{PR}_{k}^{\Sigma }.\) Por Lema 33, hay funciones \(\Sigma \)-p.r. \(\bar{g}_{1},...,\bar{g}_{n+m}\) las cuales son \( \Sigma \)-totales y cumplen
  \(\displaystyle g_{i}=\bar{g}_{i}\mid _{D_{g_{i}}}\text{, para }i=1,...,n+m. \)

  Por hipotesis inductiva los conjuntos \(D_{g}\), \(D_{g_{i}}\), \(i=1,...,n+m\), son \(\Sigma \)-p.r. y por lo tanto
  \(\displaystyle S=\bigcap_{i=1}^{n+m}D_{g_{i}} \)

  lo es. Notese que
  \(\displaystyle \chi _{D_{F}}=(\chi _{D_{g}}\circ \left( \bar{g}_{1},...,\bar{g} _{n+m}\right) \wedge \chi _{S}) \)

  lo cual nos dice que \(D_{F}\) es \(\Sigma \)-p.r.. \(\Box\)


  \textbf{\underline{Lemma 35:}} Supongamos \(f_{i}:D_{f_{i}}\subseteq \omega ^{n}\times \Sigma ^{\ast m}\rightarrow O\), \(i=1,...,k\), son funciones \(\Sigma \)-p.r. tales que \(D_{f_{i}}\cap D_{f_{j}}=\varnothing \) para \(i\neq j.\) Entonces \(f_{1}\cup ...\cup f_{k}\) es \(\Sigma \)-p.r..

  \textbf{\underline{Proof:}} Supongamos \(O=\Sigma ^{\ast }\) y \(k=2.\) Sean

  \(\displaystyle \bar{f}_{i}:\omega ^{n}\times \Sigma ^{\ast m}\rightarrow \Sigma ^{\ast },i=1,2, \)

  funciones \(\Sigma \)-p.r. tales que \(\bar{f}_{i}\mid _{D_{f_{i}}}=f_{i}\), \( i=1,2\) (Lema 33)\(.\) Por Lema 34 los conjuntos \(D_{f_{1}}\) y \(D_{f_{2}}\) son \(\Sigma \)-p.r. y por lo tanto lo es \( D_{f_{1}}\cup D_{f_{2}}\). Ya que
  \(\displaystyle f_{1}\cup f_{2}=\left( \lambda \alpha \beta \left[ \alpha \beta \right] \circ (\lambda x\alpha \left[ \alpha ^{x}\right] \circ (\chi _{D_{f_{1}}}, \bar{f}_{1}),\lambda x\alpha \left[ \alpha ^{x}\right] \circ (\chi _{D_{f_{2}}},\bar{f}_{2}))\right) \mid _{D_{f_{1}}\cup D_{f_{2}}} \)

  tenemos que \(f_{1}\cup f_{2}\) es \(\Sigma \)-p.r..
  El caso \(k >2\) puede probarse por induccion ya que

  \(\displaystyle f_{1}\cup ...\cup f_{k}=(f_{1}\cup ...\cup f_{k-1})\cup f_{k}. \)

  \(\Box\)


  \textbf{\underline{Corollary 36:}} Supongamos \(f\) es una funcion \(\Sigma \)-mixta cuyo dominio es finito. Entonces \(f\) es \(\Sigma \)-p.r..

  \textbf{\underline{Proof:}} Supongamos \(f:D_{f}\subseteq \omega ^{n}\times \Sigma ^{\ast m}\rightarrow O\) , con \(D_{f}=\{e_{1},...,e_{k}\}\). Por el Corolario 30, cada \( \{e_{i}\}\) es \(\Sigma \)-p.r. por lo cual el Lema 32 nos dice que \(C_{f(e_{i})}^{n,m}\mid _{\{e_{1}\}}\) es \(\Sigma \)-p.r.. O sea que

  \(\displaystyle f=C_{f(e_{1})}^{n,m}\mid _{\{e_{1}\}}\cup ...\cup C_{f(e_{k})}^{n,m}\mid _{\{e_{k}\}} \)

  es \(\Sigma \)-p.r.. \(\Box\)


  \textbf{\underline{Lemma 37:}} \(\lambda i\alpha \left[ \lbrack \alpha ]_{i}\right] \) es \(\Sigma \)-p.r..

  \textbf{\underline{Proof:}} Note que

  \(\displaystyle \begin{array}{rcl} \lbrack \varepsilon ]_{i} & =& \varepsilon \\ \lbrack \alpha a]_{i} & =& \left\{ \begin{array}{lll} \lbrack \alpha ]_{i} & & \text{si }i\neq \left\vert \alpha \right\vert +1 \\ a & & \text{si }i=\left\vert \alpha \right\vert +1 \end{array} \right. \end{array} \)

  lo cual dice que \(\lambda i\alpha \left[ \lbrack \alpha ]_{i}\right] =R\left( C_{\varepsilon }^{1,0},\mathcal{G}\right) \), donde \(\mathcal{G} _{a}:\omega \times \Sigma ^{\ast }\times \Sigma ^{\ast }\rightarrow \Sigma ^{\ast }\) es dada por
  \(\displaystyle \mathcal{G}_{a}(i,\alpha ,\zeta )=\left\{ \begin{array}{lll} \zeta & & \text{si }i\neq \left\vert \alpha \right\vert +1 \\ a & & \text{si }i=\left\vert \alpha \right\vert +1 \end{array} \right. \)

  O sea que solo resta probar que cada \(\mathcal{G}_{a}\) es \(\Sigma \)-p.r.. Primero note que los conjuntos
  \(\displaystyle \begin{array}{rcl} S_{1} & =& \left\{ (i,\alpha ,\zeta )\in \omega \times \Sigma ^{\ast }\times \Sigma ^{\ast }:i\neq \left\vert \alpha \right\vert +1\right\} \\ S_{2} & =& \left\{ (i,\alpha ,\zeta )\in \omega \times \Sigma ^{\ast }\times \Sigma ^{\ast }:i=\left\vert \alpha \right\vert +1\right\} \end{array} \)

  son \(\Sigma \)-p.r. ya que
  \(\displaystyle \begin{array}{rcl} \chi _{S_{1}} & =& \lambda xy\left[ x\neq y\right] \circ \left( p_{1}^{1,2},Suc\circ \lambda \alpha \left[ \left\vert \alpha \right\vert \right] \circ p_{2}^{1,2}\right) \\ \chi _{S_{2}} & =& \lambda xy\left[ x=y\right] \circ \left( p_{1}^{1,2},Suc\circ \lambda \alpha \left[ \left\vert \alpha \right\vert \right] \circ p_{2}^{1,2}\right) . \end{array} \)

  Ya que
  \(\displaystyle \mathcal{G}_{a}=p_{3}^{1,2}\mid _{S_{1}}\cup C_{a}^{1,2}\mid _{S_{2}}, \)

  el Lema 35 nos dice que \(\mathcal{G}_{a}\) es \(\Sigma \)-p.r., para cada \(a\in \Sigma \). \(\Box\)


  \textbf{\underline{Lemma 38:}} Sean \(n,m\geq 0\).
  (a) Si \(f:\omega \times S_{1}\times ...\times S_{n}\times L_{1}\times ...\times L_{m}\rightarrow \omega \) es \(\Sigma \)-p.r., con \( S_{1},...,S_{n}\subseteq \omega \) y \(L_{1},...,L_{m}\subseteq \Sigma ^{\ast } \) no vacios, entonces lo son las funciones \(\lambda xy\vec{x}\vec{\alpha} \left[ \sum_{t=x}^{t=y}f(t,\vec{x},\vec{\alpha})\right] \) y \(\lambda xy\vec{x }\vec{\alpha}\left[ \prod_{t=x}^{t=y}f(t,\vec{x},\vec{\alpha})\right] \).
  (b) Si \(f:\omega \times S_{1}\times ...\times S_{n}\times L_{1}\times ...\times L_{m}\rightarrow \Sigma ^{\ast }\) es \(\Sigma \)-p.r., con \( S_{1},...,S_{n}\subseteq \omega \) y \(L_{1},...,L_{m}\subseteq \Sigma ^{\ast } \) no vacios, entonces lo es la funcion \(\lambda xy\vec{x}\vec{\alpha}\left[ \subset _{t=x}^{t=y}f(t,\vec{x},\vec{\alpha})\right] \)

  \textbf{\underline{Proof:}} (a) Sea \(G=\lambda tx\vec{x}\vec{\alpha}\left[ \sum_{i=x}^{i=t}f(i,\vec{x}, \vec{\alpha})\right] \). Ya que

  \(\displaystyle \lambda xy\vec{x}\vec{\alpha}\left[ \sum_{i=x}^{i=y}f(i,\vec{x},\vec{\alpha}) \right] =G\circ \left( p_{2}^{n+2,m},p_{1}^{n+2,m},p_{3}^{n+2,m},...,p_{n+m+2}^{n+2,m}\right) \)

  solo tenemos que probar que \(G\) es \(\Sigma \)-p.r.. Primero note que
  \(\displaystyle \begin{array}{rcl} G(0,x,\vec{x},\vec{\alpha}) & =& \left\{ \begin{array}{lll} 0 & & \text{si }x >0 \\ f(0,\vec{x},\vec{\alpha}) & & \text{si }x=0 \end{array} \right. \\ G(t+1,x,\vec{x},\vec{\alpha}) & =& \left\{ \begin{array}{lll} 0 & & \text{si }x >t+1 \\ G(t,x,\vec{x},\vec{\alpha})+f(t+1,\vec{x},\vec{\alpha}) & & \text{si }x\leq t+1 \end{array} \right. \end{array} \)

  Sean
  \(\displaystyle \begin{array}{rcl} D_{1} & =& \left\{ (x,\vec{x},\vec{\alpha})\in \omega \times S_{1}\times ...\times S_{n}\times L_{1}\times ...\times L_{m}:x >0\right\} \\ D_{2} & =& \left\{ (x,\vec{x},\vec{\alpha})\in \omega \times S_{1}\times ...\times S_{n}\times L_{1}\times ...\times L_{m}:x=0\right\} \\ H_{1} & =& \left\{ (z,t,x,\vec{x},\vec{\alpha})\in \omega ^{3}\times S_{1}\times ...\times S_{n}\times L_{1}\times ...\times L_{m}:x >t+1\right\} \\ H_{2} & =& \left\{ (z,t,x,\vec{x},\vec{\alpha})\in \omega ^{3}\times S_{1}\times ...\times S_{n}\times L_{1}\times ...\times L_{m}:x\leq t+1\right\} . \end{array} \)

  Es facil de chequear que estos conjuntos son \(\Sigma \)-p.r.. Veamos que por ejemplo \(H_{1}\) lo es. Es decir debemos ver que \(\chi _{H_{1}}\) es \(\Sigma \) -p.r.. Ya que \(f\) es \(\Sigma \)-p.r. tenemos que \(D_{f}=\omega \times S_{1}\times ...\times S_{n}\times L_{1}\times ...\times L_{m}\) es \(\Sigma \) -p.r., lo cual por el Lema 31 nos dice que los conjuntos \( S_{1},...,S_{n}\), \(L_{1},...,L_{m}\) son \(\Sigma \)-p.r.. Ya que \(\omega \) es \( \Sigma \)-p.r., el Lema 31 nos dice que \(R=\omega ^{3}\times S_{1}\times ...\times S_{n}\times L_{1}\times ...\times L_{m}\) es \(\Sigma \)-p.r.. Notese que \(\chi _{H_{1}}=(\chi _{R}\wedge \lambda ztx\vec{x} \vec{\alpha}\left[ x >t+1\right] )\) por cual \(\chi _{H_{1}}\) es \(\Sigma \) -p.r. ya que es la conjuncion de dos predicados \(\Sigma \)-p.r.
  Ademas note que \(G=R(h,g)\), donde

  \(\displaystyle \begin{array}{rcl} h & =& C_{0}^{n+1,m}\mid _{D_{1}}\cup \lambda x\vec{x}\vec{\alpha}\left[ f(0, \vec{x},\vec{\alpha})\right] \mid _{D_{2}} \\ g & =& C_{0}^{n+3,m}\mid _{H_{1}}\cup \lambda ztx\vec{x}\vec{\alpha}\left[ z+f(t+1,\vec{x},\vec{\alpha})\right] )\mid _{H_{2}} \end{array} \)

  O sea que los Lemas 35 y 32 garantizan que \(G\) es \( \Sigma \)-p.r.. \(\Box\)


  \textbf{\underline{Lemma 39:}} Sean \(n,m\geq 0\).
  (a) Sea \(P:S\times S_{1}\times ...\times S_{n}\times L_{1}\times ...\times L_{m}\rightarrow \omega \) un predicado \(\Sigma \)-p.r. y supongamos \(\bar{S}\subseteq S\) es \(\Sigma \)-p.r.. Entonces \(\lambda x\vec{x}\vec{\alpha }\left[ (\forall t\in \bar{S})_{t\leq x}\;P(t,\vec{x},\vec{\alpha})\right] \) y \(\lambda x\vec{x}\vec{\alpha}\left[ (\exists t\in \bar{S})_{t\leq x}\;P(t, \vec{x},\vec{\alpha})\right] \) son predicados \(\Sigma \)-p.r.. (Note que el dominio de estos predicados es \(\omega \times S_{1}\times ...\times S_{n}\times L_{1}\times ...\times L_{m}\))
  (b) Sea \(P:S_{1}\times ...\times S_{n}\times L_{1}\times ...\times L_{m}\times L\rightarrow \omega \) un predicado \(\Sigma \)-p.r. y supongamos \( \bar{L}\subseteq L\) es \(\Sigma \)-p.r.. Entonces \(\lambda x\vec{x}\vec{\alpha} \left[ (\forall \alpha \in \bar{L})_{\left\vert \alpha \right\vert \leq x}\;P(\vec{x},\vec{\alpha},\alpha )\right] \) y \(\lambda x\vec{x}\vec{\alpha} \left[ (\exists \alpha \in \bar{L})_{\left\vert \alpha \right\vert \leq x}\;P(\vec{x},\vec{\alpha},\alpha )\right] \) son predicados \(\Sigma \)-p.r..

  \textbf{\underline{Proof:}} (a) Sea

  \(\displaystyle \bar{P}=P\mid _{\bar{S}\times S_{1}\times ...\times S_{n}\times L_{1}\times ...\times L_{m}}\cup C_{1}^{1+n,m}\mid _{(\omega -\bar{S})\times S_{1}\times ...\times S_{n}\times L_{1}\times ...\times L_{m}} \)

  Notese que \(\bar{P}\) es \(\Sigma \)-p.r.. Ya que
  \(\displaystyle \begin{array}{rcl} \lambda x\vec{x}\vec{\alpha}\left[ (\forall t\in \bar{S})_{t\leq x}P(t,\vec{x },\vec{\alpha})\right] & =& \lambda x\vec{x}\vec{\alpha}\left[ \prod\limits_{t=0}^{t=x}\bar{P}(t,\vec{x},\vec{\alpha})\right] \\ & =& \lambda xy\vec{x}\vec{\alpha}\left[ \prod\limits_{t=x}^{t=y}\bar{P}(t, \vec{x},\vec{\alpha})\right] \circ \left( C_{0}^{1+n,m},p_{1}^{1+n,m},...,p_{1+n+m}^{1+n,m}\right) \end{array} \)

  el Lema 38 implica que \(\lambda x\vec{x}\vec{\alpha}\left[ (\forall t\in \bar{S})_{t\leq x}\;P(t,\vec{x},\vec{\alpha})\right] \) es \( \Sigma \)-p.r..
  Finalmente note que

  \(\displaystyle \lambda x\vec{x}\vec{\alpha}\left[ (\exists t\in \bar{S})_{t\leq x}\;P(t, \vec{x},\vec{\alpha})\right] =\lnot \lambda x\vec{x}\vec{\alpha}\left[ (\forall t\in \bar{S})_{t\leq x}\;\lnot P(t,\vec{x},\vec{\alpha})\right] \)

  es \(\Sigma \)-p.r..
  (b) Sea \(< \) un orden total estricto sobre \(\Sigma .\) Sea \(k\) el cardinal de \( \Sigma \). Ya que

  \(\displaystyle \left\vert \alpha \right\vert \leq x\text{ sii }\#^{< }(\alpha )\leq \sum_{\iota =1}^{i=x}k^{i}, \)

  (ejercicio) tenemos que
  \(\displaystyle \lambda x\vec{x}\vec{\alpha}\left[ (\forall \alpha \in \bar{L})_{\left\vert \alpha \right\vert \leq x}P(\vec{x},\vec{\alpha},\alpha )\right] =\lambda x \vec{x}\vec{\alpha}\left[ (\forall t\in \#^{< }(\bar{L}))_{t\leq \sum_{\iota =1}^{i=x}k^{i}}P(\vec{x},\vec{\alpha},\ast ^{< }(t))\right] \)

  Sea \(H=\lambda t\vec{x}\vec{\alpha}\left[ P(\vec{x},\vec{\alpha},\ast ^{< }(t))\right] .\) Notese que \(H\) es \(\Sigma \)-p.r. y
  \(\displaystyle D_{H}=\#^{< }(L)\times S_{1}\times ...\times S_{n}\times L_{1}\times ...\times L_{m} \)

  Ademas note que \(\#^{< }(\bar{L})\) es \(\Sigma \)-p.r. (ejercicio), lo cual por (a) implica que
  \(\displaystyle Q=\lambda x\vec{x}\vec{\alpha}\left[ (\forall t\in \#^{< }(\bar{L}))_{t\leq x}H(t,\vec{x},\vec{\alpha})\right] \)

  es \(\Sigma \)-p.r.. O sea que
  \(\displaystyle \lambda x\vec{x}\vec{\alpha}\left[ (\forall \alpha \in \bar{L})_{\left\vert \alpha \right\vert \leq x}\;P(\vec{x},\vec{\alpha},\alpha )\right] =Q\circ \left( \lambda x\vec{x}\vec{\alpha}\left[ \sum\limits_{\iota =1}^{i=x}k^{i} \right] ,p_{1}^{1+n,m},...,p_{1+n+m}^{1+n,m}\right) \)

  es \(\Sigma \)-p.r.. \(\Box\)


  \textbf{\underline{Lemma 40:}}
  (a) El predicado \(\lambda xy\left[ x\text{ divide }y\right] \) es \( \varnothing \)-p.r..
  (b) El predicado \(\lambda x\left[ x\text{ es primo}\right] \) es \( \varnothing \)-p.r..
  (c) El predicado \(\lambda \alpha \beta \left[ \alpha \text{\ }\mathrm{ inicial}\ \beta \right] \) es \(\Sigma \)-p.r..

  \textbf{\underline{Proof:}} (a) Si tomamos \(P=\lambda tx_{1}x_{2}\left[ x_{2}=t.x_{1}\right] \in \mathrm{PR}^{\varnothing }\), tenemos que

  \(\displaystyle \begin{array}{rcl} \lambda x_{1}x_{2}\left[ x_{1}\text{ divide }x_{2}\right] & =& \lambda x_{1}x_{2}\left[ (\exists t\in \omega )_{t\leq x_{2}}\;P(t,x_{1},x_{2}) \right] \\ & =& \lambda xx_{1}x_{2}\left[ (\exists t\in \omega )_{t\leq x}\;P(t,x_{1},x_{2})\right] \circ \left( p_{2}^{2,0},p_{1}^{2,0},p_{2}^{2,0}\right) \end{array} \)

  por lo que podemos aplicar el lema anterior.
  (b) Ya que

  \(\displaystyle x\text{ es primo sii }x >1\wedge \left( (\forall t\in \omega )_{t\leq x}\;t=1\vee t=x\vee \lnot (t\text{ divide }x)\right) \)

  podemos usar un argumento similar al de la prueba de (a).
  (c) es dejado al lector. \(\Box\)

\section{Minimización y funciones $\Sigma$-recursivas}

\textbf{\underline{Lemma 41:}} Si \(P:D_{P}\subseteq \omega \times \omega ^{n}\times \Sigma ^{\ast m}\rightarrow \omega \) es un predicado \(\Sigma \)-efectivamente computable y \( D_{P}\) es \(\Sigma \)-efectivamente computable, entonces la funcion \(M(P)\) es \( \Sigma \)-efectivamente computable.

\textbf{\underline{Proof:}} Ejercicio \(\Box\)

\textbf{\underline{Theorem 42:}} Si \(f\in \mathrm{R} ^{\Sigma }\), entonces \(f\) es \(\Sigma \)-efectivamente computable.

\textbf{\underline{Proof:}} Dejamos al lector la prueba por induccion en \(k\) de que si \(f\in \mathrm{R} _{k}^{\Sigma }\), entonces \(f\) es \(\Sigma \)-efectivamente computable. \(\Box\)

\textbf{\underline{Lemma 43:}} Sean \(n,m\geq 0\). Sea \(P:D_{P}\subseteq \omega \times \omega ^{n}\times \Sigma ^{\ast m}\rightarrow \omega \) un predicado \(\Sigma \) -p.r.. Entonces
(a) \(M(P)\) es \(\Sigma \)-recursiva.
(b) Si hay una funcion \(\Sigma \)-p.r. \(f:\omega ^{n}\times \Sigma ^{\ast m}\rightarrow \omega \) tal que
\(\displaystyle M(P)(\vec{x},\vec{\alpha})=\min_{t}P(t,\vec{x},\vec{\alpha})\leq f(\vec{x}, \vec{\alpha})\text{, para cada }(\vec{x},\vec{\alpha})\in D_{M(P)}\text{,} \)
entonces \(M(P)\) es \(\Sigma \)-p.r..

\textbf{\underline{Proof:}} (a) Sea \(\bar{P}=P\mid _{D_{P}}\cup C_{0}^{n+1,m}\mid _{(\omega ^{n+1}\times \Sigma ^{\ast m})-D_{P}}\). Dejamos al lector verificar cuidadosamente que \( M(P)=M(\bar{P})\). Veremos entonces que \(M(\bar{P})\) es \(\Sigma \)-recursiva. Note que \(\bar{P}\) es \(\Sigma \)-p.r. (por que?). Sea \(k\) tal que \(\bar{P}\in \mathrm{PR}_{k}^{\Sigma }\). Ya que \(\bar{P}\) es \(\Sigma \)-total y \(\bar{P} \in \mathrm{PR}_{k}^{\Sigma }\subseteq \mathrm{R}_{k}^{\Sigma }\), tenemos que \(M(\bar{P})\in \mathrm{R}_{k+1}^{\Sigma }\) y por lo tanto \(M(\bar{P})\in \mathrm{R}^{\Sigma }\).

(b) Primero veremos que \(D_{M(\bar{P})}\) es un conjunto \(\Sigma \)-p.r.. Notese que

\(\displaystyle \chi _{D_{M(\bar{P})}}=\lambda \vec{x}\vec{\alpha}\left[ (\exists t\in \omega )_{t\leq f(\vec{x},\vec{\alpha})}\;\bar{P}(t,\vec{x},\vec{\alpha}) \right] \)

lo cual nos dice que
\(\displaystyle \chi _{D_{M(\bar{P})}}=\lambda x\vec{x}\vec{\alpha}\left[ (\exists t\in \omega )_{t\leq x}\;\bar{P}(t,\vec{x},\vec{\alpha})\right] \circ (f,p_{1}^{n,m},...,p_{n+m}^{n,m}) \)

Pero el Lema 39 nos dice que \(\lambda x\vec{x}\vec{\alpha} \left[ (\exists t\in \omega )_{t\leq x}\;\bar{P}(t,\vec{x},\vec{\alpha}) \right] \) es \(\Sigma \)-p.r. por lo cual tenemos que \(\chi _{D_{M(\bar{P})}}\) lo es.
Sea

\(\displaystyle P_{1}=\lambda t\vec{x}\vec{\alpha}\left[ \bar{P}(t,\vec{x},\vec{\alpha} )\wedge (\forall j\in \omega )_{j\leq t}\;j=t\vee \lnot \bar{P}(j,\vec{x}, \vec{\alpha})\right] \)

Note que \(P_{1}\) es \(\Sigma \)-total. Dejamos al lector usando lemas anteriores probar que \(P_{1}\) es \(\Sigma \)-p.r.. Ademas notese que para cada \((\vec{x},\vec{\alpha})\in \omega ^{n}\times \Sigma ^{\ast m}\) tenemos que
\(\displaystyle P_{1}(t,\vec{x},\vec{\alpha})=1\text{ si y solo si }t=M(\bar{P})(\vec{x}, \vec{\alpha}) \)

Esto nos dice que
\(\displaystyle M(\bar{P})=\left( \lambda \vec{x}\vec{\alpha}\left[ \prod_{t=0}^{f(\vec{x}, \vec{\alpha})}t^{P_{1}(t,\vec{x},\vec{\alpha})}\right] \right) \mid _{D_{M( \bar{P})}} \)

por lo cual para probar que \(M(\bar{P})\) es \(\Sigma \)-p.r. solo nos resta probar que
\(\displaystyle F=\lambda \vec{x}\vec{\alpha}\left[ \prod_{t=0}^{f(\vec{x},\vec{\alpha} )}t^{P_{1}(t,\vec{x},\vec{\alpha})}\right] \)

lo es. Pero
\(\displaystyle F=\lambda xy\vec{x}\vec{\alpha}\left[ \prod_{t=x}^{y}t^{P_{1}(t,\vec{x},\vec{ \alpha})}\right] \circ (C_{0}^{n,m},f,p_{1}^{n,m},...,p_{n+m}^{n,m}) \)

y por lo tanto el Lema 38 nos dice que \(F\) es \(\Sigma \)-p.r.. De esta manera hemos probado que \(M(\bar{P})\) es \(\Sigma \)-p.r. y por lo tanto \(M(P)\) lo es. \(\Box\)

\textbf{\underline{Lemma 44:}} Las siguientes funciones son \(\varnothing \)-p.r.:
(a) \( \begin{array}{rll} Q:\omega \times \mathbf{N} & \rightarrow & \omega \\ (x,y) & \rightarrow & \text{cociente de la division de }x\text{ por }y \end{array} \)
(b) \( \begin{array}{rll} R:\omega \times \mathbf{N} & \rightarrow & \omega \\ (x,y) & \rightarrow & \text{resto de la division de }x\text{ por }y \end{array} \)
(c) \( \begin{array}{rll} pr:\mathbf{N} & \rightarrow & \omega \\ n & \rightarrow & n\text{-esimo numero primo} \end{array} \)


\textbf{\underline{Proof:}} (a) Veamos primero veamos que \(Q=M(P)\), donde \(P=\lambda txy\left[ (t+1).y >x \right] \). Notar que

\(\displaystyle \begin{array}{rcl} D_{M(P)} & =& \{(x,y):(\exists t\in \omega )\;P(t,x,y)=1\} \\ & =& \{(x,y):(\exists t\in \omega )\;(t+1).y >x\} \\ & =& \omega \times \mathbf{N} \\ & =& D_{Q} \end{array} \)

Dejamos al lector la facil verificacion de que para cada \((x,y)\in \omega \times \mathbf{N}\), se tiene que
\(\displaystyle Q(x,y)=M(P)(x,y)=\min_{t}(t+1).y >x \)

Esto prueba que \(Q=M(P)\). Ya que \(P\) es \(\varnothing \)-p.r. y
\(\displaystyle Q(x,y)\leq p_{1}^{2,0}(x,y),\text{para cada }(x,y)\in \omega \times \mathbf{N } \)

(b) del Lema 43 implica que \(Q\in \mathrm{PR}^{\varnothing }\).
(b) Notese que

\(\displaystyle R=\lambda xy\left[ x\dot{-}Q(x,y).y\right] \)

y por lo tanto \(R\in \mathrm{PR}^{\varnothing }\).
(c) Para ver que \(pr\) es \(\varnothing \)-p.r., veremos que la extension \( h:\omega \rightarrow \omega \), dada por \(h(0)=0\) y \(h(n)=pr(n)\), \(n\geq 1\), es \(\varnothing \)-p.r.. Primero note que

\(\displaystyle \begin{array}{rcl} h(0) & =& 0 \\ h(x+1) & =& \min\nolimits_{t}\left( t\text{ es primo}\wedge t >h(x)\right) \end{array} \)

O sea que \(h=R\left( C_{0}^{0,0},M(P)\right) \), donde
\(\displaystyle P=\lambda tzx\left[ t\text{ es primo}\wedge t >z\right] \)

Es decir que solo nos resta ver que \(M(P)\) es \(\varnothing \)-p.r.. Claramente \( P\) es \(\varnothing \)-p.r.. Veamos que para cada \((z,x)\in \omega ^{2}\), tenemos que
\(\displaystyle M(P)(z,x)=\min\nolimits_{t}\left( t\text{ es primo}\wedge t >z\right) \leq z!+1 \)

Sea \(p\) primo tal que \(p\) divide a \(z!+1\). Es facil ver que entonces \(p >z\). Pero esto claramente nos dice que
\(\displaystyle \min\nolimits_{t}\left( t\text{ es primo}\wedge t >z\right) \leq p\leq z!+1 \)

O sea que (b) del Lema 43 implica que \(M(P)\) es \(\varnothing \) -p.r. ya que podemos tomar \(f=\lambda zx\left[ z!+1\right] \). \(\Box\)


\textbf{\underline{Lemma 45:}} Las funciones \(\lambda xi\left[ (x)_{i}\right] \) y \(\lambda x\left[ Lt(x) \right] \) son \(\varnothing \)-p.r.

\textbf{\underline{Proof:}} Note que \(D_{\lambda xi\left[ (x)_{i}\right] }=\mathbf{N}\times \mathbf{N}\). Sea

\(\displaystyle P=\lambda txi\left[ \lnot (pr(i)^{t+1}\ \text{divide }x)\right] \)

Note que \(P\) es \(\varnothing \)-p.r. y que \(D_{P}=\omega \times \omega \times \mathbf{N}\). Dejamos al lector la prueba de que \(\lambda xi\left[ (x)_{i} \right] =M(P)\). Ya que \((x)_{i}\leq x\), para todo \(x\in \mathbf{N}\), (b) del Lema 43 implica que \(\lambda xi\left[ (x)_{i}\right] \) es \( \varnothing \)-p.r..
Veamos que \(\lambda x\left[ Lt(x)\right] \) es \(\varnothing \)-p.r.. Sea

\(\displaystyle Q=\lambda tx\left[ (\forall i\in \mathbf{N})_{i\leq x}\;(i\leq t\vee (x)_{i}=0)\right] \)

Notese que \(D_{Q}=\omega \times \mathbf{N}\) y que ademas por el Lema 39 tenemos que \(Q\) es \(\varnothing \)-p.r. (dejamos al lector explicar como se aplica tal lema en este caso). Ademas notese que \(\lambda x \left[ Lt(x)\right] =M(Q)\) y que
\(\displaystyle Lt(x)\leq x,\text{para todo }x\in \mathbf{N} \)

lo cual por (b) del Lema 43 nos dice que \(\lambda x\left[ Lt(x)\right] \) es \(\varnothing \)-p.r.. \(\Box\)
Para \(x_{1},...,x_{n}\in \omega \), escribiremos \(\left\langle x_{1},...,x_{n}\right\rangle \) en lugar de \(\left\langle x_{1},...,x_{n},0,...\right\rangle \).



\textbf{\underline{Lemma 46:}} Sea \(n\geq 1\). La funcion \(\lambda x_{1}...x_{n}\left[ \left\langle x_{1},...,x_{n}\right\rangle \right] \) es \(\varnothing \)-p.r.

\textbf{\underline{Proof:}} Sea \(f_{n}=\lambda x_{1}...x_{n}\left[ \left\langle x_{1},...,x_{n}\right\rangle \right] \). Claramente \(f_{1}\) es \(\varnothing \) -p.r.. Ademas note que para cada \(n\geq 1\), tenemos

\(\displaystyle f_{n+1}=\lambda x_{1}...x_{n+1}\left[ \left( f_{n}(x_{1},...,x_{n})pr(n+1)^{x_{n+1}}\right) \right] \text{.} \)

O sea que podemos aplicar un argumento inductivo. \(\Box\)


\textbf{\underline{Lemma 47:}} Supongamos que \(\Sigma \neq \varnothing \). Sea \(< \) un orden total estricto sobre \(\Sigma \), sean \(n,m\geq 0\) y sea \( P:D_{P}\subseteq \omega ^{n}\times \Sigma ^{\ast m}\times \Sigma ^{\ast }\rightarrow \omega \) un predicado \(\Sigma \)-p.r.. Entonces
(a) \(M^{< }(P)\) es \(\Sigma \)-recursiva.
(b) Si existe una funcion \(\Sigma \)-p.r. \(f:\omega ^{n}\times \Sigma ^{\ast m}\rightarrow \omega \) tal que
\(\displaystyle \left\vert M^{< }(P)(\vec{x},\vec{\alpha})\right\vert =\left\vert \min\nolimits_{\alpha }^{< }P(\vec{x},\vec{\alpha},\alpha )\right\vert \leq f( \vec{x},\vec{\alpha})\text{, para cada }(\vec{x},\vec{\alpha})\in D_{M^{< }(P)}\text{,} \)
entonces \(M^{< }(P)\) es \(\Sigma \)-p.r..

\textbf{\underline{Proof:}} Sea \(Q=P\circ \left( p_{2}^{1+n,m},...,p_{1+n+m}^{1+n,m},\ast ^{< }\circ p_{1}^{1+n,m}\right) \). Note que

\(\displaystyle M^{< }(P)=\ast ^{< }\circ M(Q) \)

lo cual por (a) del Lema 43 implica que \(M^{< }(P)\) es \( \Sigma \)-recursiva.
Sea \(k\) el cardinal de \(\Sigma \). Ya que

\(\displaystyle \left\vert \ast ^{< }(M(Q)(\vec{x},\vec{\alpha}))\right\vert =\left\vert M^{< }(P)(\vec{x},\vec{\alpha})\right\vert \leq f(\vec{x},\vec{\alpha})\text{, } \)

para todo \((\vec{x},\vec{\alpha})\in D_{M^{< }(P)}=D_{M(Q)}\), tenemos que
\(\displaystyle M(Q)(\vec{x},\vec{\alpha}))\leq \sum_{\iota =1}^{i=f(\vec{x},\vec{\alpha} )}k^{i}\text{, para cada }(\vec{x},\vec{\alpha})\in D_{M(Q)}\text{.} \)

O sea que por (a) del Lema 43, \(M(Q)\) es \(\Sigma \)-p.r. y por lo tanto \(M^{< }(P)\) lo es. \(\Box\)

\subsection{Recursion primitiva sobre valores anteriores}

\textbf{\underline{Lemma 48:}} Supongamos
\(\displaystyle \begin{array}{rcl} f & :& U\subseteq \omega ^{n}\times \Sigma ^{\ast m}\rightarrow \omega \\ g & :& \omega \times \omega \times U\rightarrow \omega \\ h & :& \omega \times U\rightarrow \omega \end{array} \)

son funciones tales que
\(\displaystyle \begin{array}{rcl} h(0,\vec{x},\vec{\alpha}) & =& f(\vec{x},\vec{\alpha})\text{, para cada }(\vec{ x},\vec{\alpha})\in U \\ h(x+1,\vec{x},\vec{\alpha}) & =& g(h^{\downarrow }(x,\vec{x},\vec{\alpha}),x, \vec{x},\vec{\alpha})\text{, para cada }x\in \omega \text{ y }(\vec{x},\vec{ \alpha})\in U\text{.} \end{array} \)
Entonces \(h\) es \(\Sigma \)-p.r. si \(f\) y \(g\) lo son.

\textbf{\underline{Proof:}} Supongamos \(f,g\) son \(\Sigma \)-p.r.. Primero veremos que \(h^{\downarrow }\) es \(\Sigma \)-p.r.. Notese que

\(\displaystyle \begin{array}{rcl} h^{\downarrow }(0,\vec{x},\vec{\alpha}) & =& \left\langle h(0,\vec{x},\vec{ \alpha})\right\rangle \\ & =& \left\langle f(\vec{x},\vec{\alpha})\right\rangle \\ & =& 2^{f(\vec{x},\vec{\alpha})} \\ h^{\downarrow }(x+1,\vec{x},\vec{\alpha}) & =& h^{\downarrow }(x,\vec{x},\vec{ \alpha})pr(x+2)^{h(x+1,\vec{x},\vec{\alpha})} \\ & =& h^{\downarrow }(x,\vec{x},\vec{\alpha})pr(x+2)^{g(h^{\downarrow }(x,\vec{x },\vec{\alpha}),x,\vec{x},\vec{\alpha})} \end{array} \)

lo cual nos dice que \(h^{\downarrow }=R(f_{1},g_{1})\) donde
\(\displaystyle \begin{array}{rcl} f_{1} & =& \lambda \vec{x}\vec{\alpha}\left[ 2^{f(\vec{x},\vec{\alpha})}\right] \\ g_{1} & =& \lambda Ax\vec{x}\vec{\alpha}\left[ Apr(x+2)^{g(A,x,\vec{x},\vec{ \alpha})}\right] \end{array} \)

O sea que \(h^{\downarrow }\) es \(\Sigma \)-p.r. ya que \(f_{1}\) y \(g_{1}\) lo son. Finalmente notese que
\(\displaystyle h=\lambda ix[(x)_{i}]\circ (Suc\circ p_{1}^{1+n,m},h^{\downarrow }) \)

lo cual nos dice que \(h\) es \(\Sigma \)-p.r.. \(\Box\)

\subsection{Independencia del alfabeto}

\textbf{\underline{Lemma 49:}} Supongamos \(\varnothing \neq \Sigma \subseteq \Gamma \).
(a) Si \(< \) es un orden total estricto sobre \(\Sigma \), entonces las funciones \(\ast ^{< }:\omega \rightarrow \Sigma ^{\ast }\) y \(\#^{< }:\Sigma ^{\ast }\rightarrow \omega \) son \(\Gamma \)-p.r..
(b) Si \(\prec \) es un orden total estricto sobre \(\Gamma \), entonces las funciones \(\#^{\prec }\mid _{\Sigma ^{\ast }}:\Sigma ^{\ast }\rightarrow \omega \) y \(\ast ^{\prec }\mid _{\#^{\prec }(\Sigma ^{\ast })}:\#^{\prec }(\Sigma ^{\ast })\rightarrow \Sigma ^{\ast }\) son \(\Sigma \)-p.r..

\textbf{\underline{Proof:}} (a) Supongamos \(\Sigma =\{a_{1},...,a_{k}\}\) y \(< \) es dado por \( a_{1}< ...< a_{k}\). Sea \(s_{e}^{< }:\Gamma ^{\ast }\rightarrow \Gamma ^{\ast }\) dada por

\(\displaystyle \begin{array}{rcl} s_{e}^{< }(\varepsilon ) & =& a_{1} \\ s_{e}^{< }(\alpha a_{i}) & =& \alpha a_{i+1}\text{, si }i< k \\ s_{e}^{< }(\alpha a_{k}) & =& s_{e}^{< }(\alpha )a_{1} \\ s_{e}^{< }(\alpha a) & =& \varepsilon \text{, si }a\in \Gamma -\Sigma . \end{array} \)

Note que \(s_{e}^{< }\) es \(\Gamma \)-p.r. y que \(s_{e}^{< }\mid _{\Sigma ^{\ast }}=s^{< }\). Ya que \(\Sigma ^{\ast }\) es un conjunto \(\Gamma \)-p.r. tenemos que \(s^{< }\) es \(\Gamma \)-p.r.. O sea que la recursion
\(\displaystyle \begin{array}{rcl} \ast ^{< }(0) & =& \varepsilon \\ \ast ^{< }(x+1) & =& s^{< }(\ast ^{< }(x)) \end{array} \)

implica que \(\ast ^{< }\) es \(\Gamma \)-p.r..
Para ver que \(\#^{< }:\Sigma ^{\ast }\rightarrow \omega \) es \(\Gamma \)-p.r., sea \(\#_{e}^{< }:\Gamma ^{\ast }\rightarrow \omega \) dada por

\(\displaystyle \begin{array}{rcl} \#_{e}^{< }(\varepsilon ) & =& 0 \\ \#_{e}^{< }(\alpha a_{i}) & =& \#_{e}^{< }(\alpha ).k+i \\ \#_{e}^{< }(\alpha a) & =& 0\text{, si }a\in \Gamma -\Sigma . \end{array} \)

Ya que \(\#_{e}^{< }\) es \(\Gamma \)-p.r., eso es \(\#^{< }=\#_{e}^{< }\mid _{\Sigma ^{\ast }}\).
(b) Sea \(n\) el cardinal de \(\Gamma .\) Ya que

\(\displaystyle \begin{array}{rcl} \#^{\prec } & \mid & _{\Sigma ^{\ast }}(\varepsilon )=0 \\ \#^{\prec } & \mid & _{\Sigma ^{\ast }}(\alpha a)=\#^{\prec }\mid _{\Sigma ^{\ast }}(\alpha ).n+\#^{\prec }(a)\text{, para cada }a\in \Sigma \end{array} \)

la funcion \(\#^{\prec }\mid _{\Sigma ^{\ast }}\) es \(\Sigma \)-p.r.. O sea que el predicado \(P=\lambda x\alpha \left[ \#^{\prec }\mid _{\Sigma ^{\ast }}(\alpha )=x\right] \) es \(\Sigma \)-p.r.. Sea \(< \) un orden total estricto sobre \(\Sigma \). Note que \(\ast ^{\prec }\mid _{\#^{\prec }(\Sigma ^{\ast })}=M^{< }(P)\), lo cual ya que
\(\displaystyle \left\vert \ast ^{\prec }\mid _{\#^{\prec }(\Sigma ^{\ast })}(x)\right\vert \leq x \)

nos dice que \(\ast ^{\prec }\mid _{\#^{\prec }(\Sigma ^{\ast })}\) es \(\Sigma \)-p.r. (Lema 47). \(\Box\)


\textbf{\underline{Lemma 50:}} Supongamos \(\Gamma \neq \varnothing \) y sea \(< \) un orden total estricto sobre \( \Gamma \). Dada \(h\) una funcion \(\Gamma \)-mixta, son equivalentes
(1) \(h\) es \(\Gamma \)-recursiva (resp. \(\Gamma \)-p.r.)
(2) \(h^{\#^{< }}\) es \(\varnothing \)-recursiva (resp. \(\varnothing \)-p.r.)

\textbf{\underline{Proof:}} (2)\(\Rightarrow \)(1). Supongamos \(h:D_{h}\subseteq \omega ^{n}\times \Gamma ^{\ast m}\rightarrow \Gamma ^{\ast }\). Ya que \(h^{\#^{< }}\) es \(\Gamma \) -recursiva (resp. \(\Gamma \)-p.r.) y

\(\displaystyle h=\ast ^{< }\circ h^{\#^{< }}\circ \left( p_{1}^{n,m},...,p_{n}^{n,m},\#^{< }\circ p_{n+1}^{n,m},...,\#^{< }\circ p_{n+m}^{n,m}\right) \text{,} \)

tenemos que \(h\) es \(\Gamma \)-recursiva (resp. \(\Gamma \)-p.r.).
(1)\(\Rightarrow \)(2). Probaremos por induccion en \(k\) que

(*) Si \(h\in \mathrm{R}_{k}^{\Gamma }\) (resp. \(h\in \mathrm{PR} _{k}^{\Gamma })\), entonces \(h^{\#^{< }}\) es \(\varnothing \)-recursiva (resp. \( \varnothing \)-p.r.).
El caso \(k=0\) es facil y dejado al lector. Supongamos (*) vale para un \(k\) fijo. Veremos que vale para \(k+1\). Sea \(h\in \mathrm{R} _{k+1}^{\Gamma }\) (resp. \(h\in \mathrm{PR}_{k+1}^{\Gamma }\)). Hay varios casos

Caso 1. Supongamos \(h=f\circ (f_{1},...,f_{n})\), con \(f,f_{1},...,f_{n}\in \mathrm{R}_{k}^{\Gamma }\) (resp. \(f,f_{1},...,f_{n}\in \mathrm{PR} _{k}^{\Gamma }\)). Por hipotesis inductiva tenemos que \(f^{\#^{< }},f_{1}^{ \#^{< }},...,f_{n}^{\#^{< }}\) son \(\varnothing \)-recursivas (resp. \(\varnothing \) -p.r.). Ya que \(h^{\#^{< }}=f^{\#^{< }}\circ \left( f_{1}^{\#^{< }},...,f_{n}^{\#^{< }}\right) \), tenemos que \(h^{\#^{< }}\) es \( \varnothing \)-recursiva (resp. \(\varnothing \)-p.r.).

Caso 2. Supongamos \(h=M(P)\), con \(P:\omega \times \omega ^{n}\times \Gamma ^{\ast m}\rightarrow \omega \), un predicado en \(\mathrm{R}_{k}^{\Gamma }\). Ya que \(h^{\#^{< }}=M(P^{\#^{< }})\), tenemos que \(h^{\#^{< }}\) es \(\varnothing \) -recursiva.

Caso 3. Supongamos \(h=R(f,\mathcal{G})\), con

\(\displaystyle \begin{array}{rcl} f & :& \omega ^{n}\times \Gamma ^{\ast m}\rightarrow \Gamma ^{\ast } \\ \mathcal{G}_{a} & :& \omega ^{n}\times \Gamma ^{\ast m}\times \Gamma ^{\ast }\times \Gamma ^{\ast }\rightarrow \Gamma ^{\ast }\text{, }a\in \Gamma \end{array} \)

funciones en \(\mathrm{R}_{k}^{\Gamma }\) (resp. \(\mathrm{PR}_{k}^{\Gamma }\)). Sea \(\Gamma =\{a_{1},...,a_{r}\}\), con \(a_{1}< a_{2}< ...< a_{r}\). Por hipotesis inductiva tenemos que \(f^{\#^{< }}\) y cada \(\mathcal{G} _{a}^{\#^{< }} \) son \(\varnothing \)-recursivas (resp. \(\varnothing \)-p.r.). Sea
\(\displaystyle \begin{array}{lll} i_{0}:\omega & \rightarrow & \omega \\ \;\;\;\;\;x & \rightarrow & \left\{ \begin{array}{lll} r & & \text{si }r\text{ divide }x \\ R(x,r) & & \text{caso contrario} \end{array} \right. \end{array} \)

y sea
\(\displaystyle B=\lambda x\left[ Q(x\dot{-}i_{0}(x),r)\right] \)

(\(R\) y \(Q\) son definidas en el Lema 44). Note que \(i_{0}\) y \(B\) son \(\varnothing \)-p.r. y que
\(\displaystyle \ast ^{< }(x)=\ast ^{< }(B(x))a_{i_{0}(x)}\text{, para }x\geq 1 \)

(ver Lema 6). Tambien tenemos
\(\displaystyle \begin{array}{rcl} h^{\#^{< }}(\vec{x},\vec{y},t+1) & =& \#^{< }(h(\vec{x},\ast ^{< }(\vec{y}),\ast ^{< }(t+1))) \\ & =& \#^{< }(h(\vec{x},\ast ^{< }(\vec{y}),\ast ^{< }(B(t+1))a_{i_{0}(t+1)})) \\ & =& \#^{< }\left( \mathcal{G}_{a_{i_{0}(t+1)}}(\vec{x},\ast ^{< }(\vec{y}),\ast ^{< }(B(t+1)),h(\vec{x},\ast ^{< }(\vec{y}),\ast ^{< }(B(t+1)))\right) \\ & =& \#^{< }\left( \mathcal{G}_{a_{i_{0}(t+1)}}(\vec{x},\ast ^{< }(\vec{y}),\ast ^{< }(B(t+1)),\ast ^{< }(h^{\#^{< }}(\vec{x},\vec{y},B(t+1))))\right) \\ & =& \mathcal{G}_{a_{i_{0}(t+1)}}^{\#^{< }}(\vec{x},\vec{y},B(t+1),h^{\#^{< }}( \vec{x},\vec{y},B(t+1))) \end{array} \)

y ya que \(B(t+1)< t+1\), tenemos que
(**) \(h^{\#^{< }}(\vec{x},\vec{y},t+1)=\mathcal{G}_{a_{i_{0}(t+1)}}^{ \#^{< }}(\vec{x},\vec{y},B(t+1),\left\langle h^{\#^{< }}(\vec{x},\vec{y} ,0),...,h^{\#^{< }}(\vec{x},\vec{y},t)\right\rangle )\)
A continuacion definamos

\(\displaystyle H=\lambda t\vec{x}\vec{y}\left[ \left\langle h^{\#^{< }}(\vec{x},\vec{y} ,0),...,h^{\#^{< }}(\vec{x},\vec{y},t)\right\rangle \right] \)

Por (**) tenemos que
\(\displaystyle \begin{array}{rcl} H(0,\vec{x},\vec{y}) & =& \left\langle h^{\#^{< }}(\vec{x},\vec{y} ,0)\right\rangle =\left\langle f^{\#^{< }}(\vec{x},\vec{y})\right\rangle =2^{f^{\#^{< }}(\vec{x},\vec{y})} \\ H(t+1,\vec{x},\vec{y}) & =& \left( (H(t,\vec{x},\vec{y})+1).pr(t+2)^{\mathcal{G }_{a_{i_{0}(t+1)}}^{\#^{< }}(\vec{x},\vec{y},B(t+1),(H(t,\vec{x},\vec{y} ))_{B(t+1)})}\right) \end{array} \)

O sea que si definimos \(g:\omega \times \omega \times \omega ^{n}\times \omega ^{m}\rightarrow \omega \) por
\(\displaystyle g(z,t,\vec{x},\vec{y})=\left\{ \begin{array}{clc} \left( (z+1).pr(t+2)^{\mathcal{G}_{a_{1}}^{\#^{< }}(\vec{x},\vec{y} ,B(t+1),(z)_{B(t+1)})}\right) & \text{si} & i_{0}(t+1)=1 \\ \vdots & & \vdots \\ \left( (z+1).pr(t+2)^{\mathcal{G}_{a_{r}}^{\#^{< }}(\vec{x},\vec{y} ,B(t+1),(z)_{B(t+1)})}\right) & \text{si} & i_{0}(t+1)=r \end{array} \right. \)

tenemos que \(H=R(\lambda x\left[ 2^{x}\right] \circ f^{\#^{< }},g)\). Note que \(g\) es \(\varnothing \)-recursiva (resp. \(\varnothing \)-p.r.), ya que
\(\displaystyle g=f_{1}(z,t,\vec{x},\vec{y})P_{1}(z,t,\vec{x},\vec{y})+...+f_{r}(z,t,\vec{x}, \vec{y})P_{r}(z,t,\vec{x},\vec{y})\text{,} \)

con
\(\displaystyle \begin{array}{rcl} f_{i} & =& \lambda zt\vec{x}\vec{y}\left[ \left( (z+1).pr(t+2)^{\mathcal{G} _{a_{i}}^{\#^{< }}(\vec{x},\vec{y},B(t+1),(z)_{B(t+1)})}\right) \right] \\ P_{i} & =& \lambda zt\vec{x}\vec{y}\left[ i_{0}(t+1)=i\right] \end{array} \)

y estas funciones son totales y \(\varnothing \)-recursivas (resp. \(\varnothing \) -p.r.). O sea que \(H\) es \(\varnothing \)-recursiva (resp. \(\varnothing \)-p.r.) y por lo tanto lo es
\(\displaystyle h^{\#^{< }}=\lambda \vec{x}\vec{y}t\left[ (H(t,\vec{x},\vec{y}))_{t+1}\right] \)

Los otros casos en los cuales \(h\) es obtenida por recursion primitiva son similares. \(\Box\)


\textbf{\underline{Theorem 51:}} Sean \(\Sigma \) y \(\Gamma \) alfabetos cualesquiera.
(a) Supongamos una funcion \(f\) es \(\Sigma \)-mixta y \(\Gamma \)-mixta, entonces \(f\) es \(\Sigma \)-recursiva (resp. \(\Sigma \)-p.r.) sii \(f\) es \( \Gamma \)-recursiva (resp. \(\Gamma \)-p.r.).
(b) Supongamos un conjunto \(S\) es \(\Sigma \)-mixto y \(\Gamma \)-mixto, entonces \(S\) es \(\Sigma \)-p.r. sii \(S\) es \(\Gamma \)-p.r..


\textbf{\underline{Proof:}} (a) Ya que \(f\) es \((\Sigma \cap \Gamma )\)-mixta, podemos suponer sin perdida de generalidad que \(\Sigma \subseteq \Gamma \). Primero haremos el caso en que \(\Sigma =\varnothing \) y \(\Gamma \neq \varnothing \). Sea \(< \) un orden total estricto sobre \(\Gamma \). Ya que \(f\) es \(\varnothing \)-mixta, tenemos \( f=f^{\#^{< }}\) y por lo tanto podemos aplicar el lema anterior.

Supongamos ahora que \(\Sigma \neq \varnothing \). O sea que \(f:D_{f}\subseteq \omega ^{n}\times \Sigma ^{\ast m}\rightarrow O\), con \(O\in \{\omega ,\Sigma ^{\ast }\}.\) Haremos el caso \(O=\Sigma ^{\ast }.\) Supongamos \(f\) es \(\Sigma \) -recursiva (resp. \(\Sigma \)-p.r.). Sea \(\prec \) un orden total estricto sobre \(\Gamma .\) Ya que las funciones \(\#^{\prec }\mid _{\Sigma ^{\ast }}\) y \(\ast ^{\prec }\mid _{\#^{\prec }(\Sigma ^{\ast })}\) son \(\Sigma \)-p.r. (Lema 49) y

\(\displaystyle \begin{array}{rcl} f^{\#^{\prec }} & =& \#^{\prec }\circ f\circ \left( p_{1}^{n+m,0},...,p_{n}^{n+m,0},\ast ^{\prec }\circ p_{n+1}^{n+m,0},...,\ast ^{\prec }\circ p_{n+m}^{n+m,0}\right) \\ & =& \#^{\prec }\mid _{\Sigma ^{\ast }}\circ f\circ \left( p_{1}^{n+m,0},...,p_{n}^{n+m,0},\ast ^{\prec }\mid _{\#^{\prec }(\Sigma ^{\ast })}\circ p_{n+1}^{n+m,0},...,\ast ^{\prec }\mid _{\#^{\prec }(\Sigma ^{\ast })}\circ p_{n+m}^{n+m,0}\right) \end{array} \)

tenemos que \(f^{\#^{\prec }}\) es \(\Sigma \)-recursiva (resp. \(\Sigma \)-p.r.). O sea que por el caso ya probado de (a), \(f^{\#^{\prec }}\) es \(\varnothing \) -recursiva (resp. \(\varnothing \)-p.r.) lo cual por el lema anterior nos dice que \(f\) es \(\Gamma \)-recursiva (resp. \(\Gamma \)-p.r.).
Supongamos ahora que \(f\) es \(\Gamma \)-recursiva (resp. \(\Gamma \)-p.r.). Sea \( < \) un orden total estricto sobre \(\Sigma .\) Ya que \(\#^{< }\) y \(\ast ^{< }\) son \(\Gamma \)-p.r. (Lema 49), la funcion

\(\displaystyle f^{\#^{< }}=\#^{< }\circ f\circ \left( p_{1}^{n+m,0},...,p_{n}^{n+m,0},\ast ^{< }\circ p_{n+1}^{n+m,0},...,\ast ^{< }\circ p_{n+m}^{n+m,0}\right) \)

es \(\Gamma \)-recursiva (resp. \(\Gamma \)-p.r.). Por el caso ya probado de (a), \(f^{\#^{< }}\) es \(\varnothing \)-recursiva (resp. \(\varnothing \)-p.r.), lo cual por el lema anterior nos dice que \(f\) es \(\Sigma \)-recursiva (resp. \( \Sigma \)-p.r.).
(b) es dejado al lector (use (a)). \(\Box\)

\section{Sintaxis de $\mathcal{S}^{\Sigma}$}

\textbf{\underline{Lemma 52:}} Se tiene que:
(a) Si \(I_{1}...I_{n}=J_{1}...J_{m}\), con \( I_{1},...,I_{n},J_{1},...,J_{m}\in \mathrm{Ins}^{\Sigma }\), entonces \(n=m\) y \(I_{j}=J_{j}\) para cada \(j\geq 1\).
(b) Si \(\mathcal{P}\in \mathrm{Pro}^{\Sigma }\), entonces existe una unica sucesion de instrucciones \(I_{1},...,I_{n}\) tal que \(\mathcal{P} =I_{1}...I_{n}\)


\textbf{\underline{Proof:}} (a) Supongamos \(I_{n}\) es un tramo final propio de \(J_{m}.\) Notar que entonces \(n >1\). Es facil ver que entonces ya sea \(J_{m}=\mathrm{L}\bar{u} I_{n}\) para algun \(u\in \mathbf{N}\), o \(I_{n}\) es de la forma \(\mathrm{GOTO} \;\mathrm{L}\bar{n}\) y \(J_{m}\) es de la forma \(w\mathrm{IF}\;\mathrm{P}\bar{k }\;\mathrm{BEGINS}\;a\;\mathrm{GOTO}\;\mathrm{L}\bar{n}\) donde \(w\in \{ \mathrm{L}\bar{n}:n\in \mathbf{N}\}\cup \{\varepsilon \}\). El segundo caso no puede darse porque entonces el anteultimo simbolo de \(I_{n-1}\) deberia ser \(\mathrm{S}\) lo cual no sucede para ninguna instruccion. O sea que

\(\displaystyle I_{1}...I_{n}=J_{1}...J_{m-1}\mathrm{L}\bar{u}I_{n} \)

lo cual dice que
(*) \(I_{1}...I_{n-1}=J_{1}...J_{m-1}\mathrm{L}\bar{u}.\)
Es decir que \(\mathrm{L}\bar{u}\) es tramo final de \(I_{n-1}\) y por lo tanto \(\mathrm{GOTO}\;\mathrm{L}\bar{u}\) es tramo final de \(I_{n-1}.\) Por (*), \(\mathrm{GOTO}\) es tramo final de \(J_{1}...J_{m-1}\), lo cual es impossible. Hemos llegado a una contradiccion lo cual nos dice que \(I_{n}\) no es un tramo final propio de \(J_{m}.\) Por simetria tenemos que \( I_{n}=J_{m} \), lo cual usando un razonamiento inductivo nos dice que \(n=m\) y \(I_{j}=J_{j} \) para cada \(j\geq 1\).

(b) Es consecuencia directa de (a). \(\Box\)

\section{Funciones $\Sigma$-computables}


\textbf{\underline{Theorem 53:}} Si \(f\) es \(\Sigma \)-computable, entonces \(f\) es \(\Sigma \)-efectivamente computable.


\textbf{\underline{Proof:}} Supongamos por ejemplo que \(f:S\subseteq \omega ^{n}\times \Sigma ^{\ast m}\rightarrow \omega \) es computada por \(\mathcal{P}\in \mathrm{Pro}^{\Sigma }\). Es claro que el procedimiento que consiste en realizar las sucesivas instrucciones de \(\mathcal{P}\) (partiendo de \(((x_{1},...,x_{n},0,0,...),( \alpha _{1},...,\alpha _{m},\varepsilon ,\varepsilon ,...))\)) y eventualmente terminar en caso de que nos toque realizar la instruccion \(n( \mathcal{P})+1\), y dar como salida el contenido de la variable \(\mathrm{N}1\) , es un procedimiento efectivo que computa a \(f\). \(\Box\)


\textbf{\underline{Proposition 54:}}

(a) Sea \(f:S\subseteq \omega ^{n}\times \Sigma ^{\ast m}\rightarrow \omega \) una funcion \(\Sigma \)-computable. Entonces hay un macro
\(\displaystyle \left[ \mathrm{V}\overline{n+1}\leftarrow f(\mathrm{V}1,...,\mathrm{V}\bar{n} ,\mathrm{W}1,...,\mathrm{W}\bar{m})\right] \)
(b) Sea \(f:S\subseteq \omega ^{n}\times \Sigma ^{\ast m}\rightarrow \Sigma ^{\ast }\) una funcion \(\Sigma \)-computable. Entonces hay un macro
\(\displaystyle \left[ \mathrm{W}\overline{m+1}\leftarrow f(\mathrm{V}1,...,\mathrm{V}\bar{n} ,\mathrm{W}1,...,\mathrm{W}\bar{m})\right] \)

\textbf{\underline{Proof:}} (b) Sea \(\mathcal{P}\) un programa que compute a \(f\). Tomemos un \(k\) tal que \( k\geq n,m\) y tal que todas las variables y labels de \(\mathcal{P}\) estan en el conjunto

\(\displaystyle \{\mathrm{N}1,...,\mathrm{N}\bar{k},\mathrm{P}1,...,\mathrm{P}\bar{k}, \mathrm{L}1,...,\mathrm{L}\bar{k}\}\text{.} \)

Sea \(\mathcal{P}^{\prime }\) la palabra que resulta de reemplazar en \( \mathcal{P}\):
- la variable \(\mathrm{N}\overline{j}\) por \(\mathrm{V}\overline{n+j}\) , para cada \(j=1,...,k\)
- la variable \(\mathrm{P}\overline{j}\) por \(\mathrm{W}\overline{m+j}\) , para cada \(j=1,...,k\)
- el label \(\mathrm{L}\overline{j}\) por \(\mathrm{A}\overline{j}\), para cada \(j=1,...,k\)
Notese que

\(\displaystyle \begin{array}{l} \mathrm{V}\overline{n+1}\leftarrow \mathrm{V}1 \\ \ \ \ \ \ \ \ \ \ \vdots \\ \mathrm{V}\overline{n+n}\leftarrow \mathrm{V}\overline{n} \\ \mathrm{V}\overline{n+n+1}\leftarrow 0 \\ \ \ \ \ \ \ \ \ \ \vdots \\ \mathrm{V}\overline{n+k}\leftarrow 0 \\ \mathrm{W}\overline{m+1}\leftarrow \mathrm{W}1 \\ \ \ \ \ \ \ \ \ \ \vdots \\ \mathrm{W}\overline{m+m}\leftarrow \mathrm{W}\overline{m} \\ \mathrm{W}\overline{m+m+1}\leftarrow \varepsilon \\ \ \ \ \ \ \ \ \ \ \vdots \\ \mathrm{W}\overline{m+k}\leftarrow \varepsilon \\ \mathcal{P}^{\prime } \end{array} \)

es el macro buscado, el cual tendra sus variables auxiliares y labels en la lista
\(\displaystyle \mathrm{V}\overline{n+1},...,\mathrm{V}\overline{n+k},\mathrm{W}\overline{m+2 },...,\mathrm{V}\overline{m+k},\mathrm{A}1,...,\mathrm{A}\overline{k}. \)

\(\Box\)




\textbf{\underline{Proposition 55:}} Sea \(P:S\subseteq \omega ^{n}\times \Sigma ^{\ast m}\rightarrow \omega \) un predicado \(\Sigma \)-computable. Entonces hay un macro
\(\displaystyle \left[ \mathrm{IF}\;P(\mathrm{V}1,...,\mathrm{V}\bar{n},\mathrm{W}1,..., \mathrm{W}\bar{m})\;\mathrm{GOTO}\;\mathrm{A}1\right] \)



\textbf{\underline{Theorem 56:}} Si \(h\) es \(\Sigma \)-recursiva, entonces \(h\) es \(\Sigma \) -computable.

\textbf{\underline{Proof:}} Probaremos por induccion en \(k\) que

(*) Si \(h\in \mathrm{R}_{k}^{\Sigma }\), entonces \(h\) es \(\Sigma \) -computable.
El caso \(k=0\) es dejado al lector. Supongamos (*) vale para \(k\), veremos que vale para \(k+1\). Sea \(h\in \mathrm{R}_{k+1}^{\Sigma }-\mathrm{R} _{k}^{\Sigma }.\) Hay varios casos

Caso 1. Supongamos \(h=M(P)\), con \(P:\omega \times \omega ^{n}\times \Sigma ^{\ast m}\rightarrow \omega \), un predicado perteneciente a \(\mathrm{R} _{k}^{\Sigma }\). Por hipotesis inductiva, \(P\) es \(\Sigma \)- computable y por lo tanto tenemos un macro

\(\displaystyle \left[ \mathrm{IF}\;P(\mathrm{V}1,...,\mathrm{V}\overline{n+1},\mathrm{W} 1,...,\mathrm{W}\bar{m})\;\mathrm{GOTO}\;\mathrm{A}1\right] \)

lo cual nos permite realizar el siguiente programa
\(\displaystyle \begin{array}{ll} \mathrm{L}2 & \mathrm{IF}\;P(\mathrm{N}\overline{n+1},\mathrm{N}1,..., \mathrm{N}\bar{n},\mathrm{P}1,...,\mathrm{P}\bar{m})\text{\ }\mathrm{GOTO}\; \mathrm{L}1 \\ & \mathrm{N}\overline{n+1}\leftarrow \mathrm{N}\overline{n+1}+1 \\ & \mathrm{GOTO}\;\mathrm{L}2 \\ \mathrm{L}1 & \mathrm{N}1\leftarrow \mathrm{N}\overline{n+1} \end{array} \)

Es facil chequear que este programa computa \(h.\)
Caso 2. Supongamos \(h=R(f,\mathcal{G})\), con

\(\displaystyle \begin{array}{rcl} f & :& S_{1}\times ...\times S_{n}\times L_{1}\times ...\times L_{m}\rightarrow \Sigma ^{\ast } \\ \mathcal{G}_{a} & :& S_{1}\times ...\times S_{n}\times L_{1}\times ...\times L_{m}\times \Sigma ^{\ast }\times \Sigma ^{\ast }\rightarrow \Sigma ^{\ast } \text{, }a\in \Sigma \end{array} \)

elementos de \(\mathrm{R}_{k}^{\Sigma }\). Sea \(\Sigma =\{a_{1},...,a_{r}\}.\) Por hipotesis inductiva, las funciones \(f\), \(\mathcal{G}_{a}\), \(a\in \Sigma \) , son \(\Sigma \)-computables y por lo tanto podemos hacer el siguiente programa via el uso de macros
\(\displaystyle \begin{array}{rl} & \left[ \mathrm{P}\overline{m+3}\leftarrow f(\mathrm{N}1,...,\mathrm{N}\bar{ n},\mathrm{P}1,...,\mathrm{P}\bar{m})\right] \\ \mathrm{L}\overline{r+1} & \mathrm{IF}\;\mathrm{P}\overline{m+1}\ \text{ {B}}\mathrm{EGINS\ }a_{1}\text{ }\mathrm{GOTO}\;\mathrm{L}1 \\ & \ \ \ \ \ \ \ \ \ \ \ \ \vdots \\ & \mathrm{IF}\;\mathrm{P}\overline{m+1}\ \mathrm{BEGINS\ }a_{r}\text{ } \mathrm{GOTO}\;\mathrm{L}\bar{r} \\ & \mathrm{GOTO}\;\mathrm{L}\overline{r+2} \\ \mathrm{L}1 & \mathrm{P}\overline{m+1}\leftarrow \text{ }^{\curvearrowright } \mathrm{P}\overline{m+1} \\ & \left[ \mathrm{P}\overline{m+3}\leftarrow \mathcal{G}_{a_{1}}(\mathrm{N} 1,...,\mathrm{N}\bar{n},\mathrm{P}1,...,\mathrm{P}\bar{m},\mathrm{P} \overline{m+2},\mathrm{P}\overline{m+3})\right] \\ & \mathrm{P}\overline{m+2}\leftarrow \mathrm{P}\overline{m+2}a_{1} \\ & \mathrm{GOTO}\;\mathrm{L}\overline{r+1} \\ & \ \ \ \ \ \ \ \ \ \ \ \ \vdots \\ \mathrm{L}\bar{r} & \mathrm{P}\overline{m+1}\leftarrow \text{ } ^{\curvearrowright }\mathrm{P}\overline{m+1} \\ & \mathrm{P}\overline{m+3}\leftarrow \mathcal{G}_{a_{r}}(\mathrm{N}1,..., \mathrm{N}\bar{n},\mathrm{P}1,...,\mathrm{P}\bar{m},\mathrm{P}\overline{m+2}, \mathrm{P}\overline{m+3}) \\ & \mathrm{P}\overline{m+2}\leftarrow \mathrm{P}\overline{m+2}a_{r} \\ & \mathrm{GOTO}\;\mathrm{L}\overline{r+1} \\ \mathrm{L}\overline{r+2} & \mathrm{P}1\leftarrow \mathrm{P}\overline{m+3} \end{array} \)

Es facil chequear que este programa computa \(h.\)
El resto de los casos son dejados al lector. \(\Box\)

\section{Análisis de la recursividad de $\mathcal{S}^{\Sigma}$}


\textbf{\underline{Lemma 57:}} Sea \(\Sigma \) un alfabeto cualquiera. Las funciones \(S\) y \(\overline{\ \;}\) son \((\Sigma \cup \Sigma _{p})\)-p.r..

\textbf{\underline{Proof:}} Use el Teorema 51. \(\Box\)

\textbf{\underline{Lemma 58:}} Para cada \(n,x\in \omega \), tenemos que \( \left\vert \bar{n}\right\vert \leq x\) si y solo si \(n\leq 10^{x}-1\)

\textbf{\underline{Lemma 59:}} \(\mathrm{Ins}^{\Sigma }\) es un conjunto \((\Sigma \cup \Sigma _{p})\)-p.r..

\textbf{\underline{Proof:}} Para simplificar la Proof asumiremos que \(\Sigma =\{@,\& \}\). Ya que \( \mathrm{Ins}^{\Sigma }\) es union de los siguientes conjuntos

\(\displaystyle \begin{array}{rcl} L_{1} & =& \left\{ \mathrm{N}\bar{k}\leftarrow \mathrm{N}\bar{k}+1:k\in \mathbf{N}\right\} \\ L_{2} & =& \left\{ \mathrm{N}\bar{k}\leftarrow \mathrm{N}\bar{k}\dot{-}1:k\in \mathbf{N}\right\} \\ L_{3} & =& \left\{ \mathrm{N}\bar{k}\leftarrow \mathrm{N}\bar{n}:k,n\in \mathbf{N}\right\} \\ L_{4} & =& \left\{ \mathrm{N}\bar{k}\leftarrow 0:k\in \mathbf{N}\right\} \\ L_{5} & =& \left\{ \mathrm{IF}\;\mathrm{N}\bar{k}\neq 0\;\mathrm{GOTO}\; \mathrm{L}\bar{m}:k,m\in \mathbf{N}\right\} \\ L_{6} & =& \left\{ \mathrm{P}\bar{k}\leftarrow \mathrm{P}\bar{k}.@:k\in \mathbf{N}\right\} \\ L_{7} & =& \left\{ \mathrm{P}\bar{k}\leftarrow \mathrm{P}\bar{k}.\& :k\in \mathbf{N}\right\} \\ L_{8} & =& \left\{ \mathrm{P}\bar{k}\leftarrow \text{ }^{\curvearrowright } \mathrm{P}\bar{k}:k\in \mathbf{N}\right\} \\ L_{9} & =& \left\{ \mathrm{P}\bar{k}\leftarrow \mathrm{P}\bar{n}:k,n\in \mathbf{N}\right\} \\ L_{9} & =& \left\{ \mathrm{P}\bar{k}\leftarrow \varepsilon :k\in \mathbf{N} \right\} \\ L_{10} & =& \left\{ \mathrm{IF}\;\mathrm{P}\bar{k}\;\mathrm{BEGINS}\;@\; \mathrm{GOTO}\;\mathrm{L}\bar{m}:k,m\in \mathbf{N}\right\} \\ L_{11} & =& \left\{ \mathrm{IF}\;\mathrm{P}\bar{k}\;\mathrm{BEGINS}\;\& \; \mathrm{GOTO}\;\mathrm{L}\bar{m}:k,m\in \mathbf{N}\right\} \\ L_{12} & =& \left\{ \mathrm{GOTO}\;\mathrm{L}\bar{m}:m\in \mathbf{N}\right\} \\ L_{13} & =& \left\{ \mathrm{SKIP}\right\} \\ L_{14} & =& \left\{ \mathrm{L}\bar{k}\alpha :k\in \mathbf{N\;}\text{y }\alpha \in L_{1}\cup ...\cup L_{13}\right\} \end{array} \)

solo debemos probar que \(L_{1},...,L_{14}\) son \((\Sigma \cup \Sigma _{p})\) -p.r.. Veremos primero por ejemplo que
\(\displaystyle L_{10}=\left\{ \mathrm{IFP}\bar{k}\mathrm{BEGINS}@\mathrm{GOTOL}\bar{m} :k,m\in \mathbf{N}\right\} \)

es \((\Sigma \cup \Sigma _{p})\)-p.r.. Primero notese que \(\alpha \in L_{10}\) si y solo si existen \(k,m\in \mathbf{N}\) tales que
\(\displaystyle \alpha =\mathrm{IFP}\bar{k}\mathrm{BEGINS}@\mathrm{GOTOL}\bar{m} \)

Mas formalmente tenemos que \(\alpha \in L_{10}\) si y solo si
\(\displaystyle (\exists k\in \mathbf{N})(\exists m\in \mathbf{N})\;\alpha =\mathrm{IFP}\bar{ k}\mathrm{BEGINS}@\mathrm{GOTOL}\bar{m} \)

Ya que cuando existen tales \(k,m\) tenemos que \(\bar{k}\) y \(\bar{m}\) son subpalabras de \(\alpha \), el lema anterior nos dice que \(\alpha \in L_{10}\) si y solo si
\(\displaystyle (\exists k\in \mathbf{N})_{k\leq 10^{\left\vert \alpha \right\vert }}(\exists m\in \mathbf{N})_{m\leq 10^{\left\vert \alpha \right\vert }}\;\alpha =\mathrm{IFP}\bar{k}\mathrm{BEGINS}@\mathrm{GOTOL}\bar{m} \)

Sea
\(\displaystyle P=\lambda mk\alpha \left[ \alpha =\mathrm{IFP}\bar{k}\mathrm{BEGINS}@\mathrm{ GOTOL}\bar{m}\right] \)

Ya que \(D_{\lambda k\left[ \bar{k}\right] }=\omega \), tenemos que \( D_{P}=\omega \times (\Sigma \cup \Sigma _{p})^{\ast }\times (\Sigma \cup \Sigma _{p})^{\ast }\). Notese que
\(\displaystyle P=\lambda \alpha \beta \left[ \alpha =\beta \right] \circ \left( p_{3}^{2,1},f\right) \)

donde
\(\displaystyle f=\lambda \alpha _{1}\alpha _{2}\alpha _{3}\alpha _{4}\left[ \alpha _{1}\alpha _{2}\alpha _{3}\alpha _{4}\right] \circ \left( C_{\mathrm{IFP} }^{2,1},\lambda k\left[ \bar{k}\right] \circ p_{2}^{2,1},C_{\mathrm{BEGINS}@ \mathrm{GOTOL}}^{2,1},\lambda k\left[ \bar{k}\right] \circ p_{1}^{2,1}\right) \)

lo cual nos dice que \(P\) es \((\Sigma \cup \Sigma _{p})\)-p.r..
Notese que

\(\displaystyle \chi _{L_{10}}=\lambda \alpha \left[ (\exists k\in \mathbf{N})_{k\leq 10^{\left\vert \alpha \right\vert }}(\exists m\in \mathbf{N})_{m\leq 10^{\left\vert \alpha \right\vert }}\;P(m,k,\alpha )\right] \)

Esto nos dice que podemos usar dos veces el Lema 39 para ver que \(\chi _{L_{10}}\) es \((\Sigma \cup \Sigma _{p})\)-p.r.. Veamos como. Sea
\(\displaystyle Q=\lambda k\alpha \left[ (\exists m\in \mathbf{N})_{m\leq 10^{\left\vert \alpha \right\vert }}\;P(m,k,\alpha )\right] \)

Por el Lema 39 tenemos que
\(\displaystyle \lambda xk\alpha \left[ (\exists m\in \mathbf{N})_{m\leq x}\;P(m,k,\alpha ) \right] \)

es \((\Sigma \cup \Sigma _{p})\)-p.r. lo cual nos dice que
\(\displaystyle Q=\lambda xk\alpha \left[ (\exists m\in \mathbf{N})_{m\leq x}\;P(m,k,\alpha ) \right] \circ (\lambda \alpha \left[ 10^{\left\vert \alpha \right\vert } \right] \circ p_{2}^{1,1},p_{1}^{1,1},p_{2}^{1,1}) \)

lo es. Ya que
\(\displaystyle \chi _{L_{10}}=\lambda \alpha \left[ (\exists k\in \mathbf{N})_{k\leq 10^{\left\vert \alpha \right\vert }}\;Q(k,\alpha )\right] \)

podemos en forma similar aplicar el Lema 39 y obtener finalmente que \(\chi _{L_{10}}\) es \((\Sigma \cup \Sigma _{p})\)-p.r..
En forma similar podemos probar que \(L_{1},...,L_{13}\) son \((\Sigma \cup \Sigma _{p})\)-p.r.. Esto nos dice que \(L_{1}\cup ...\cup L_{13}\) es \((\Sigma \cup \Sigma _{p})\)-p.r.. Notese que \(L_{1}\cup ...\cup L_{13}\) es el conjunto de las instrucciones basicas de \(\mathcal{S}^{\Sigma }\). Llamemos \( \mathrm{InsBas}^{\Sigma }\) a dicho conjunto. Para ver que \(L_{14}\) es \( (\Sigma \cup \Sigma _{p})\)-p.r. notemos que

\(\displaystyle \chi _{L_{14}}=\lambda \alpha \left[ (\exists k\in \mathbf{N})_{k\leq 10^{\left\vert \alpha \right\vert }}(\exists \beta \in \mathrm{InsBas} ^{\Sigma })_{\left\vert \beta \right\vert \leq \left\vert \alpha \right\vert }\;\alpha =\mathrm{L}\bar{k}\beta \right] \)

lo cual nos dice que aplicando dos veces el Lema 39 obtenemos que \(\chi _{L_{14}}\) es \((\Sigma \cup \Sigma _{p})\)-p.r.. Ya que \( \mathrm{Ins}^{\Sigma }=\mathrm{InsBas}^{\Sigma }\cup L_{14}\) tenemos que \( \mathrm{Ins}^{\Sigma }\) es \((\Sigma \cup \Sigma _{p})\)-p.r.. \(\Box\)


\textbf{\underline{Lemma 60:}} \(Bas\) y \(Lab\) son funciones \((\Sigma \cup \Sigma _{p})\)-p.r.

\textbf{\underline{Proof:}} Sea \(< \) un orden total estricto sobre \(\Sigma \cup \Sigma _{p}\). Sea \(L=\{ \mathrm{L}\bar{k}:k\in \mathbf{N}\}\cup \{\varepsilon \}\). Dejamos al lector probar que \(L\) es un conjunto \((\Sigma \cup \Sigma _{p})\)-p.r.. Sea

\(\displaystyle P=\lambda I\alpha \left[ \alpha \in \mathrm{Ins}^{\Sigma }\wedge I\in \mathrm{Ins}^{\Sigma }\wedge \lbrack \alpha ]_{1}\neq \mathrm{L}\wedge (\exists \beta \in L)\ I=\beta \alpha \right] \)

Note que \(D_{P}=(\Sigma \cup \Sigma _{p})^{\ast 2}\). Dejamos al lector probar que \(P\) es \((\Sigma \cup \Sigma _{p})\)-p.r.. Notese ademas que cuando \(I\in \mathrm{Ins}^{\Sigma }\) tenemos que \(P(I,\alpha )=1\) sii \(\alpha =Bas(I)\). Dejamos al lector probar que \(Bas=M^{< }\left( P\right) \) por lo que para ver que \(Bas\) es \((\Sigma \cup \Sigma _{p})\)-p.r., solo nos falta ver que la funcion \(Bas\) es acotada por alguna funcion \((\Sigma \cup \Sigma _{p})\)-p.r. y \((\Sigma \cup \Sigma _{p})\)-total. Pero esto es trivial ya que \(\left\vert Bas(I)\right\vert \leq \left\vert I\right\vert =p_{1}^{0,1}(I)\) para cada \(I\in \mathrm{Ins}^{\Sigma }\).
Finalmente note que

\(\displaystyle Lab=M^{< }\left( \lambda I\alpha \left[ \alpha Bas(I)=I\right] \right) \)

lo cual nos dice que \(Lab\) es \((\Sigma \cup \Sigma _{p})\)-p.r.. \(\Box\)
Recordemos que dado un programa \(\mathcal{P}\) habiamos definido \(I_{i}^{ \mathcal{P}}=\varepsilon \), para \(i=0\) o \(i >n(\mathcal{P}).\) O sea que la funcion \((\Sigma \cup \Sigma _{p})\)-mixta \(\lambda i\mathcal{P}\left[ I_{i}^{ \mathcal{P}}\right] \) tiene dominio igual a \(\omega \times \mathrm{Pro} ^{\Sigma }\).

\textbf{\underline{Lemma 61:}}
(a) \(\mathrm{Pro}^{\Sigma }\) es un conjunto \((\Sigma \cup \Sigma _{p}) \)-p.r.
(b) \(\lambda \mathcal{P}\left[ n(\mathcal{P})\right] \) y \(\lambda i \mathcal{P}\left[ I_{i}^{\mathcal{P}}\right] \) son funciones \((\Sigma \cup \Sigma _{p})\)-p.r..

\textbf{\underline{Proof:}} Ya que \(\mathrm{Pro}^{\Sigma }=D_{\lambda \mathcal{P}\left[ n(\mathcal{P}) \right] }\) tenemos que (b) implica (a). Para probar (b) sea \(< \) un orden total estricto sobre \(\Sigma \cup \Sigma _{p}\). Sea \(P\) el siguiente predicado

\(\lambda x\left[ Lt(x) >0\wedge (\forall t\in \mathbf{N})_{t\leq Lt(x)}\;\ast ^{< }((x)_{t})\in \mathrm{Ins}^{\Sigma }\wedge \right. \)
\(\ \ \ \ \ \ \ \ \ \ \ \ \ \ \ \ \ \ \ \ \ \ \ \ \ \ \ \ \ (\forall t\in \mathbf{N})_{t\leq Lt(x)}(\forall m\in \mathbf{N})\;\lnot (\mathrm{L} \bar{m}\ \)t-final \(\ast ^{< }((x)_{t}))\vee \)
\(\ \ \ \ \ \ \ \ \ \ \ \ \ \ \ \ \ \ \ \ \ \ \ \ \ \ \ \ \ \ \ \ \ \ \ \ \ \ \ \ \ \ \ \ \ \ \ \ \ \ \ \ \ \ \ \ \ \ \ \left. (\exists j\in \mathbf{ N})_{j\leq Lt(x)}(\exists \alpha \in (\Sigma \cup \Sigma _{p})-Num)\;\mathrm{ L}\bar{m}\alpha \ \text{t-inicial}\ast ^{< }((x)_{j})\right] \)
Notese que \(D_{P}=\mathbf{N}\) y que \(P(x)=1\) sii \(Lt(x) >0\), \(\ast ^{< }((x)_{t})\in \mathrm{Ins}^{\Sigma }\), para cada \(t=1,...,Lt(x)\) y ademas \(\subset _{t=1}^{t=Lt(x)}\ast ^{< }((x)_{t})\in \mathrm{Pro}^{\Sigma }\). Para ver que \(P\) es \((\Sigma \cup \Sigma _{p})\)-p.r. solo nos falta acotar el cuantificador \((\forall m\in \mathbf{N})\) de la expresion lambda que define a \(P\). Ya que nos interesan los valores de \(m\) para los cuales \(\bar{m}\) es posiblemente una subpalabra de alguna de las palabras \(\ast ^{< }((x)_{j})\), el Lema 58 nos dice que una cota posible es \( 10^{\max \{\left\vert \ast ^{< }((x)_{j})\right\vert :1\leq j\leq Lt(x)\}}-1\) . Dejamos al lector los detalles de la Proof de que \(P\) es \((\Sigma \cup \Sigma _{p})\)-p.r.. Sea

\(\displaystyle Q=\lambda x\alpha \left[ P(x)\wedge \alpha =\subset _{t=1}^{t=Lt(x)}\ast ^{< }((x)_{t})\right] \text{.} \)

Note que \(D_{Q}=\mathbf{N}\times (\Sigma \cup \Sigma _{p})^{\ast }\). Claramente \(Q\) es \((\Sigma \cup \Sigma _{p})\)-p.r.. Ademas note que \( D_{M(Q)}=\mathrm{Pro}^{\Sigma }\). Notese que para \(\mathcal{P}\in \mathrm{Pro }^{\Sigma }\), tenemos que \(M(Q)(\mathcal{P})\) es aquel numero tal que pensado como infinitupla (via mirar su secuencia de exponentes) codifica la secuencia de instrucciones que forman a \(\mathcal{P}\). Es decir
\(\displaystyle M(Q)(\mathcal{P})=\left\langle \#^{< }(I_{1}^{\mathcal{P}}),\#^{< }(I_{2}^{ \mathcal{P}}),...,\#^{< }(I_{n(\mathcal{P})}^{\mathcal{P}}),0,0,...\right \rangle \)

Por (b) del Lema 43, \(M(Q)\) es \((\Sigma \cup \Sigma _{p})\) -p.r. ya que para cada \(\mathcal{P}\in \mathrm{Pro}^{\Sigma }\) tenemos que
\(\displaystyle \begin{array}{rcl} M(Q)(\mathcal{P}) & =& \left\langle \#^{< }(I_{1}^{\mathcal{P}}),\#^{< }(I_{2}^{ \mathcal{P}}),...,\#^{< }(I_{n(\mathcal{P})}^{\mathcal{P}}),0,0,...\right \rangle \\ & =& \underset{i=1}{\overset{n(\mathcal{P})}{\Pi }}pr(i)^{\#^{< }(I_{1}^{ \mathcal{P}})} \\ & \leq & \underset{i=1}{\overset{\left\vert \mathcal{P}\right\vert }{\Pi }} pr(i)^{\#^{< }(\mathcal{P})} \end{array} \)

Ademas tenemos que
\(\displaystyle \begin{array}{rcl} \lambda \mathcal{P}\left[ n(\mathcal{P})\right] & =& \lambda x\left[ Lt(x) \right] \circ M(Q) \\ \lambda i\mathcal{P}\left[ I_{i}^{\mathcal{P}}\right] & =& \ast ^{< }\circ g\circ \left( p_{1}^{1,1},M(Q)\circ p_{2}^{1,1}\right) \end{array} \)

donde \(g=C_{0}^{1,1}\mid _{\{0\}\times \omega }\cup \lambda ix\left[ (x)_{i} \right] \), lo cual dice que \(\lambda \mathcal{P}\left[ n(\mathcal{P})\right] \) y \(\lambda i\mathcal{P}\left[ I_{i}^{\mathcal{P}}\right] \) son funciones \( (\Sigma \cup \Sigma _{p})\)-p.r.. \(\Box\)


\textbf{\underline{Lemma 62:}} Dado un orden total estricto \( < \) sobre \(\Sigma \cup \Sigma _{p}\), las funciones \(s\), \(S_{\#}\) y \(S_{\ast } \) son \((\Sigma \cup \Sigma _{p})\)-p.r..

\textbf{\underline{Proof:}} Necesitaremos algunas funciones \((\Sigma \cup \Sigma _{p})\)-p.r.. Dada una instruccion \(I\) en la cual al menos ocurre una variable, usaremos \(\#Var1(I)\) para denotar el numero de la primer variable que ocurre en \(I\). Por ejemplo

\(\displaystyle \#Var1\left( \mathrm{L}\bar{n}\;\mathrm{IF\;N}\bar{k}\neq 0\;\mathrm{GOTO\;L} \bar{m}\right) =k \)

Notese que \(\lambda I[\#Var1(I)]\) tiene dominio igual a \(\mathrm{Ins} ^{\Sigma }-L\), donde \(L\) es la union de los siguientes conjuntos \begin{gather*} \{\mathrm{GOTO L}\bar{m}:m\in \mathbf{N\}\cup }\{\mathrm{L}\bar{k} \mathrm{ GOTO L}\bar{m}:k,m\in \mathbf{N\}}
\left\{ \mathrm{SKIP}\right\} \mathbf{\cup }\{\mathrm{L}\bar{k} \mathrm{SKIP }:k\in \mathbf{N\}} \end{gather*} Dada una instruccion \(I\) en la cual ocurren dos variables, usaremos \( \#Var2(I)\) para denotar el numero de la segunda variable que ocurre en \(I\). Por ejemplo
\(\displaystyle \#Var2\left( \mathrm{N}\bar{k}\leftarrow \mathrm{N}\bar{m}\right) =m \)

Notese que el dominio de \(\lambda I[\#Var2(I)]\) es igual a la union de los siguientes conjuntos
\(\displaystyle \begin{array}{rcl} \{\mathrm{N}\bar{k} & \leftarrow & \mathrm{N}\bar{m}:k,m\in \mathbf{N\}\cup }\{ \mathrm{L}\bar{j}\ \mathrm{N}\bar{k}\leftarrow \mathrm{N}\bar{m}:j,k,m\in \mathbf{N\}} \\ \{\mathrm{P}\bar{k} & \leftarrow & \mathrm{P}\bar{m}:k,m\in \mathbf{N\}\cup }\{ \mathrm{L}\bar{j}\ \mathrm{P}\bar{k}\leftarrow \mathrm{P}\bar{m}:j,k,m\in \mathbf{N\}} \end{array} \)

Ademas notese que para una instruccion \(I\) tenemos que
\(\displaystyle \begin{array}{rcl} \#Var1(I) & =& \min_{k}(\mathrm{N}\bar{k}\mathrm{\leftarrow }\text{ }\mathrm{ ocu}\text{ }I\vee \mathrm{N}\bar{k}\mathrm{\neq }\text{ }\mathrm{ocu}\text{ } I\vee \mathrm{P}\bar{k}\mathrm{\leftarrow }\text{ }\mathrm{ocu}\text{ }I\vee \mathrm{P}\bar{k}\mathrm{B}\;\mathrm{ocu}\text{ }I) \\ \#Var2(I) & =& \min_{k}(\mathrm{N}\bar{k}\ \text{t-final }I\vee \mathrm{N}\bar{ k}\mathrm{+}\text{ }\mathrm{ocu}\text{ }I\vee \mathrm{N}\bar{k}\mathrm{\dot{- }}\text{ }\mathrm{ocu}\text{ }I\vee \mathrm{P}\bar{k}\ \text{t-final }I\vee \mathrm{P}\bar{k}.\text{ }\mathrm{ocu}\text{ }I) \end{array} \)

Esto nos dice que si llamamos \(P\) al predicado
\(\displaystyle \lambda k\alpha \left[ \alpha \in \mathrm{Ins}^{\Sigma }\wedge (\mathrm{N} \bar{k}\mathrm{\leftarrow }\text{ }\mathrm{ocu}\text{ }\alpha \vee \mathrm{N} \bar{k}\mathrm{\neq }\text{ }\mathrm{ocu}\text{ }\alpha \vee \mathrm{P}\bar{k }\mathrm{\leftarrow }\text{ }\mathrm{ocu}\text{ }\alpha \vee \mathrm{P}\bar{k }\mathrm{B}\;\mathrm{ocu}\text{ }\alpha )\right] \)

entonces \(\lambda I[\#Var1(I)]=M(P)\) por lo cual \(\lambda I[\#Var1(I)]\) es \( (\Sigma \cup \Sigma _{p})\)-p.r. Similarmente se puede ver que \(\lambda I[\#Var2(I)]\) es \((\Sigma \cup \Sigma _{p})\)-p.r.. Sea
\(\displaystyle \begin{array}{rll} F_{\dot{-}}:\mathbf{N}\times \mathbf{N} & \rightarrow & \omega \\ (x,j) & \rightarrow & \left\langle (x)_{1},....,(x)_{j-1},(x)_{j}\dot{-} 1,(x)_{j+1},...\right\rangle \end{array} \)

Ya que
\(\displaystyle F_{\dot{-}}(x,j)=\left\{ \begin{array}{lll} Q(x,pr(j)) & & \text{si }pr(j)\text{ divide }x \\ x & & \text{caso contrario} \end{array} \right. \)

tenemos que \(F_{\dot{-}}\) es \((\Sigma \cup \Sigma _{p})\)-p.r.. Sea
\(\displaystyle \begin{array}{rll} F_{+}:\mathbf{N}\times \mathbf{N} & \rightarrow & \omega \\ (x,j) & \rightarrow & \left\langle (x)_{1},....,(x)_{j-1},(x)_{j}+1,(x)_{j+1},...\right\rangle \end{array} \)

Ya que \(F_{+}(x,j)=x.pr(j)\) tenemos que \(F_{+}\) es \((\Sigma \cup \Sigma _{p}) \)-p.r.. Sea
\(\displaystyle \begin{array}{rll} F_{\leftarrow }:\mathbf{N}\times \mathbf{N}\times \mathbf{N} & \rightarrow & \omega \\ (x,j,k) & \rightarrow & \left\langle (x)_{1},....,(x)_{j-1},(x)_{k},(x)_{j+1},...\right\rangle \end{array} \)

Ya que \(F_{\leftarrow }(x,j,k)=Q(x,pr(j)^{(x)_{j}}).pr(j)^{(x)_{k}}\) tenemos que \(F_{\leftarrow }\) es \((\Sigma \cup \Sigma _{p})\)-p.r.. Sea
\(\displaystyle \begin{array}{rll} F_{0}:\mathbf{N}\times \mathbf{N} & \rightarrow & \omega \\ (x,j) & \rightarrow & \left\langle (x)_{1},....,(x)_{j-1},0,(x)_{j+1},...\right\rangle \end{array} \)

Es facil ver que \(F_{0}\) es \((\Sigma \cup \Sigma _{p})\)-p.r.. Para cada \( a\in \Sigma \), sea
\(\displaystyle \begin{array}{rll} F_{a}:\mathbf{N}\times \mathbf{N} & \rightarrow & \omega \\ (x,j) & \rightarrow & \left\langle (x)_{1},....,(x)_{j-1},\#^{< }(\ast ^{< }((x)_{j})a),(x)_{j+1},...\right\rangle \end{array} \)

Es facil ver que \(F_{a}\) es \((\Sigma \cup \Sigma _{p})\)-p.r.. En forma similar puede ser probado que
\(\displaystyle \begin{array}{rll} F_{\curvearrowright }:\mathbf{N}\times \mathbf{N} & \rightarrow & \omega \\ (x,j) & \rightarrow & \left\langle (x)_{1},....,(x)_{j-1},\#^{< }(^{\curvearrowright }(\ast ^{< }((x)_{j}))),(x)_{j+1},...\right\rangle \end{array} \)

es \((\Sigma \cup \Sigma _{p})\)-p.r.
Dado \((i,x,y,\mathcal{P})\in \omega \times \mathbf{N}\times \mathbf{N}\times \mathrm{Pro}^{\Sigma }\), tenemos varios casos en los cuales los valores \( s(i,x,y,\mathcal{P}),S_{\#}(i,x,y,\mathcal{P})\) y \(S_{\ast }(i,x,y,\mathcal{P })\) pueden ser obtenidos usando las funciones antes definidas:

(1) CASO \(i=0\vee i >n(\mathcal{P})\). Entonces
\(\displaystyle \begin{array}{rcl} s(i,x,y,\mathcal{P}) & =& i \\ S_{\#}(i,x,y,\mathcal{P}) & =& x \\ S_{\ast }(i,x,y,\mathcal{P}) & =& y \end{array} \)

(2) CASO \((\exists j\in \omega )\;Bas(I_{i}^{\mathcal{P}})=\mathrm{N} \bar{j}\leftarrow \mathrm{N}\bar{j}+1\). Entonces
\(\displaystyle \begin{array}{rcl} s(i,x,y,\mathcal{P}) & =& i+1 \\ S_{\#}(i,x,y,\mathcal{P}) & =& F_{+}(x,\#Var1(I_{i}^{\mathcal{P}})) \\ S_{\ast }(i,x,y,\mathcal{P}) & =& y \end{array} \)

(3) CASO \((\exists j\in \omega )\;Bas(I_{i}^{\mathcal{P}})=\mathrm{N} \bar{j}\leftarrow \mathrm{N}\bar{j}\dot{-}1\). Entonces
\(\displaystyle \begin{array}{rcl} s(i,x,y,\mathcal{P}) & =& i+1 \\ S_{\#}(i,x,y,\mathcal{P}) & =& F_{\dot{-}}(x,\#Var1(I_{i}^{\mathcal{P}})) \\ S_{\ast }(i,x,y,\mathcal{P}) & =& y \end{array} \)

(4) CASO \((\exists j,k\in \omega )\;Bas(I_{i}^{\mathcal{P}})=\mathrm{N }\bar{j}\leftarrow \mathrm{N}\bar{k}\). Entonces
\(\displaystyle \begin{array}{rcl} s(i,x,y,\mathcal{P}) & =& i+1 \\ S_{\#}(i,x,y,\mathcal{P}) & =& F_{\leftarrow }(x,\#Var1(I_{i}^{\mathcal{P} }),\#Var2(I_{i}^{\mathcal{P}})) \\ S_{\ast }(i,x,y,\mathcal{P}) & =& y \end{array} \)

(5) CASO \((\exists j,k\in \omega )\;Bas(I_{i}^{\mathcal{P}})=\mathrm{N }\bar{j}\leftarrow 0\). Entonces
\(\displaystyle \begin{array}{rcl} s(i,x,y,\mathcal{P}) & =& i+1 \\ S_{\#}(i,x,y,\mathcal{P}) & =& F_{0}(x,\#Var1(I_{i}^{\mathcal{P}})) \\ S_{\ast }(i,x,y,\mathcal{P}) & =& y \end{array} \)

(6) CASO \((\exists j,m\in \omega )\;\left( Bas(I_{i}^{\mathcal{P}})= \mathrm{IF}\;\mathrm{N}\bar{j}\neq 0\;\mathrm{GOTO}\;\mathrm{L}\bar{m}\wedge (x)_{j}=0\right) \). Entonces
\(\displaystyle \begin{array}{rcl} s(i,x,y,\mathcal{P}) & =& i+1 \\ S_{\#}(i,x,y,\mathcal{P}) & =& x \\ S_{\ast }(i,x,y,\mathcal{P}) & =& y \end{array} \)

(7) CASO \((\exists j,m\in \omega )\;\left( Bas(I_{i}^{\mathcal{P}})= \mathrm{IF}\;\mathrm{N}\bar{j}\neq 0\;\mathrm{GOTO}\;\mathrm{L}\bar{m}\wedge (x)_{j}\neq 0\right) \). Entonces
\(\displaystyle \begin{array}{rcl} s(i,x,y,\mathcal{P}) & =& \min_{l}\left( Lab(I_{l}^{\mathcal{P}})\neq \varepsilon \wedge Lab(I_{l}^{\mathcal{P}})\text{ }\mathrm{t}\text{ { -final} }I_{i}^{\mathcal{P}}\right) \\ S_{\#}(i,x,y,\mathcal{P}) & =& x \\ S_{\ast }(i,x,y,\mathcal{P}) & =& y \end{array} \)

(8) CASO \((\exists j\in \omega )\;Bas(I_{i}^{\mathcal{P}})=\mathrm{P} \bar{j}\leftarrow \mathrm{P}\bar{j}.a\). Entonces
\(\displaystyle \begin{array}{rcl} s(i,x,y,\mathcal{P}) & =& i+1 \\ S_{\#}(i,x,y,\mathcal{P}) & =& x \\ S_{\ast }(i,x,y,\mathcal{P}) & =& F_{a}(y,\#Var1(I_{i}^{\mathcal{P}})) \end{array} \)

(9) CASO \((\exists j\in \omega )\;Bas(I_{i}^{\mathcal{P}})=\mathrm{P} \bar{j}\leftarrow \) \(^{\curvearrowright }\mathrm{P}\bar{j}\). Entonces
\(\displaystyle \begin{array}{rcl} s(i,x,y,\mathcal{P}) & =& i+1 \\ S_{\#}(i,x,y,\mathcal{P}) & =& x \\ S_{\ast }(i,x,y,\mathcal{P}) & =& F_{\curvearrowright }(y,\#Var1(I_{i}^{ \mathcal{P}})) \end{array} \)

(10) CASO \((\exists j,k\in \omega )\;Bas(I_{i}^{\mathcal{P}})=\mathrm{ P}\bar{j}\leftarrow \mathrm{P}\bar{k}\). Entonces
\(\displaystyle \begin{array}{rcl} s(i,x,y,\mathcal{P}) & =& i+1 \\ S_{\#}(i,x,y,\mathcal{P}) & =& x \\ S_{\ast }(i,x,y,\mathcal{P}) & =& F_{\leftarrow }(y,\#Var1(I_{i}^{\mathcal{P} }),\#Var2(I_{i}^{\mathcal{P}})) \end{array} \)

(11) CASO \((\exists j\in \omega )\;Bas(I_{i}^{\mathcal{P}})=\mathrm{P} \bar{j}\leftarrow \varepsilon \). Entonces
\(\displaystyle \begin{array}{rcl} s(i,x,y,\mathcal{P}) & =& i+1 \\ S_{\#}(i,x,y,\mathcal{P}) & =& x \\ S_{\ast }(i,x,y,\mathcal{P}) & =& F_{0}(y,\#Var1(I_{i}^{\mathcal{P}})) \end{array} \)

(12) CASO \((\exists j,m\in \omega )(\exists a\in \Sigma )\;\left( Bas(I_{i}^{\mathcal{P}})=\mathrm{IF}\;\mathrm{P}\bar{j}\;\mathrm{BEGINS}\;a\; \mathrm{GOTO}\;\mathrm{L}\bar{m}\wedge \lbrack \ast ^{< }((y)_{j})]_{1}\neq a\right) \). Entonces
\(\displaystyle \begin{array}{rcl} s(i,x,y,\mathcal{P}) & =& i+1 \\ S_{\#}(i,x,y,\mathcal{P}) & =& x \\ S_{\ast }(i,x,y,\mathcal{P}) & =& y \end{array} \)

(13) CASO \((\exists j,m\in \omega )(\exists a\in \Sigma )\;\left( Bas(I_{i}^{\mathcal{P}})=\mathrm{IF\;P}\bar{j}\;\mathrm{BEGINS\;}a\;\mathrm{ GOTO\;L}\bar{m}\wedge \lbrack \ast ^{< }((y)_{j})]_{1}=a\right) \). Entonces
\(\displaystyle \begin{array}{rcl} s(i,x,y,\mathcal{P}) & =& \min_{l}\left( Lab(I_{l}^{\mathcal{P}})\neq \varepsilon \wedge Lab(I_{l}^{\mathcal{P}})\text{ }\mathrm{t}\text{ { -final} }I_{i}^{\mathcal{P}}\right) \\ S_{\#}(i,x,y,\mathcal{P}) & =& x \\ S_{\ast }(i,x,y,\mathcal{P}) & =& y \end{array} \)

(14) CASO \((\exists j\in \omega )\;Bas(I_{i}^{\mathcal{P}})=\mathrm{ GOTO}\) \(\mathrm{L}\bar{j}\). Entonces
\(\displaystyle \begin{array}{rcl} s(i,x,y,\mathcal{P}) & =& \min_{l}\left( Lab(I_{l}^{\mathcal{P}})\neq \varepsilon \wedge Lab(I_{l}^{\mathcal{P}})\text{ }\mathrm{t}\text{ { -final} }I_{i}^{\mathcal{P}}\right) \\ S_{\#}(i,x,y,\mathcal{P}) & =& x \\ S_{\ast }(i,x,y,\mathcal{P}) & =& y \end{array} \)

(15) CASO \(Bas(I_{i}^{\mathcal{P}})=\mathrm{SKIP}\). Entonces
\(\displaystyle \begin{array}{rcl} s(i,x,y,\mathcal{P}) & =& k+1 \\ S_{\#}(i,x,y,\mathcal{P}) & =& x \\ S_{\ast }(i,x,y,\mathcal{P}) & =& y \end{array} \)

O sea que los casos anteriores nos permiten definir conjuntos \( S_{1},...,S_{15}\), los cuales son disjuntos de a pares y cuya union da el conjunto \(\omega \times \mathbf{N}\times \mathbf{N}\times \mathrm{Pro} ^{\Sigma }\), de manera que cada una de las funciones \(s,S_{\#}\) y \(S_{\ast }\) pueden escribirse como union disjunta de funciones \((\Sigma \cup \Sigma _{p}) \)-p.r. restrinjidas respectivamente a los conjuntos \(S_{1},...,S_{15}\) . Ya que los conjuntos \(S_{1},...,S_{15}\) son \((\Sigma \cup \Sigma _{p})\) -p.r. el Lema 35 nos dice que \(s,S_{\#}\) y \(S_{\ast }\) lo son. \(\Box\)


\textbf{\underline{Lemma 63:}} Sean: ... hacer!!

\textbf{\underline{Proof:}} Hacer

\textbf{\underline{Proposition 64:}} Sean $n, m \leq 0$, las funciones $i^{n, m}, E_{\#j}^{n, m}, j = 1, 2, ...$ son
  $\Sigma \cup \Sigma_{p}$-PR.

\textbf{\underline{Proof:}} Sea $<$ un orden total estricto sobre $\Sigma \cup \Sigma_{p}$ y sean $s, S_{\#}, S_{\*}$ las funciones previamente definidas en el Lemma 62, definamos:

\(\displaystyle \begin{array}{rcl} C_{\#}^{n,m} & =& \lambda t\vec{x}\vec{\alpha}\mathcal{P}\left[ \left\langle E_{\#1}^{n,m}(t,\vec{x},\vec{\alpha},\mathcal{P}),E_{\#2}^{n,m}(t,\vec{x}, \vec{\alpha},\mathcal{P}),...\right\rangle \right] \\ C_{\ast }^{n,m} & =& \lambda t\vec{x}\vec{\alpha}\mathcal{P}\left[ \left\langle \#^{< }(E_{\ast 1}^{n,m}(t,\vec{x},\vec{\alpha},\mathcal{P} )),\#^{< }(E_{\ast 2}^{n,m}(t,\vec{x},\vec{\alpha},\mathcal{P} )),...\right\rangle \right] \end{array} \)

Notese que
\(\displaystyle \begin{array}{rcl} i^{n,m}(0,\vec{x},\vec{\alpha},\mathcal{P}) & =& 1 \\ C_{\#}^{n,m}(0,\vec{x},\vec{\alpha},\mathcal{P}) & =& \left\langle x_{1},...,x_{n}\right\rangle \\ C_{\ast }^{n,m}(0,\vec{x},\vec{\alpha},\mathcal{P}) & =& \left\langle \#^{< }(\alpha _{1}),...,\#^{< }(\alpha _{m})\right\rangle \\ i^{n,m}(t+1,\vec{x},\vec{\alpha},\mathcal{P}) & =& s(i^{n,m}(t,\vec{x},\vec{ \alpha},\mathcal{P}),C_{\#}^{n,m}(t,\vec{x},\vec{\alpha},\mathcal{P} ),C_{\ast }^{n,m}(t,\vec{x},\vec{\alpha},\mathcal{P})) \\ C_{\#}^{n,m}(t+1,\vec{x},\vec{\alpha},\mathcal{P}) & =& S_{\#}(i^{n,m}(t,\vec{x },\vec{\alpha},\mathcal{P}),C_{\#}^{n,m}(t,\vec{x},\vec{\alpha},\mathcal{P} ),C_{\ast }^{n,m}(t,\vec{x},\vec{\alpha},\mathcal{P})) \\ C_{\ast }^{n,m}(t+1,\vec{x},\vec{\alpha},\mathcal{P}) & =& S_{\ast }(i^{n,m}(t, \vec{x},\vec{\alpha},\mathcal{P}),C_{\#}^{n,m}(t,\vec{x},\vec{\alpha}, \mathcal{P}),C_{\ast }^{n,m}(t,\vec{x},\vec{\alpha},\mathcal{P})) \end{array} \)

Por el Lema 63 tenemos que \(i^{n,m}\), \( C_{\#}^{n,m}\) y \(C_{\ast }^{n,m}\) son \((\Sigma \cup \Sigma _{p})\)-p.r.. Ademas notese que
\(\displaystyle \begin{array}{rcl} E_{\#j}^{n,m} & =& \lambda t\vec{x}\vec{\alpha}\mathcal{P}\left[ (C_{\#}^{n,m}(t,\vec{x},\vec{\alpha},\mathcal{P}))_{j}\right] \\ E_{\ast j}^{n,m} & =& \lambda t\vec{x}\vec{\alpha}\mathcal{P}\left[ \ast ^{< }((C_{\ast }^{n,m}(t,\vec{x},\vec{\alpha},\mathcal{P}))_{j})\right] \end{array} \)

por lo cual las funciones \(E_{\#j}^{n,m}\), \(E_{\ast j}^{n,m}\), \(j=1,2,...\), son \((\Sigma \cup \Sigma _{p})\)-p.r. \(\Box\)


\textbf{\underline{Theorem 65:}} Las funciones \(\Phi _{\#}^{n,m}\) y \(\Phi _{\ast }^{n,m}\) son \((\Sigma \cup \Sigma _{p})\)-recursivas.

\textbf{\underline{Proof:}} Veremos que \(\Phi _{\#}^{n,m}\) es \((\Sigma \cup \Sigma _{p})\)-recursiva. Sea \(H\) el predicado \((\Sigma \cup \Sigma _{p})\)-mixto

\(\displaystyle \lambda t\vec{x}\vec{\alpha}\mathcal{P}\left[ i^{n,m}(t,x_{1},...,x_{n}, \alpha _{1},...,\alpha _{m},\mathcal{P})=n(\mathcal{P})+1\right] \text{.} \)

Note que \(D_{H}=\omega ^{n+1}\times \Sigma ^{\ast m}\times \mathrm{Pro} ^{\Sigma }\). Ya que the functiones \(i^{n,m}\) y \(\lambda \mathcal{P}\left[ n( \mathcal{P})\right] \) son \((\Sigma \cup \Sigma _{p})\)-p.r., \(H\) lo es. Notar que \(D_{M(H)}=D_{\Phi _{\#}^{n,m}}\). Ademas para \((\vec{x},\vec{\alpha}, \mathcal{P})\in D_{M(H)}\), tenemos que \(M(H)(\vec{x},\vec{\alpha},\mathcal{P} )\) es la menor cantidad de pasos necesarios para que \(\mathcal{P}\) termine partiendo del estado \(((x_{1},...,x_{n},0,0,...),(\alpha _{1},...,\alpha _{m},\varepsilon ,\varepsilon ,...))\). Ya que \(H\) es \((\Sigma \cup \Sigma _{p})\)-p.r., tenemos que \(M(H)\) es \((\Sigma \cup \Sigma _{p})\)-r.. Notese que para \((\vec{x},\vec{\alpha},\mathcal{P})\in D_{M(H)}=D_{\Phi _{\#}^{n,m}} \) tenemos que
\(\displaystyle \Phi _{\#}^{n,m}(\vec{x},\vec{\alpha},\mathcal{P})=E_{\#1}^{n,m}\left( M(H)( \vec{x},\vec{\alpha},\mathcal{P}),\vec{x},\vec{\alpha},\mathcal{P}\right) \)

lo cual con un poco mas de trabajo nos permite probar que
\(\displaystyle \Phi _{\#}^{n,m}=E_{\#1}^{n,m}\circ \left( M(H),p_{1}^{n,m+1},...,p_{n+m+1}^{n,m+1}\right) \)

Ya que la funcion \(E_{\#1}^{n,m}\) es \((\Sigma \cup \Sigma _{p})\)-r., lo es \( \Phi _{\#}^{n,m}\). \(\Box\)



\textbf{\underline{Corollary 66:}} Si \(f:D_{f}\subseteq \omega ^{n}\times \Sigma ^{\ast m}\rightarrow O\) es \( \Sigma \)-computable, entonces \(f\) es \(\Sigma \)-recursiva.

\textbf{\underline{Proof:}} Haremos el caso \(O=\Sigma ^{\ast }\). Sea \(\mathcal{P}_{0}\) un programa que compute a \(f\). Primero veremos que \(f\) es \((\Sigma \cup \Sigma _{p})\) -recursiva. Note que

\(\displaystyle f=\Phi _{\ast }^{n,m}\circ \left( p_{1}^{n,m},...,p_{n+m}^{n,m},C_{\mathcal{P }_{0}}^{n,m}\right) \)

donde cabe destacar que \(p_{1}^{n,m},...,p_{n+m}^{n,m}\) son las proyecciones respecto del alfabeto \(\Sigma \cup \Sigma _{p}\), es decir que tienen dominio \(\omega ^{n}\times (\Sigma \cup \Sigma _{p})^{\ast m}\). Ya que \(\Phi _{\ast }^{n,m}\) es \((\Sigma \cup \Sigma _{p})\)-recursiva tenemos que \(f\) lo es. O sea que el Teorema 51 nos dice que \(f\) es \(\Sigma \) -recursiva. \(\Box\)


\textbf{\underline{Tesis de Church:}} Toda funcion \(\Sigma \)-efectivamente computable es \(\Sigma \)-recursiva.


\textbf{\underline{Corollary 67:}} Si \(f:D_{f}\subseteq \omega ^{n}\times \Sigma ^{\ast m}\rightarrow O\) es \( \Sigma \)-recursiva, entonces existe un predicado \(\Sigma \)-p.r. \(P:\mathbf{N} \times \omega ^{n}\times \Sigma ^{\ast m}\rightarrow \omega \) y una funcion \( \Sigma \)-p.r. \(g:\mathbf{N}\rightarrow O\) tales que \(f=g\circ M(P).\)

\textbf{\underline{Proof:}} Supongamos que \(O=\Sigma ^{\ast }\). Sea \(\mathcal{P}_{0}\) un programa el cual compute a \(f\). Sea \(< \) un orden total estricto sobre \(\Sigma \). Note que podemos tomar

\(\displaystyle \begin{array}{rcl} P & =& \lambda t\vec{x}\vec{\alpha}[i^{n,m}\left( (t)_{1},\vec{x},\vec{\alpha}, \mathcal{P}_{0}\right) =n(\mathcal{P}_{0})+1\wedge (t)_{2}=\#^{< }(E_{\ast 1}^{n,m}((t)_{1},\vec{x},\vec{\alpha},\mathcal{P}_{0}))] \\ g & =& \lambda t\left[ \ast ^{< }((t)_{2})\right] \text{.} \end{array} \)

(Justifique por que \(P\) es \(\Sigma \)-p.r..) \(\Box\)


\textbf{\underline{Lemma 68:}} Supongamos \(f_{i}:D_{f_{i}}\subseteq \omega ^{n}\times \Sigma ^{\ast m}\rightarrow O\), \(i=1,...,k\), son funciones \(\Sigma \)-recursivas tales que \(D_{f_{i}}\cap D_{f_{j}}=\varnothing \) para \(i\neq j\). Entonces la funcion \(f_{1}\cup ...\cup f_{k}\) es \(\Sigma \)-recursiva.

\textbf{\underline{Proof:}} Probaremos el caso \(k=2\) y \(O=\Sigma ^{\ast }\). Sean \(\mathcal{P}_{1}\) y \( \mathcal{P}_{2}\) programas que computen las funciones \(f_{1}\) y \(f_{2}\), respectivamente. Sean

\(\displaystyle \begin{array}{rcl} P_{1} & =& \lambda t\vec{x}\vec{\alpha}\left[ i^{n,m}(t,\vec{x},\vec{\alpha}, \mathcal{P}_{1})=n(\mathcal{P}_{1})+1\right] \\ P_{2} & =& \lambda t\vec{x}\vec{\alpha}\left[ i^{n,m}(t,\vec{x},\vec{\alpha}, \mathcal{P}_{2})=n(\mathcal{P}_{2})+1\right] \end{array} \)

Notese que \(D_{P_{1}}=D_{P_{2}}=\omega \times \omega ^{n}\times \Sigma ^{\ast m}\) y que \(P_{1}\) y \(P_{2}\) son \((\Sigma \cup \Sigma _{p})\)-p.r.. Ya que son \(\Sigma \)-mixtos, el Teorema 51 nos dice que son \( \Sigma \)-p.r.. Tambien notese que \(D_{M((P_{1}\vee P_{2}))}=D_{f_{1}}\cup D_{f_{2}}\). Definamos
\(\displaystyle \begin{array}{rcl} g_{1} & =& \lambda \vec{x}\vec{\alpha}\left[ E_{\ast 1}^{n,m}(M\left( (P_{1}\vee P_{2})\right) (\vec{x},\vec{\alpha}),\vec{x},\vec{\alpha}, \mathcal{P}_{1})^{P_{i}(M\left( (P_{1}\vee P_{2})\right) (\vec{x},\vec{\alpha }),\vec{x},\vec{\alpha})}\right] \\ g_{2} & =& \lambda \vec{x}\vec{\alpha}\left[ E_{\ast 1}^{n,m}(M\left( (P_{1}\vee P_{2})\right) (\vec{x},\vec{\alpha}),\vec{x},\vec{\alpha}, \mathcal{P}_{2})^{P_{i}(M\left( (P_{1}\vee P_{2})\right) (\vec{x},\vec{\alpha }),\vec{x},\vec{\alpha})}\right] \end{array} \)

Notese que \(g_{1}\) y \(g_{2}\) son \(\Sigma \)-recursivas y que \( D_{g_{1}}=D_{g_{2}}=D_{f_{1}}\cup D_{f_{2}}\), Ademas notese que
\(\displaystyle g_{1}(\vec{x},\vec{\alpha})=\left\{ \begin{array}{lll} f_{1}(\vec{x},\vec{\alpha}) & & \text{si }(\vec{x},\vec{\alpha})\in D_{f_{1}} \\ \varepsilon & & \text{caso contrario} \end{array} \right. \)

\(\displaystyle g_{2}(\vec{x},\vec{\alpha})=\left\{ \begin{array}{lll} f_{2}(\vec{x},\vec{\alpha}) & & \text{si }(\vec{x},\vec{\alpha})\in D_{f_{2}} \\ \varepsilon & & \text{caso contrario} \end{array} \right. \)

O sea que \(f_{1}\cup f_{2}=\lambda \alpha \beta \left[ \alpha \beta \right] \circ (g_{1},g_{2})\) es \(\Sigma \)-recursiva. \(\Box\)


\textbf{\underline{Lemma 69:}} Supongamos \(\Sigma \supseteq \Sigma _{p}\). Entonces \( Halt^{\Sigma }\) es no \(\Sigma \)-recursivo.

\textbf{\underline{Proof:}} Supongamos \(Halt^{\Sigma }\) es \(\Sigma \)-recursivo y por lo tanto \(\Sigma \) -computable. Por la proposicion de existencia de macros tenemos que hay un macro

\(\displaystyle \left[ \mathrm{IF}\;Halt^{\Sigma }(\mathrm{W}1)\;\mathrm{GOTO}\;\mathrm{A}1 \right] \)

Sea \(\mathcal{P}_{0}\) el siguiente programa de \(\mathcal{S}^{\Sigma }\)
\(\displaystyle \mathrm{L}1\;\left[ \mathrm{IF}\;Halt^{\Sigma }(\mathrm{P}1)\;\mathrm{GOTO}\; \mathrm{L}1\right] \)

Note que
- \(\mathcal{P}_{0}\) termina partiendo desde \(\left( (0,0,...),( \mathcal{P}_{0},\varepsilon ,\varepsilon ,...)\right) \) sii \(Halt^{\Sigma }( \mathcal{P}_{0})=0\),
lo cual produce una contradiccion si tomamos en (*) \(\mathcal{P}= \mathcal{P}_{0}\). \(\Box\)

\section{Conjuntos $\Sigma$-recursivamente enumerables}

\textbf{\underline{Theorem 70:}} Sea \(S\subseteq \omega ^{n}\times \Sigma ^{\ast m}\). Entonces \(S\) es \(\Sigma \)-efectivamente enumerable sii \(S\) es \(\Sigma \)-recursivamente enumerable


\textbf{\underline{Proof:}} (\(\Rightarrow \)) Use la Tesis de Church.

(\(\Leftarrow \)) Use el Theorem 42. \(\Box\)



\textbf{\underline{Theorem 71:}} Dado \(S\subseteq \omega ^{n}\times \Sigma ^{\ast m} \), son equivalentes
(1) \(S\) es \(\Sigma \)-recursivamente enumerable
(2) \(S=I_{F}\), para alguna \(F:D_{F}\subseteq \omega ^{k}\times \Sigma ^{\ast l}\rightarrow \omega ^{n}\times \Sigma ^{\ast m}\) tal que cada \(F_{i}\) es \(\Sigma \)-recursiva.
(3) \(S=D_{f}\), para alguna funcion \(\Sigma \)-recursiva \(f\)
(4) \(S=\varnothing \) o \(S=I_{F}\), para alguna \(F:\omega \rightarrow \omega ^{n}\times \Sigma ^{\ast m}\) tal que cada \(F_{i}\) es \(\Sigma \)-p.r.


\textbf{\underline{Proof:}} (2)\(\Rightarrow \)(3). Para \(i=1,...,n+m\), sea \(\mathcal{P}_{i}\) un programa el cual computa a \(F_{i}\) y sea \(< \) un orden total estricto sobre \(\Sigma \). Sea \(P:\mathbf{N}\times \omega ^{n}\times \Sigma ^{\ast m}\rightarrow \omega \) dado por \(P(t,\vec{x},\vec{\alpha})=1\) sii se cumplen las siguientes condiciones

\(\displaystyle \begin{array}{rcl} i^{k,l}(\left( (t)_{k+l+1},(t)_{1},...,(t)_{k},\ast ^{< }((t)_{k+1}),...,\ast ^{< }((t)_{k+l})),\mathcal{P}_{1}\right) & =& n(\mathcal{P}_{1})+1 \\ & & \vdots \\ i\left( (t)_{k+l+1},(t)_{1}...(t)_{k},\ast ^{< }((t)_{k+1})...\ast ^{< }((t)_{k+l})),\mathcal{P}_{n+m}\right) & =& n(\mathcal{P}_{n+m})+1 \\ E_{\#1}^{k,l}((t)_{k+l+1},(t)_{1},...,(t)_{k},\ast ^{< }((t)_{k+1}),...,\ast ^{< }((t)_{k+l})),\mathcal{P}_{1}) & =& x_{1} \\ & & \vdots \\ E_{\#1}^{k,l}((t)_{k+l+1},(t)_{1},...,(t)_{k},\ast ^{< }((t)_{k+1}),...,\ast ^{< }((t)_{k+l})),\mathcal{P}_{n}) & =& x_{n} \\ E_{\ast 1}^{k,l}((t)_{k+l+1},(t)_{1},...,(t)_{k},\ast ^{< }((t)_{k+1}),...,\ast ^{< }((t)_{k+l})),\mathcal{P}_{n+1}) & =& \alpha _{1} \\ & & \vdots \\ E_{\ast 1}^{k,l}((t)_{k+l+1},(t)_{1},...,(t)_{k},\ast ^{< }((t)_{k+1}),...,\ast ^{< }((t)_{k+l})),\mathcal{P}_{n+m}) & =& \alpha _{m} \end{array} \)

Note que \(P\) es \((\Sigma \cup \Sigma _{p})\)-p.r. y por lo tanto \(P\) es \( \Sigma \)-p.r.. Pero entonces \(M(P)\) es \(\Sigma \)-r. lo cual nos dice que se cumple (3) ya que \(D_{M(P)}=I_{F}=S\).
(3)\(\Rightarrow \)(4). Supongamos \(S\neq \varnothing \). Sea \( (z_{1},...,z_{n},\gamma _{1},...,\gamma _{m})\in S\) fijo. Sea \(\mathcal{P}\) un programa el cual compute a \(f\) y sea \(< \) un orden total estricto sobre \( \Sigma \). Sea \(P:\mathbf{N}\rightarrow \omega \) dado por \(P(x)=1\) sii

\(\displaystyle i^{n,m}\left( (x)_{n+m+1},(x)_{1},...,(x)_{n},\ast ^{< }((x)_{n+1}),...,\ast ^{< }((x)_{n+m})),\mathcal{P}\right) =n(\mathcal{P})+1 \)

Es facil ver que \(P\) es \((\Sigma \cup \Sigma _{p})\)-p.r. por lo cual es \( \Sigma \)-p.r.. Sea \(\bar{P}=P\cup C_{0}^{1,0}\mid _{\{0\}}\). Para \(i=1,...,n\) , definamos \(F_{i}:\omega \rightarrow \omega \) de la siguiente manera
\(\displaystyle F_{i}(x)=\left\{ \begin{array}{ccc} (x)_{i} & \text{si} & \bar{P}(x)=1 \\ z_{i} & \text{si} & \bar{P}(x)\neq 1 \end{array} \right. \)

Para \(i=n+1,...,n+m\), definamos \(F_{i}:\omega \rightarrow \Sigma ^{\ast }\) de la siguiente manera
\(\displaystyle F_{i}(x)=\left\{ \begin{array}{lll} \ast ^{< }((x)_{i}) & \text{si} & \bar{P}(x)=1 \\ \gamma _{i-n} & \text{si} & \bar{P}(x)\neq 1 \end{array} \right. \)

Por el lema de division por casos, cada \(F_{i}\) es \(\Sigma \)-p.r.. Es facil ver que \(F=(F_{1},...,F_{n+m})\) cumple (4). \(\Box\)


\textbf{\underline{Corollary 72:}} Supongamos \(f:D_{f}\subseteq \omega ^{n}\times \Sigma ^{\ast m}\rightarrow O\) es \(\Sigma \)-recursiva y \(S\subseteq D_{f}\) es \( \Sigma \)-r.e., entonces \(f\mid _{S}\) es \(\Sigma \)-recursiva.


\textbf{\underline{Proof:}} Supongamos \(O=\Sigma ^{\ast }.\) Por el Theorem anterior \(S=D_{g}\), para alguna funcion \(\Sigma \)-recursiva \(g.\) Notese que componiendo adecuadamente podemos suponer que \(I_{g}=\{\varepsilon \}.\) O sea que tenemos \(f\mid _{S}=\lambda \alpha \beta \left[ \alpha \beta \right] \circ (f,g)\). \(\Box\)


\textbf{\underline{Corollary 73:}} Supongamos \(f:D_{f}\subseteq \omega ^{n}\times \Sigma ^{\ast m}\rightarrow O\) es \(\Sigma \)-recursiva y \(S\subseteq I_{f}\) es \(\Sigma \)-r.e., entonces \( f^{-1}(S)=\{(\vec{x},\vec{\alpha}):f(\vec{x},\vec{\alpha})\in S\}\) es \( \Sigma \)-r.e..

\textbf{\underline{Proof:}} Por el Theorem anterior \(S=D_{g}\), para alguna funcion \(\Sigma \)-recursiva \( g \). O sea que \(f^{-1}(S)=D_{g\circ f}\) es \(\Sigma \)-r.e.. \(\Box\)


\textbf{\underline{Corollary 74:}} Supongamos \(S_{1},S_{2}\subseteq \omega ^{n}\times \Sigma ^{\ast m}\) son conjuntos \(\Sigma \)-r.e.. Entonces \(S_{1}\cap S_{2}\) es \(\Sigma \)-r.e..

\textbf{\underline{Proof:}} Por el Theorem anterior \(S_{i}=D_{g_{i}}\), con \(g_{1},g_{2}\) funciones \( \Sigma \)-recursivas\(.\) Notese que podemos suponer que \(I_{g_{1}},I_{g_{2}} \subseteq \omega \) por lo que \(S_{1}\cap S_{2}=D_{\lambda xy\left[ xy\right] \circ (g_{1},g_{2})}\) es \(\Sigma \)-r.e.\(.\) \(\Box\)


\textbf{\underline{Corollary 75:}} Supongamos \(S_{1},S_{2}\subseteq \omega ^{n}\times \Sigma ^{\ast m}\) son conjuntos \(\Sigma \)-r.e.. Entonces \(S_{1}\cup S_{2}\) es \(\Sigma \)-r.e.

\textbf{\underline{Proof:}} Supongamos \(S_{1}\neq \varnothing \neq S_{2}.\) Sean \(F,G:\omega \rightarrow \omega ^{n}\times \Sigma ^{\ast m}\) tales que \(I_{F}=S_{1}\), \(I_{G}=S_{2}\) y las funciones \(F_{i} {\acute{}} s\) y \(G_{i} {\acute{}} s\) son \(\Sigma \)-recursivas. Sean \(f=\lambda x\left[ Q(x,2)\right] \) y \( g=\lambda x\left[ Q(x\dot{-}1,2)\right] .\) Sea \(H:\omega \rightarrow \omega ^{n}\times \Sigma ^{\ast m}\) dada por

\(\displaystyle H_{i}=(F_{i}\circ f)\mathrm{\mid }_{\{x:x\mathrm{\ es\ par}\}}\cup (G_{i}\circ g)\mathrm{\mid }_{\{x:x\mathrm{\ es\ impar}\}} \)

Por el Corollary 72 y el Lema 68, cada \(H_{i}\) es \( \Sigma \)-recursiva. Ya que \(I_{H}=S_{1}\cup S_{2}\).tenemos que \(S_{1}\cup S_{2}\) es \(\Sigma \)-r.e. \(\Box\)


\textbf{\underline{Theorem 76:}} Sea \(S\subseteq \omega ^{n}\times \Sigma ^{\ast m}\). Entonces \(S\) es \(\Sigma \)-efectivamente computable sii \(S\) es \(\Sigma \)-recursivo
\textbf{\underline{Proof:}} (\(\Rightarrow \)) Use la Tesis de Church.

(\(\Leftarrow \)) Use el Teorema 42. \(\Box\)




\textbf{\underline{Theorem 77:}} Sea \(S\subseteq \omega ^{n}\times \Sigma ^{\ast m}.\) Son equivalentes
(a) \(S\) es \(\Sigma \)-recursivo
(b) \(S\) y \((\omega ^{n}\times \Sigma ^{\ast m})-S\) son \(\Sigma \) -recursivamente enumerables


\textbf{\underline{Proof:}} (a)\(\Rightarrow \)(b)\(.\) Note que \(S=D_{Pred\circ \chi _{S}}.\)

(b)\(\Rightarrow \)(a). Note que \(\chi _{S}=C_{1}^{n,m}\mathrm{\mid }_{S}\cup C_{0}^{n,m}\mathrm{\mid }_{\omega ^{n}\times \Sigma ^{\ast m}-S}\). \(\Box\)

\textbf{\underline{Lemma 78:}} Supongamos que \(\Sigma \supseteq \Sigma _{p}.\) Entonces
\(\displaystyle A=\left\{ \mathcal{P}\in \mathrm{Pro}^{\Sigma }:Halt^{\Sigma }(\mathcal{P} )\right\} \)

es \(\Sigma \)-r.e. y no es \(\Sigma \)-recursivo. Mas aun el conjunto
\(\displaystyle N=\left\{ \mathcal{P}\in \mathrm{Pro}^{\Sigma }:\lnot Halt^{\Sigma }( \mathcal{P})\right\} \)
no es \(\Sigma \)-r.e.


\textbf{\underline{Proof:}} Sea \(P=\lambda t\mathcal{P}\left[ i^{0,1}(t,\mathcal{P},\mathcal{P})=n( \mathcal{P})+1\right] \). Note que \(P\) es \(\Sigma \)-p.r. por lo que \(M(P)\) es \(\Sigma \)-r.. Ademas note que \(D_{M(P)}=A\), lo cual implica que \(A\) es \( \Sigma \)-r.e.. Ya que \(Halt^{\Sigma }\) es no \(\Sigma \)-recursivo (Lema 69) y

\(\displaystyle Halt^{\Sigma }=C_{1}^{0,1}\mid _{A}\cup C_{0}^{0,1}\mid _{N} \)

el Lema 68 nos dice que \(N\) no es \(\Sigma \)-r.e.. Finalmente supongamos \(A\) es \(\Sigma \)-recursivo. Entonces el conjunto
\(\displaystyle N=\left( \Sigma ^{\ast }-A\right) \cap \mathrm{Pro}^{\Sigma } \)

deberia serlo, lo cual es absurdo. \(\Box\)

\section{Maquinas de Turing}


\textbf{\underline{Lemma 79:}} Sea \(L\subseteq \Sigma ^{\ast }.\) entonces \(L=L(M)\) para alguna maquina de Turing \(M\) sii \(L=H(M)\) para alguna maquina de Turing \(M=(Q,\Sigma ,\Gamma ,\delta ,q_{0},B,F)\).

\textbf{\underline{Proof:}} (\(\Rightarrow \)) Dada una maquina \(M=(Q,\Sigma ,\Gamma ,\delta ,q_{0},B,F)\), costruiremos una maquina \(M_{1}=(Q_{1},\Sigma ,\Gamma _{1},\delta _{1}, \tilde{q}_{0},B,\varnothing )\) tal que \(L(M)=H(M_{1}).\) Tomaremos \(\Gamma _{1}=\Gamma \cup \{X\}\), con \(X\) un simbolo nuevo no perteneciente a \(\Gamma \). Para cada \(a\in \Sigma \), sea \(q_{a}\) un estado nuevo, no perteneciente a \(Q.\) Sean \(\tilde{q}_{0},q_{r},q_{d},q_{B}\) estados nuevos no pertenecientes a \(Q.\) Tomemos entonces

\(\displaystyle Q_{1}=Q\cup \{\tilde{q}_{0},q_{r},q_{d},q_{B}\}\cup \{q_{a}:a\in \Sigma \} \)

Finalmente definamos \(\delta _{1}\) de la siguiente manera:
\(\displaystyle \begin{array}{rcl} \delta _{1}(\tilde{q}_{0},B) & =& \{(q_{B},X,R)\} \\ \delta _{1}(q_{B},a) & =& \{(q_{a},B,R)\}\text{, para }a\in \Sigma \\ \delta _{1}(q_{B},B) & =& \{(q_{0},B,K)\} \\ \delta _{1}(q_{a},b) & =& \{(q_{b},a,R)\}\text{, para }a,b\in \Sigma \\ \delta _{1}(q_{a},B) & =& \{(q_{r},a,L)\}\text{, para }a\in \Sigma \\ \delta _{1}(q_{r},a) & =& \{(q_{r},a,L)\}\text{, para }a\in \Sigma \\ \delta _{1}(q_{r},B) & =& \{(q_{0},B,K)\} \\ \delta _{1}(q,X) & =& \{(q,X,K)\}\text{, para }q\in Q \\ \delta _{1}(q,\sigma ) & =& \delta (q,\sigma )\cup \{(q_{d},\sigma ,K)\}\text{ , para }q\in F\text{ y }\sigma \in \Gamma \\ \delta _{1}(q,\sigma ) & =& \delta (q,\sigma )\text{, para }q\in Q-F\text{ y } \sigma \in \Gamma \\ \delta _{1}(q_{d},\sigma ) & =& \varnothing \text{, para }\sigma \in \Gamma \end{array} \)

(\(\delta _{1}\) se define igual a vacio para los casos no contemplados arriba).
(\(\Leftarrow \)) Dada \(M=(Q,\Sigma ,\Gamma ,\delta ,q_{0},B,F)\), dejamos al lector la construccion de una maquina \(M_{1}=(Q_{1},\Sigma ,\Gamma _{1},\delta _{1},\tilde{q}_{0},B,\varnothing )\) tal que \(H(M)=L(M_{1})\). \(\Box\)


\textbf{\underline{Lemma 80:}} El predicado \(\lambda ndd^{\prime }\left[ d\vdash d^{\prime }\right] \) es \( (\Gamma \cup Q)\)-p.r..

\textbf{\underline{Proof:}} Note que \(D_{\lambda dd^{\prime }\left[ d\vdash d^{\prime }\right] }=Des\times Des\). Tambien notese que los predicados

\(\displaystyle \begin{array}{rcl} & & \lambda p\sigma q\gamma \left[ (p,\sigma ,L)\in \delta (q,\gamma )\right] \\ & & \lambda p\sigma q\gamma \left[ (p,\sigma ,R)\in \delta (q,\gamma )\right] \\ & & \lambda p\sigma q\gamma \left[ (p,\sigma ,K)\in \delta (q,\gamma )\right] \end{array} \)

son \((\Gamma \cup Q)\)-p.r. ya que los tres tienen dominio igual a \(Q\times \Gamma \times Q\times \Gamma \) el cual es finito (Corolario 36 ). Sea \(P_{R}:Des\times Des\times \Gamma \times \Gamma ^{\ast }\times \Gamma ^{\ast }\times Q\times Q\rightarrow \omega \) definido por \(P_{R}(d,d^{\prime },\sigma ,\alpha ,\beta ,p,q)=1\) sii
\(\displaystyle d=\alpha p\beta \wedge (q,\sigma ,R)\in \delta \left( p,\left[ \beta B\right] _{1}\right) \wedge d^{\prime }=\alpha \sigma q^{\curvearrowright }\beta \)

Sea \(P_{L}:Des\times Des\times \Gamma \times \Gamma ^{\ast }\times \Gamma ^{\ast }\times Q\times Q\rightarrow \omega \) definido por \(P_{L}(d,d^{\prime },\sigma ,\alpha ,\beta ,p,q)=1\) sii
\(\displaystyle d=\alpha p\beta \wedge (q,\sigma ,L)\in \delta \left( p,\left[ \beta B\right] _{1}\right) \wedge \alpha \neq \varepsilon \wedge d^{\prime }=\left\lfloor \alpha ^{\curvearrowleft }q\left[ \alpha \right] _{\left\vert \alpha \right\vert }\sigma ^{\curvearrowright }\beta \right\rfloor \)

Sea \(P_{K}:Des\times Des\times \Gamma \times \Gamma ^{\ast }\times \Gamma ^{\ast }\times Q\times Q\rightarrow \omega \) definido por \(P_{K}(d,d^{\prime },\sigma ,\alpha ,\beta ,p,q)=1\) sii
\(\displaystyle d=\alpha p\beta \wedge (q,\sigma ,K)\in \delta \left( p,\left[ \beta B\right] _{1}\right) \wedge d^{\prime }=\left\lfloor \alpha q\sigma ^{\curvearrowright }\beta \right\rfloor \)

Se deja al lector la verificacion de que estos predicados son \((\Gamma \cup Q)\)-p.r.. Notese que \(\lambda dd^{\prime }\left[ d\vdash d^{\prime }\right] \) es igual al predicado
\(\displaystyle \lambda dd^{\prime }\left[ (\exists \sigma \in \Gamma )(\exists \alpha ,\beta \in \Gamma ^{\ast })(\exists p,q\in Q)(P_{R}\vee P_{L}\vee P_{K})(d,d^{\prime },\sigma ,\alpha ,\beta ,p,q)\right] \)

lo cual por el Lema 39 nos dice que \(\lambda dd^{\prime } \left[ d\vdash d^{\prime }\right] \) es \((\Gamma \cup Q)\)-p.r. \(\Box\)


\textbf{\underline{Proposition 81:}} \(\lambda ndd^{\prime }\left[ d\overset{n}{\vdash }d^{\prime }\right] \) es \( (\Gamma \cup Q)\)-p.r..

\textbf{\underline{Proof:}} Sea \(Q=\lambda dd^{\prime }\left[ d\vdash d^{\prime }\right] \cup C_{0}^{0,2}\mid _{(\Gamma \cup Q)^{\ast 2}-Des^{2}}\) es decir \(Q\) es el resultado de extender con el valor \(0\) al predicado \(\lambda dd^{\prime } \left[ d\vdash d^{\prime }\right] \) de manera que este definido en todo \( (\Gamma \cup Q)^{\ast 2}\). Sea \(< \) un orden total estricto sobre \(\Gamma \cup Q\) y sea \(Q_{1}:\mathbf{N}\times Des\times Des\rightarrow \omega \) definido por \(Q_{1}(x,d,d^{\prime })=1\) sii

\(\left( (\forall i\in \mathbf{N})_{i\leq Lt(x)}\ast ^{< }((x)_{i})\in Des\right) \wedge \ast ^{< }((x)_{1})=d\wedge \)

\(\ \ \ \ \ \ \ \ \ \ \ \ \ \ \ \ \ \ \ \ \ \ \ast ^{< }((x)_{Lt(x)})=d^{\prime }\wedge \left( (\forall i\in \mathbf{N})_{i\leq Lt(x)\dot{-}1}\;Q(\ast ^{< }((x)_{i}),\ast ^{< }((x)_{i+1}))\right) \)

Notese que dicho rapidamente \(Q_{1}(x,d,d^{\prime })=1\) sii \(x\) codifica una computacion que parte de \(d\) y llega a \(d^{\prime }\). Se deja al lector la verificacion de que este predicado es \((\Gamma \cup Q)\)-p.r.. Notese que

\(\displaystyle \lambda ndd^{\prime }\left[ d\overset{n}{\vdash }d^{\prime }\right] =\lambda ndd^{\prime }\left[ \left( \exists x\in \mathbf{N}\right) \;Lt(x)=n+1\wedge Q_{1}(x,d,d^{\prime })\right] \)

Es decir que solo nos falta acotar el cuantificador existencial, para poder aplicar el lema de cuantificacion acotada. Ya que cuando \( d_{1},...,d_{n+1}\in Des\) son tales que \(d_{1}\vdash d_{2}\vdash ...\vdash d_{n+1}\) tenemos que
\(\displaystyle \left\vert d_{i}\right\vert \leq \left\vert d_{1}\right\vert +n\text{, para } i=1,...,n \)

una posible cota para dicho cuantificador es
\(\displaystyle \prod_{i=1}^{n+1}pr(i)^{\left\vert \Gamma \cup Q\right\vert ^{\left\vert d\right\vert +n}}\text{.} \)

O sea que, por el lema de cuantificacion acotada, tenemos que el predicado \( \lambda ndd^{\prime }\left[ d\overset{n}{\vdash }d^{\prime }\right] \) es \( (\Gamma \cup Q)\)-p.r. \(\Box\)

\textbf{\underline{Theorem 82:}} Sea \(M=\left( Q,\Sigma ,\Gamma ,\delta ,q_{0},B,F\right) \) una maquina de Turing. Entonces \(L(M)\) es \(\Sigma \)-recursivamente enumerable.


\textbf{\underline{Proof:}} Sea \(P\) el siguiente predicado \((\Gamma \cup Q)\)-mixto

\(\displaystyle \lambda n\alpha \left[ (\exists d\in Des)\;\left\lfloor q_{0}B\alpha \right\rfloor \overset{n}{\vdash }d\wedge St(d)\in F\right] \)

Notese que \(D_{P}=\omega \times \Gamma ^{\ast }\). Dejamos al lector probar que \(P\) es \((\Gamma \cup Q)\)-p.r.. Sea \(P^{\prime }=P\mid _{\omega \times \Sigma ^{\ast }}\). Notese que \(P^{\prime }(n,\alpha )=1\) sii \(\alpha \in L(M) \) atestiguado por una computacion de longitud \(n\). Ya que \(P^{\prime }\) es \((\Gamma \cup Q)\)-p.r. (por que?) y ademas es \(\Sigma \)-mixto, el Teorema 51 nos dice que \(P^{\prime }\) es \(\Sigma \)-p.r.. Ya que \( L(M)=D_{M(P^{\prime })}\), el Teorema 71 nos dice que \( L(M)\) es \(\Sigma \)-r.e.. \(\Box\)

\section{Funciones $\Sigma$-Turing computables}

\textbf{\underline{Theorem 83:}} Supongamos \(f:S\subseteq \omega ^{n}\times \Sigma ^{\ast }{}^{m}\rightarrow O \) es \(\Sigma \)-Turing computable. Entonces \(f\) es \(\Sigma \)-recursiva.

\textbf{\underline{Proof:}} Supongamos \(O=\Sigma ^{\ast }\) y sea \(M=\left( Q,\Sigma ,\Gamma ,\delta ,q_{0},B,\shortmid ,F\right) \) una maquina de Turing deterministica con unit la cual compute a \(f\). Sea \(< \) un orden total estricto sobre \(\Gamma \cup Q\) . Sea \(P:\mathbf{N}\times \omega ^{n}\times \Sigma ^{\ast m}\rightarrow \omega \) dado por \(P(x,\vec{x},\vec{\alpha})=1\) sii

\((\exists q\in Q)\;\left\lfloor q_{0}B\shortmid ^{x_{1}}...B\shortmid ^{x_{n}}B\alpha _{1}...B\alpha _{m}\right\rfloor \overset{(x)_{1}}{\vdash } \left\lfloor qB\ast ^{< }((x)_{2})\right\rfloor \wedge \)
\(\ \ \ \ \ \ \ \ \ \ \ \ \ \ \ \ \ \ \ \ \ \ \ \ \ \ \ \ \ \ \ \ \ \ \ \ \ \ \ \ \ \ \ \ \ \ \ \ \ \ \ \ \ \ \ \ \ \ \ \ \ \ \ \ \ \ \ \ \ \ \ \ \ \ \ \ \ \wedge (\forall d\in Des)_{\left\vert d\right\vert \leq \left\vert \ast ^{< }((x)_{2})\right\vert +2}\;\left\lfloor qB\ast ^{< }((x)_{2})\right\rfloor \nvdash d\)
Es facil ver que \(P\) es \((\Gamma \cup Q)\)-p.r. por lo que \(P\) es \( \Sigma \)-p.r. ya que es \(\Sigma \)-mixto. Notese que

\(\displaystyle f=\lambda \vec{x}\vec{\alpha}\left[ \left( \min_{x}P(x,\vec{x},\vec{\alpha} )\right) _{2}\right] \text{,} \)

lo cual nos dice que \(f\) es \(\Sigma \)-recursiva. \(\Box\)


\textbf{\underline{Lemma 84:}} Sea \(\mathcal{P}\in \mathrm{Pro}^{\Sigma }\) y sea \(k\) tal que las variables que ocurren en \(\mathcal{P}\) estan todas en la lista \(\mathrm{N}1,..., \mathrm{N}\bar{k},\mathrm{P}1,...,\mathrm{P}\bar{k}.\) Para cada \(a\in \Sigma \cup \{\shortmid \}\), sea \(\tilde{a}\) un nuevo simbolo. Sea \(\Gamma =\Sigma \cup \{B,\shortmid \}\cup \{\tilde{a}:a\in \Sigma \cup \{\shortmid \}\}\). Entonces hay una maquina de Turing deterministica con unit \(M=\left( Q,\Gamma ,\Sigma ,\delta ,q_{0},B,\shortmid ,\{q_{f}\}\right) \) la cual satisface
(1) \(\delta (q_{f},\sigma )=\varnothing \), para cada \(\sigma \in \Gamma \).
(2) Cualesquiera sean \(x_{1},...,x_{k}\in \omega \) y \(\alpha _{1},...,\alpha _{k}\in \Sigma ^{\ast }\), el programa \(\mathcal{P}\) se detiene partiendo del estado
\(\displaystyle \left( (x_{1},...,x_{k},0,...),(\alpha _{1},...,\alpha _{k},\varepsilon ,...)\right) \)

sii \(M\) se detiene partiendo de la descripcion instantanea
\(\displaystyle \left\lfloor q_{0}B\shortmid ^{x_{1}}B...B\shortmid ^{x_{k}}B\alpha _{1}B...B\alpha _{k}B\right\rfloor \)
(3) Si \(x_{1},...,x_{k}\in \omega \) y \(\alpha _{1},...,\alpha _{k}\in \Sigma ^{\ast }\) son tales que \(\mathcal{P}\) se detiene partiendo del estado
\(\displaystyle \left( (x_{1},...,x_{k},0,...),(\alpha _{1},...,\alpha _{k},\varepsilon ,...)\right) \)

y llega al estado
\(\displaystyle \left( (y_{1},...,y_{k},0,...),(\beta _{1},...,\beta _{k},\varepsilon ,...)\right) \)

entonces
\(\displaystyle \left\lfloor q_{0}B\shortmid ^{x_{1}}B...B\shortmid ^{x_{k}}B\alpha _{1}B...B\alpha _{k}B\right\rfloor \overset{\ast }{\vdash }\left\lfloor q_{f}B\shortmid ^{y_{1}}B...B\shortmid ^{y_{k}}B\beta _{1}B...B\beta _{k}B\right\rfloor \)

\textbf{\underline{Proof:}} Dado un estado \(((x_{1},...,x_{k},0,...),(\alpha _{1},...,\alpha _{k},\varepsilon ,...))\) lo representaremos en la cinta de la siguiente manera


\(\displaystyle B\shortmid ^{x_{1}}...B\shortmid ^{x_{k}}B\alpha _{1}...B\alpha _{k}BBBB.... \)

A continuacion describiremos una serie de maquinas las cuales simularan, via la representacion anterior, las distintas clases de instrucciones que pueden ocurrir en \(\mathcal{P}\). Todas las maquinas definidas tendran a \(\shortmid \) como unit y a \(B\) como blanco, tendran a \(\Sigma \) como su alfabeto terminal y su alfabeto mayor sera \(\Gamma =\Sigma \cup \{B,\shortmid \}\cup \{\tilde{a }:a\in \Sigma \cup \{\shortmid \}\}\). Ademas tendran uno o dos estados finales con la propiedad de que si \(q\) es un estado final, entonces \(\delta (q,\sigma )=\varnothing \), para cada \(\sigma \in \Gamma \). Esta propiedad es importante ya que nos permitira concatenar pares de dichas maquinas identificando algun estado final de la primera con el inicial de la segunda.
Para \(1\leq i\leq k\), sea \(M_{i}^{+}\) una maquina tal que

\(\displaystyle \begin{array}{lcl} B\shortmid ^{x_{1}}...B\shortmid ^{x_{k}}B\alpha _{1}...B\alpha _{k} & \overset{\ast }{\vdash } & B\shortmid ^{x_{1}}...B\shortmid ^{x_{i-1}}B\shortmid ^{x_{i}+1}B\shortmid ^{x_{i+1}}...B\shortmid ^{x_{k}}B\alpha _{1}...B\alpha _{k} \\ \uparrow & & \uparrow \\ q_{0} & & q_{f} \end{array} \)

Es claro que la maquina \(M_{i}^{+}\) simula la instruccion \(\mathrm{N}\bar{ \imath}\leftarrow \mathrm{N}\bar{\imath}+1\).
Para \(1\leq i\leq k\), sea \(M_{i}^{\dot{-}}\) una maquina tal que

\(\displaystyle \begin{array}{lcl} B\shortmid ^{x_{1}}...B\shortmid ^{x_{k}}B\alpha _{1}...B\alpha _{k} & \overset{\ast }{\vdash } & B\shortmid ^{x_{1}}...B\shortmid ^{x_{i-1}}B\shortmid ^{x_{i}\dot{-}1}B\shortmid ^{x_{i+1}}...B\shortmid ^{x_{k}}B\alpha _{1}...B\alpha _{k} \\ \uparrow & & \uparrow \\ q_{0} & & q_{f} \end{array} \)

Para \(1\leq i\leq k\) y \(a\in \Sigma \), sea \(M_{i}^{a}\) una maquina tal que
\(\displaystyle \begin{array}{lcl} B\shortmid ^{x_{1}}...B\shortmid ^{x_{k}}B\alpha _{1}...B\alpha _{k} & \overset{\ast }{\vdash } & B\shortmid ^{x_{1}}...B\shortmid ^{x_{k}}B\alpha _{1}...B\alpha _{i-1}B\alpha _{i}aB\alpha _{i+1}...B\alpha _{k} \\ \uparrow & & \uparrow \\ q_{0} & & q_{f} \end{array} \)

Para \(1\leq i\leq k\), sea \(M_{i}^{\curvearrowright }\) una maquina tal que
\(\displaystyle \begin{array}{lcl} B\shortmid ^{x_{1}}...B\shortmid ^{x_{k}}B\alpha _{1}...B\alpha _{k} & \overset{\ast }{\vdash } & B\shortmid ^{x_{1}}...B\shortmid ^{x_{k}}B\alpha _{1}...B\alpha _{i-1}B^{\curvearrowright }\alpha _{i}B\alpha _{i+1}...B\alpha _{k} \\ \uparrow & & \uparrow \\ q_{0} & & q_{f} \end{array} \)

Para \(j=1,...,k\), y \(a\in \Sigma \), sea \(IF_{j}^{a}\) una maquina con dos estados finales \(q_{si}\) y \(q_{no}\) tal que si \(\alpha _{j}\) comienza con \(a\) , entonces
\(\displaystyle \begin{array}{lcl} B\shortmid ^{x_{1}}...B\shortmid ^{x_{k}}B\alpha _{1}...B\alpha _{k} & \overset{\ast }{\vdash } & B\shortmid ^{x_{1}}...B\shortmid ^{x_{k}}B\alpha _{1}...B\alpha _{k} \\ \uparrow & & \uparrow \\ q_{0} & & q_{si} \end{array} \)

y en caso contrario
\(\displaystyle \begin{array}{lcl} B\shortmid ^{x_{1}}...B\shortmid ^{x_{k}}B\alpha _{1}...B\alpha _{k} & \overset{\ast }{\vdash } & B\shortmid ^{x_{1}}...B\shortmid ^{x_{k}}B\alpha _{1}...B\alpha _{k} \\ \uparrow & & \uparrow \\ q_{0} & & q_{no} \end{array} \)

Analogamente para \(j=1,...,k\), sea \(IF_{j}\) una maquina tal que si \( x_{j}\neq 0\), entonces
\(\displaystyle \begin{array}{lcl} B\shortmid ^{x_{1}}...B\shortmid ^{x_{k}}B\alpha _{1}...B\alpha _{k} & \overset{\ast }{\vdash } & B\shortmid ^{x_{1}}...B\shortmid ^{x_{k}}B\alpha _{1}...B\alpha _{k} \\ \uparrow & & \uparrow \\ q_{0} & & q_{si} \end{array} \)

y si \(x_{j}=0\), entonces
\(\displaystyle \begin{array}{lcl} B\shortmid ^{x_{1}}...B\shortmid ^{x_{k}}B\alpha _{1}...B\alpha _{k} & \overset{\ast }{\vdash } & B\shortmid ^{x_{1}}...B\shortmid ^{x_{k}}B\alpha _{1}...B\alpha _{k} \\ \uparrow & & \uparrow \\ q_{0} & & q_{no} \end{array} \)

Para \(1\leq i,j\leq k\), sea \(M_{i\leftarrow j}^{\ast }\) una maquina tal que
\(\displaystyle \begin{array}{lcl} B\shortmid ^{x_{1}}...B\shortmid ^{x_{k}}B\alpha _{1}...B\alpha _{k} & \overset{\ast }{\vdash } & B\shortmid ^{x_{1}}...B\shortmid ^{x_{k}}B\alpha _{1}...B\alpha _{i-1}B\alpha _{j}B\alpha _{i+1}...B\alpha _{k} \\ \uparrow & & \uparrow \\ q_{0} & & q_{f} \end{array} \)

Para \(1\leq i,j\leq k\), sea \(M_{i\leftarrow j}^{\#}\) una maquina tal que
\(\displaystyle \begin{array}{lcl} B\shortmid ^{x_{1}}...B\shortmid ^{x_{k}}B\alpha _{1}...B\alpha _{k} & \overset{\ast }{\vdash } & B\shortmid ^{x_{1}}...B\shortmid ^{x_{i-1}}B\shortmid ^{x_{j}}B\shortmid ^{x_{i+1}}...B\shortmid ^{x_{k}}B\alpha _{1}...B\alpha _{k} \\ \uparrow & & \uparrow \\ q_{0} & & q_{f} \end{array} \)

Para \(1\leq i\leq k\), sea \(M_{i\leftarrow 0}\) una maquina tal que
\(\displaystyle \begin{array}{lcl} B\shortmid ^{x_{1}}...B\shortmid ^{x_{k}}B\alpha _{1}...B\alpha _{k} & \overset{\ast }{\vdash } & B\shortmid ^{x_{1}}...B\shortmid ^{x_{i-1}}BB\shortmid ^{x_{i+1}}...B\shortmid ^{x_{k}}B\alpha _{1}...B\alpha _{k} \\ \uparrow & & \uparrow \\ q_{0} & & q_{f} \end{array} \)

Para \(1\leq i\leq k\), sea \(M_{i\leftarrow \varepsilon }\) una maquina tal que
\(\displaystyle \begin{array}{lcl} B\shortmid ^{x_{1}}...B\shortmid ^{x_{k}}B\alpha _{1}...B\alpha _{k} & \overset{\ast }{\vdash } & B\shortmid ^{x_{1}}...B\shortmid ^{x_{k}}B\alpha _{1}...B\alpha _{i-1}BB\alpha _{i+1}...B\alpha _{k} \\ \uparrow & & \uparrow \\ q_{0} & & q_{f} \end{array} \)

Sea
\(\displaystyle M_{\mathrm{SKIP}}=\left( \{q_{0},q_{f}\},\Gamma ,\Sigma ,\delta ,q_{0},B,\shortmid ,\{q_{f}\}\right) , \)

con \(\delta (q_{0},B)=\{(q_{f},B,K)\}\) y \(\delta =\varnothing \) en cualquier otro caso.
Finalmente sea

\(\displaystyle M_{\mathrm{GOTO}}=\left( \{q_{0},q_{si},q_{no}\},\Gamma ,\Sigma ,\delta ,q_{0},B,\shortmid ,\{q_{si},q_{no}\}\right) , \)

con \(\delta (q_{0},B)=\{(q_{si},B,K)\}\) y \(\delta =\varnothing \) en cualquier otro caso.
Para poder hacer concretamente las maquinas recien descriptas deberemos dise \~{n}ar antes algunas maquinas auxiliares. Para cada \(j\geq 1\), sea \(D_{j}\) la maquina descripta en la Figura 1. Notese que

\(\displaystyle \begin{array}{lcr} \alpha B\beta _{1}B\beta _{2}B...B\beta _{j}B\gamma & \overset{\ast }{\vdash } & \alpha B\beta _{1}B\beta _{2}B...B\beta _{j}B\gamma \\ \ \ \uparrow & & \uparrow \ \ \\ \ \ q_{0} & & q_{f}\ \ \end{array} \)

siempre que \(\alpha ,\gamma \in \Gamma ^{\ast }\), \(\beta _{1},...,\beta _{j}\in (\Gamma -\{B\})^{\ast }\). Analogamente tenemos definidas las maquinas \(I_{j}.\)
Para \(j\geq 1\), sea \(TD_{j}\) una maquina con un solo estado final \(q_{f}\) y tal que

\(\displaystyle \begin{array}{ccc} \alpha B\gamma & \overset{\ast }{\vdash } & \alpha BB\gamma \\ \uparrow & & \uparrow \ \ \\ q_{0} & & q_{f}\ \ \end{array} \)

cada vez que \(\alpha ,\gamma \in \Gamma ^{\ast }\) y \(\gamma \) tiene exactamente \(j\) ocurrencias de \(B\). Es decir la maquina \(TD_{j}\) corre un espacio a la derecha todo el bloque \(\gamma \) y agrega un blanco en el espacio que se genera a la izquierda de dicho bloque. Por ejemplo, para el caso de \(\Sigma =\{\& \}\) podemos tomar \(TD_{3}\) igual a la maquina de la Figura 3.
Analogamente, para \(j\geq 1\), sea \(TI_{j}\) una maquina tal que

\(\displaystyle \begin{array}{ccc} \alpha B\sigma \gamma & \overset{\ast }{\vdash } & \alpha B\gamma \\ \uparrow \ & & \uparrow \\ q_{0}\ \ & & q_{f} \end{array} \)

cada vez que \(\alpha \in \Gamma ^{\ast }\), \(\sigma \in \Gamma \) y \(\gamma \) tiene exactamente \(j\) ocurrencias de \(B\). Es decir la maquina \(TI_{j}\) corre un espacio a la izquierda todo el bloque \(\gamma \) (por lo cual en el lugar de \(\sigma \) queda el primer simbolo de \(\gamma \)).
Teniendo las maquinas auxiliares antes definidas podemos combinarlas para obtener las maquinas simuladoras de instrucciones. Por ejemplo \(M_{i}^{a}\) puede ser la maquina descripta en la Figura 4. En la Figura 2 tenemos una posible forma de dise\~{n}ar la maquina \(IF_{i}^{a}\). En la Figura 7 tenemos una posible forma de dise\~{n}ar la maquina \(M_{i\leftarrow j}^{\ast }\) para el caso \(\Sigma =\{a,b\}\) y \(i< j\).

Supongamos ahora que \(\mathcal{P}=I_{1}...I_{n}\). Para cada \(i=1,...,n\), definiremos una maquina \(M_{i}\) que simulara la instruccion \(I_{i}\). Luego uniremos adecuadamente estas maquinas para formar la maquina que simulara a \( \mathcal{P}\)

- Si \(Bas(I_{i})=\mathrm{N}\bar{j}\leftarrow \mathrm{N}\bar{j}+1\) tomaremos \(M_{i}=M_{j}^{+}\)
- Si \(Bas(I_{i})=\mathrm{N}\bar{j}\leftarrow \mathrm{N}\bar{j}\dot{-} 1 \) tomaremos \(M_{i}=M_{j}^{\dot{-}}\)
- Si \(Bas(I_{i})=\mathrm{N}\bar{j}\leftarrow 0\) tomaremos \( M_{i}=M_{j\leftarrow 0}\).
- Si \(Bas(I_{i})=\mathrm{N}\bar{j}\leftarrow \mathrm{N}\bar{m}\) tomaremos \(M_{i}=M_{j\leftarrow m}^{\#}\).
- Si \(Bas(I_{i})=\mathrm{IF}\;\mathrm{N}\bar{j}\not=0\) \(\mathrm{GOTO} \;\mathrm{L}\bar{m}\) tomaremos \(M_{i}=IF_{j}.\)
- Si \(Bas(I_{i})=\mathrm{P}\bar{j}\leftarrow \mathrm{P}\bar{j}.a\) tomaremos \(M_{i}=M_{j}^{a}\).
- Si \(Bas(I_{i})=\mathrm{P}\bar{j}\leftarrow \ ^{\curvearrowright } \mathrm{P}\bar{j}\) tomaremos \(M_{i}=M_{j}^{\curvearrowright }\).
- Si \(Bas(I_{i})=\mathrm{P}\bar{j}\leftarrow \varepsilon \) tomaremos \( M_{i}=M_{j\leftarrow \varepsilon }\).
- Si \(Bas(I_{i})=\mathrm{P}\bar{j}\leftarrow \mathrm{P}\bar{m}\) tomaremos \(M_{i}=M_{j\leftarrow m}^{\ast }\).
- Si \(Bas(I_{i})=\mathrm{IF}\;\mathrm{P}\bar{j}\;\mathrm{BEGINS}\;a\; \mathrm{GOTO}\;\mathrm{L}\bar{m}\) tomaremos \(M_{i}=IF_{j}^{a}\).
- Si \(Bas(I_{i})=\mathrm{SKIP}\) tomaremos \(M_{i}=M_{\mathrm{SKIP}}\).
- Si \(Bas(I_{i})=\mathrm{GOTO}\;\mathrm{L}\bar{m}\) tomaremos \( M_{i}=M_{\mathrm{GOTO}}\).
Ya que la maquina \(M_{i}\) puede tener uno o dos estados finales, la representaremos como se muestra en la Figura 5, entendiendo que en el caso en que \(M_{i}\) tiene un solo estado final, este esta representado por el circulo de abajo a la izquierda y en el caso en que \(M_{i}\) tiene dos estados finales, el estado final representado con lineas punteadas corresponde al estado \(q_{si}\) y el otro al estado \(q_{no}\).

Para armar la maquina que simulara a \(\mathcal{P}\) hacemos lo siguiente. Primero unimos las maquinas \(M_{1},...,M_{n}\) como lo muestra la Figura 6. Luego para cada \(i\) tal que \(Bas(I_{i})\) es de la forma \(\alpha \mathrm{GOTO} \;\mathrm{L}\bar{m}\), ligamos con una flecha de la forma

\(\displaystyle \underrightarrow{\;\;\;\;\;\;B,B,K\;\;\;\;\;\;} \)

el estado final \(q_{si}\) de la \(M_{i}\) con el estado inicial de la \(M_{h}\), donde \(h\) es tal que \(I_{h}\) es la primer instruccion que tiene label \( \mathrm{L}\bar{m}\).
Es intuitivamente claro que la maquina asi obtenida cumple con lo requerido aunque una Proof formal de esto puede resultar extremadamente tediosa. \(\Box\)

\textbf{\underline{Theorem 85:}} Si \(f:D_{f}\subseteq \omega ^{n}\times \Sigma ^{\ast m}\rightarrow O\) es \( \Sigma \)-recursiva, entonces \(f\) es \(\Sigma \)-Turing computable.

\textbf{\underline{Proof:}} Supongamos \(O=\Sigma ^{\ast }.\) Ya que \(f\) es \(\Sigma \)-computable, existe \( \mathcal{P}\in \mathrm{Pro}^{\Sigma }\) el cual computa \(f\). Note que podemos suponer que \(\mathcal{P}\) tiene la propiedad de que cuando \(\mathcal{P}\) termina, en el estado alcansado las variables numericas tienen todas el valor \(0\) y las alfabeticas distintas de \(\mathrm{P}1\) todas el valor \( \varepsilon \). Sea \(M\) la maquina de Turing con unit dada por el lema anterior, donde elejimos el numero \(k\) con la propiedad adicional de ser mayor que \(n\) y \(m\). Sea \(M_{1}\) una maquina tal que para cada \((\vec{x}, \vec{\alpha})\in \omega ^{n}\times \Sigma ^{\ast m}\),

\(\displaystyle \left\lfloor q_{0}B\shortmid ^{x_{1}}B...B\shortmid ^{x_{n}}B\alpha _{1}B...B\alpha _{n}B\right\rfloor \overset{\ast }{\vdash }\left\lfloor qB\shortmid ^{x_{1}}B...B\shortmid ^{x_{n}}B^{k-n}B\alpha _{1}B...B\alpha _{m}B\right\rfloor \)

donde \(q_{0}\) es el estado inicial de \(M_{1}\) y \(q\) es un estado tal que \( \delta (q,\sigma )=\varnothing \), para cada \(\sigma .\) Sea \(M_{2}\) una maquina tal que para cada \(\alpha \in \Sigma ^{\ast }\),
\(\displaystyle \left\lfloor q_{0}B^{k+1}\alpha \right\rfloor \overset{\ast }{\vdash } \left\lfloor qB\alpha \right\rfloor \)

donde \(q_{0}\) es el estado inicial de \(M_{2}\) y \(q\) es un estado tal que \( \delta (q,\sigma )=\varnothing \), para cada \(\sigma \). Note que la concatenacion de \(M_{1}\), \(M\) y \(M_{2}\) (en ese orden) produce una maquina de Turing la cual computa \(f\). \(\Box\)



\textbf{\underline{Theorem 86:}} Si \(L\subseteq \Sigma ^{\ast }\) es \(\Sigma \)-r.e., entonces \(L=L(M)=H(M)\) para alguna maquina de Turing deterministica \(M.\)


\textbf{\underline{Proof:}} Por el Teorema 71 hay una funcion \(f:L\rightarrow \omega \), la cual es \(\Sigma \)-recursiva\(.\) Sea \(\mathcal{P}\) un programa el cual compute a \(f\). Sea \(M\) la maquina de Turing deterministica dada en el lema anterior. Sea \(M_{1}\) una maquina de Turing deterministica tal que para todo \(\alpha \in \Sigma ^{\ast }\),

\(\displaystyle \left\lfloor q_{0}B\alpha \right\rfloor \overset{\ast }{\vdash }\left\lfloor qB^{k+1}\alpha \right\rfloor \)

donde \(q_{0}\) es el estado inicial de \(M\) y \(q\) es un estado tal que \(\delta (q,\sigma )=\varnothing \), para cada \(\sigma \). Note que la concatenacion de \( M_{1}\) con \(M\) (en ese orden) produce una maquina de Turing deterministica \( M_{2}\) tal que \(H(M_{2})=L(M_{2})=L\). \(\Box\)


\begin{thebibliography}{X}
\bibitem{Baz} \textsc{Diego Vaggione},
<<Apunte de Clase, 2017>>,
\textit{FaMAF, UNC}.
\bibitem{Baz} \textsc{Agustín Curto},
<<Carpeta de Clase, 2017>>,
\textit{FaMAF, UNC}.
\end{thebibliography}

\vspace{\fill}
\begin{center}
Por favor, mejorá este documento en github
\includegraphics[width=1cm]{graphics/github.png} \\
https://github.com/acurto714/resumenLengForm
\end{center}
\end{document}
